\chapter{Multivariable calculus}
\label{vol02:ch11}

\epigraph{The calculus of one variable handles a road; the calculus
of many handles the country the road runs through.}{A working
paraphrase of the engineer's working register, in the vocabulary of
this volume.}

\chapmeta{Half-life: foundational. Archetypes: optimisation
(multivariable, building on Volume II Chapter 5's one-variable
case), scaling (high-dimensional volume concentration as the
canonical case). Exercise target: 35. Project track: analysis only.
Hazard class: none. Reader path default: standard, with sections
explicitly tagged. Named cases: none.}

\noindent
A function of one variable describes a quantity that depends on a
single input. Most engineering quantities depend on several. The
temperature in a room depends on three spatial coordinates and on
time. The cost of a product depends on the cost of every input. The
loss of a fitted model depends on every parameter. Volume II
Chapter 5 introduced the derivative for one variable. Volume II
Chapter 8 introduced the gradient as a vector field. This chapter
joins the two: the gradient is the derivative when there are many
variables, the Hessian is the second derivative, the critical-point
condition becomes a vector equation, and the optimisation archetype
acquires its working multivariable form. The reader leaves able to
locate a critical point in two or more variables, classify it,
constrain it by equalities and inequalities through Lagrange
multipliers and the Karush-Kuhn-Tucker conditions, set up a multiple
integral in whatever coordinate system the geometry asks for, and
notice when high-dimensional intuition breaks.

\section{Functions of many variables}\pathtag{core}

A \engterm{function of many variables} is a rule that assigns to
each point of a region in $\R^{n}$ a single real number. We write
$f : \R^{n} \to \R$ and read it as $f$ takes $n$ inputs and returns
one output. The inputs are written as a vector $\vect{x} = (x_{1},
x_{2}, \ldots, x_{n})$, and the output is $f(\vect{x})$. The two
extremes the reader has already met: $n = 1$ in Volume II
Chapter 5, and $n = 3$ in the scalar fields of Volume II
Chapter 8.

\subsection{Examples the engineer keeps to hand}

The temperature inside a heated wall is a function of position
and time, $T(x, y, z, t) : \R^{4} \to \R$. The kinetic energy of a
rigid body depends on three velocity components and three angular-
velocity components, $K(\vect{v}, \boldsymbol{\omega}) : \R^{6} \to
\R$. The training loss of a neural network with $p$ parameters is a
function $L(\boldsymbol{\theta}) : \R^{p} \to \R$, where $p$ may be
in the millions. The volume of a closed cylinder is a function of
two design parameters, $V(r, h) = \pi r^{2} h$, $V : \R^{2} \to \R$.
The chapter develops the apparatus on small examples (mostly
$n = 2$ or $n = 3$) and the reader applies it at the scale the
problem demands.

\subsection{Domain, range, level sets}

The \engterm{domain} of $f$ is the set of $\vect{x}$ for which
$f(\vect{x})$ is defined; the \engterm{range} is the set of values
the function takes. The most useful geometric object attached to
$f : \R^{n} \to \R$ is the family of \engterm{level sets} (or
\engterm{contours} when $n = 2$): for each value $c$ in the range,
the level set is
\[
S_{c} = \set{\vect{x} \in \R^{n} : f(\vect{x}) = c}.
\]
For $n = 2$, the level sets of a smooth $f$ are typically curves;
the reader has seen them on a topographic map (constant elevation),
on a weather map (constant pressure), and on a phase diagram
(constant Gibbs free energy). For $n = 3$, level sets are typically
surfaces. The contours of $f(x, y) = x^{2} + y^{2}$ are circles of
radius $\sqrt{c}$; the level sets of $f(x, y, z) = x^{2} + y^{2} +
z^{2}$ are spheres.

\paragraph{Worked example: natural domains.} The domain is often
forced by the formula. For $f(x, y) = \ln(1 - x^{2} - y^{2})$, the
logarithm requires its argument positive, so the domain is the open
unit disk $x^{2} + y^{2} < 1$; the function diverges to $-\infty$ as
the boundary circle is approached, and the boundary itself is not in
the domain. For $g(x, y) = \sqrt{x - y}$, the square root requires
$x \geq y$, so the domain is the closed half-plane on and below the
line $y = x$. For $h(x, y) = 1/(x^{2} + y^{2})$, the domain is the
plane with the origin removed, where the function would divide by
zero. The reader states the domain before differentiating or
integrating; an integral set up over a region that strays outside
the natural domain is meaningless, and a great many "the answer is
infinite" surprises trace to a domain boundary the setup ignored.

The contour map is the engineer's primary visualisation in two
dimensions. A function increases fastest in the direction
perpendicular to its contours, decreases fastest in the opposite
direction, and is locally constant along the contours. Section 11.2
makes this precise as the gradient theorem.
\Cref{fig:vol02:ch11:contour-gradient} shows the contours of
$f(x, y) = x^{2} + 2 y^{2}$ together with the gradient at four
sample points, anticipating the geometric content of the next
section.

% Figure: contour plot of f(x,y) = x^2 + 2 y^2 with gradient arrows
% perpendicular to the contours, illustrating the steepest-ascent
% theorem. Used in section 11.2 prose.
\begin{figure}[htbp]
\centering
\begin{tikzpicture}[scale=1.0,
  contour/.style={engblue, thick},
  arrow/.style={->, engred, very thick},
  axis/.style={engblack, thick, ->},
  label/.style={font=\small},
]
% Axes
\draw[axis] (-3.2, 0) -- (3.4, 0) node[right, label] {$x$};
\draw[axis] (0, -2.4) -- (0, 2.6) node[above, label] {$y$};

% Concentric ellipses: f(x,y) = x^2 + 2 y^2 = c, so x^2/c + y^2/(c/2) = 1
% Semi-axes (sqrt c, sqrt(c/2)). Choose c = 0.5, 1, 2, 4.
\draw[contour] (0,0) ellipse ({sqrt(0.5)} and {sqrt(0.25)});
\draw[contour] (0,0) ellipse ({sqrt(1.0)} and {sqrt(0.5)});
\draw[contour] (0,0) ellipse ({sqrt(2.0)} and {sqrt(1.0)});
\draw[contour] (0,0) ellipse ({sqrt(4.0)} and {sqrt(2.0)});

% Contour labels
\node[label, engblue] at (2.15, 0.18) {$c=4$};
\node[label, engblue] at (1.50, 0.18) {$c=2$};
\node[label, engblue] at (1.05, 0.18) {$c=1$};

% Gradient arrows at a few sample points: grad f = (2x, 4y).
% Scale arrows for visual clarity (length = 0.15 * grad).
% Point A: (1, 0.5). grad = (2, 2). Length 0.25.
\draw[arrow] (1, 0.5) -- ++({2*0.25}, {2*0.25});
\filldraw[engred] (1, 0.5) circle (0.04);
% Point B: (-1, 0.7). grad = (-2, 2.8).
\draw[arrow] (-1, 0.7) -- ++({-2*0.25}, {2.8*0.25});
\filldraw[engred] (-1, 0.7) circle (0.04);
% Point C: (1.4, -0.6). grad = (2.8, -2.4).
\draw[arrow] (1.4, -0.6) -- ++({2.8*0.25}, {-2.4*0.25});
\filldraw[engred] (1.4, -0.6) circle (0.04);
% Point D: (0, 1.0). grad = (0, 4).
\draw[arrow] (0, 1.0) -- ++(0, {4*0.18});
\filldraw[engred] (0, 1.0) circle (0.04);

% Annotation: gradient is perpendicular to the contour at every point
\node[label, anchor=west] at (1.55, 1.1) {$\grad f \perp$ contour};

\end{tikzpicture}
\caption[Contours of $f(x,y) = x^{2} + 2y^{2}$ with gradient
arrows.]{Contours of $f(x, y) = x^{2} + 2y^{2}$ (blue ellipses) with
$\grad f = (2x, 4y)$ drawn at four sample points (red). At every
point, the gradient is perpendicular to the contour through that
point, and its magnitude reports the rate of steepest ascent. The
asymmetry of the ellipses reflects the asymmetric coefficient on
$y^{2}$: the function increases twice as fast per unit of $y$ as
per unit of $x$.}
\label{fig:vol02:ch11:contour-gradient}
\end{figure}


Reading a contour map is a skill the reader builds by repetition.
The spacing between adjacent contours, drawn at equal increments of
$c$, reports the magnitude of the slope: closely spaced contours
mean a steep function, widely spaced contours a flat one. A circular
nest of contours that shrink to a point marks a local maximum or
minimum; contours that cross in an X mark a saddle. The reader who
has walked a ridge line on a topographic map already knows the
geometry: the trail that stays on one contour is level, the trail
that cuts across contours climbs or descends, and the steepest path
crosses contours at right angles.

\paragraph{Worked example: reading level sets of three surfaces.}
Take three functions and identify their level sets without plotting.
For $f(x, y) = 3x - 2y$, each level set $3x - 2y = c$ is a straight
line of slope $3/2$; the level sets are parallel lines, equally
spaced for equal increments of $c$, because the function is linear
and its slope is constant. For $f(x, y) = x^{2} - y^{2}$, the level
set $x^{2} - y^{2} = c$ is a hyperbola opening left-right when
$c > 0$, opening up-down when $c < 0$, and degenerating to the pair
of lines $y = \pm x$ when $c = 0$. The two asymptotic lines are the
signature of a saddle, which section 11.3 confirms analytically. For
$f(x, y) = e^{-(x^{2} + y^{2})}$, the level set $e^{-(x^{2} + y^{2})}
= c$ requires $x^{2} + y^{2} = -\ln c$, a circle of radius
$\sqrt{-\ln c}$ that exists only for $0 < c \leq 1$; the contours are
concentric circles crowding together near the central peak at the
origin and spreading out in the tails. The Gaussian bump is the
canonical shape of a two-dimensional probability density (Volume II
Chapter 13) and of a diffraction-limited light spot (Volume V).

\paragraph{Worked example: the graph as a surface.} The
\engterm{graph} of $f : \R^{2} \to \R$ is the surface $z = f(x, y)$
in $\R^{3}$. The graph and the contour map carry the same
information in two presentations. For $f(x, y) = x^{2} + y^{2}$, the
graph is a paraboloid of revolution opening upward, and the contours
are the circles obtained by slicing the paraboloid with horizontal
planes $z = c$. For $f(x, y) = \sqrt{x^{2} + y^{2}}$, the graph is a
cone with vertex at the origin, and the contours are again circles,
but now equally spaced in radius because the cone climbs linearly.
The reader distinguishes the two surfaces by their contour spacing:
the paraboloid's circles crowd together as $c$ grows (radius
$\sqrt{c}$, growing sublinearly), while the cone's circles are evenly
spaced (radius $c$). The contour spacing is the visual reading of
the gradient magnitude, formalised in the next section.

\subsection{Continuity, the multivariable form}

A function $f : \R^{n} \to \R$ is \engterm{continuous} at a point
$\vect{a}$ if $f(\vect{x}) \to f(\vect{a})$ as $\vect{x} \to
\vect{a}$ along every path. The "every path" clause distinguishes
the multivariable case from the one-variable case. In one variable,
$x \to a$ has only two approach directions; in two or more, the
approach can be along a line, along a curve, or along an oscillating
spiral. A function may have a limit along every straight line and
fail to be continuous because some non-straight path gives a
different limit. The standard cautionary example is
\[
f(x, y) =
\begin{cases}
\dfrac{xy}{x^{2} + y^{2}} & (x, y) \neq (0, 0), \\
0 & (x, y) = (0, 0).
\end{cases}
\]
Along any line $y = mx$, the function takes the constant value
$m / (1 + m^{2})$; along the line $x = 0$, it is identically zero.
Different lines give different limits at the origin, so $f$ is not
continuous there. The engineer keeps this in mind when fitting a
surface to scattered data: a fit that looks reasonable along every
sampled direction can still be discontinuous between directions.

\paragraph{Worked example: a limit that exists.} Continuity can also
be confirmed, not only refuted, by the path test followed by an
estimate. Consider
\[
g(x, y) = \frac{x^{2} y}{x^{2} + y^{2}}, \qquad g(0, 0) = 0.
\]
Along $y = mx$, $g = m x^{3} / (x^{2}(1 + m^{2})) = m x / (1 + m^{2})
\to 0$ as $x \to 0$, so every straight line gives the limit $0$. The
straight-line test is suggestive but not conclusive; the cautionary
$xy/(x^{2} + y^{2})$ above passed every straight-line test and still
failed. The conclusive bound uses polar coordinates: with $x =
r\cos\theta$, $y = r\sin\theta$,
\[
\abs{g(x, y)} = \abs{\frac{r^{3}\cos^{2}\theta\sin\theta}{r^{2}}}
= r\,\abs{\cos^{2}\theta\sin\theta} \leq r \to 0
\]
as $r \to 0$, uniformly in $\theta$ because the trigonometric factor
is bounded by $1$. The bound holds for every approach path, so
$g \to 0$ and $g$ is continuous at the origin. The working rule: a
straight-line limit that disagrees between lines proves
discontinuity, but only a path-independent bound (typically through
the polar radius $r$) proves continuity.

\paragraph{Worked example: the parabola that betrays the lines.} The
straight-line test is not merely inconclusive, it can be actively
misleading; a function can pass every straight-line test and still fail
along a curve. Consider
\[
h(x, y) = \frac{x^{2} y}{x^{4} + y^{2}}, \qquad h(0, 0) = 0.
\]
Along any line $y = mx$, $h = m x^{3}/(x^{4} + m^{2} x^{2}) =
m x/(x^{2} + m^{2}) \to 0$ as $x \to 0$, so every straight line gives
the limit $0$, and along the axes $h \equiv 0$ as well. Yet along the
parabola $y = x^{2}$,
\[
h(x, x^{2}) = \frac{x^{2}\cdot x^{2}}{x^{4} + x^{4}}
= \frac{x^{4}}{2 x^{4}} = \frac{1}{2}
\]
for every $x \neq 0$, so the limit along that parabola is $\tfrac{1}{2}
\neq 0$. The function is discontinuous at the origin, and no collection
of straight-line tests, however large, could have caught it. The
parabola matched the denominator's two powers of $x$ against the
numerator's, which no line could do. The engineering warning sharpens the
earlier one: a response surface swept along straight rays through an
operating point can look perfectly smooth and still hide a ridge along a
curved path that a coupled two-factor move would find. The polar-radius
bound that proved continuity for $g$ above fails here precisely because
no bound of the form $\abs{h} \leq C r^{p}$ holds: the $x^{4}$ in the
denominator makes $h$ scale differently along different-curvature paths,
which is the analytic signature of curve-dependent discontinuity.

\begin{mastery}
The distinction between joint continuity (continuity as a function
of $\vect{x}$) and separate continuity (continuity in each variable
with the others fixed) is the source of the path-dependence above.
The function $xy/(x^{2} + y^{2})$ is separately continuous at the
origin (each one-variable slice through the origin is continuous)
yet not jointly continuous. Separate continuity is the property a
naive coordinate-by-coordinate sampler tests; joint continuity is
the property the engineer actually needs. The two coincide for
functions with bounded partial derivatives in a neighbourhood (a
Lipschitz condition), which is why every smooth model in this volume
escapes the pathology. The reader who fits a response surface by
sweeping one factor at a time is testing separate continuity and
silently assuming joint continuity; an interaction term that the
one-factor-at-a-time sweep never excites is the engineering version
of a non-straight approach path.
\end{mastery}

\section{Partial derivatives and gradients}\pathtag{core}

Volume II Chapter 5 defined the derivative of $f : \R \to \R$ as
the limit of the difference quotient. The same construction applied
one variable at a time gives the partial derivative.

\subsection{Partial derivatives}

\begin{definition}[Partial derivative]
Let $f : \R^{n} \to \R$ be defined on a region containing
$\vect{x}$. The \engterm{partial derivative} of $f$ with respect to
$x_{i}$ at $\vect{x}$ is
\[
\pdv{f}{x_{i}}(\vect{x}) = \lim_{h \to 0}
\frac{f(\vect{x} + h \unitvec{e}_{i}) - f(\vect{x})}{h},
\]
where $\unitvec{e}_{i}$ is the unit vector in the $i$th coordinate
direction, provided the limit exists.
\end{definition}

The partial derivative measures the rate of change of $f$ when only
the $i$th input varies and all others are held fixed. The notation
$\pdv{f}{x_{i}}$ is standard; the alternative $f_{x_{i}}$ is shorter
and used when the subscripts are unambiguous. The reader computes a
partial derivative by treating every other variable as a constant
and applying the one-variable rules of Volume II Chapter 5.

The geometric reading of $\partial f / \partial x$ is the slope of
the curve cut from the graph $z = f(x, y)$ by the vertical plane
$y = y_{0}$. \Cref{fig:vol02:ch11:partial-slice} shows the
construction: holding all coordinates but the $i$th fixed reduces
the multivariable problem to a one-variable problem on the slice,
and the partial derivative inherits the one-variable derivative of
that slice.

% Figure: a partial derivative as the slope of the slice obtained by
% holding all but one variable fixed. Surface z = f(x, y), the slice
% y = y0, and the tangent line whose slope is f_x(x0, y0).
\begin{figure}[htbp]
\centering
\begin{tikzpicture}[scale=1.0, line cap=round,
  axis/.style={engblack, thick, ->},
  slice/.style={engred, thick},
  tangent/.style={engblue, very thick},
  surf/.style={engblack!50, thin},
  label/.style={font=\small},
]
% 3D-look axes (oblique projection): x right, y back, z up.
\draw[axis] (0,0) -- (4.5, 0) node[right, label] {$x$};
\draw[axis] (0,0) -- (-1.8, -1.3) node[below, label] {$y$};
\draw[axis] (0,0) -- (0, 3.5) node[above, label] {$z = f(x, y)$};

% Sketch a surface as a grid of curves (saddle-ish: f = x^2/2 - y^2/4 + 1).
% Background curves: y = const lines
\foreach \yv in {-1.5, -0.5, 0.5} {
  \draw[surf] plot[domain=0.1:3.5, samples=30, smooth]
    ({\x + (-1.8/4.5)*\yv}, {-1.3/4.5*\yv + (\x*\x)/2*0.18 - (\yv*\yv)/4*0.4 + 0.3});
}
% x = const lines
\foreach \xv in {1.0, 2.0, 3.0} {
  \draw[surf] plot[domain=-1.5:0.5, samples=30, smooth]
    ({\xv + (-1.8/4.5)*\x}, {-1.3/4.5*\x + (\xv*\xv)/2*0.18 - (\x*\x)/4*0.4 + 0.3});
}

% Highlight slice y = y0 (here y0 = -0.5): traces out a curve in the (x, z) plane
\draw[slice] plot[domain=0.1:3.8, samples=40, smooth]
  ({\x + (-1.8/4.5)*(-0.5)}, {-1.3/4.5*(-0.5) + (\x*\x)/2*0.18 - 0.1 + 0.3});

% Mark the point (x0, y0) at x0 = 2: z = 4/2*0.18 - 0.1 + 0.3 = 0.56
\coordinate (P) at ({2 + (-1.8/4.5)*(-0.5)}, {-1.3/4.5*(-0.5) + (2*2)/2*0.18 - 0.1 + 0.3});
\filldraw[engblue] (P) circle (0.06);

% Tangent to the slice at P: slope = d/dx (x^2/2)*0.18 = 0.18 x = 0.36 at x=2
% In our display coords, the tangent line has slope (dz/dx)/(1) (since dy = 0 along the slice)
\draw[tangent] ($(P) + (-1, {-1*0.36})$) -- ($(P) + (1.2, {1.2*0.36})$);

% Labels
\node[label, engred, anchor=west] at (3.7, 0.95) {slice $y = y_{0}$};
\node[label, engblue, anchor=south west] at (P) {$\partial f / \partial x \big|_{(x_{0}, y_{0})}$};

% Mark x0 on the x-axis
\draw[densely dotted, engblack] (2, 0) -- (P);
\node[label, anchor=north] at (2, 0) {$x_{0}$};

\end{tikzpicture}
\caption[Partial derivative as the slope of a one-variable
slice.]{The partial derivative $\partial f / \partial x$ at the
point $(x_{0}, y_{0})$ is the slope of the tangent to the curve cut
from the surface $z = f(x, y)$ by the vertical plane $y = y_{0}$.
The slice (red) reduces the bivariate problem to a one-variable
problem in $x$; the tangent line (blue) inherits its slope from the
one-variable derivative of that slice at $x_{0}$.}
\label{fig:vol02:ch11:partial-slice}
\end{figure}


\paragraph{A worked example.} For $f(x, y) = x^{2} y + \sin(xy)$,
\begin{align*}
\pdv{f}{x} &= 2 x y + y \cos(xy), \\
\pdv{f}{y} &= x^{2} + x \cos(xy).
\end{align*}
The product rule and the chain rule apply variable by variable.

\paragraph{Worked example: an ideal-gas state function.} The ideal-
gas law $PV = nRT$ gives pressure as a function of volume and
temperature, $P(V, T) = nRT/V$. The two partials report distinct
physical sensitivities,
\[
\pdv{P}{V} = -\frac{nRT}{V^{2}}, \qquad
\pdv{P}{T} = \frac{nR}{V}.
\]
The first is negative: at fixed temperature, compressing the gas
(decreasing $V$) raises the pressure, with sensitivity scaling as
$1/V^{2}$. The second is positive: at fixed volume, heating the gas
raises the pressure linearly in temperature. The reader carries the
sign and the scaling, not only the formula. A thermodynamics chapter
(Volume IV) builds the apparatus of partial derivatives of state
functions into the machinery of Maxwell relations, where the
equality of mixed partials below becomes a physical law rather than
a calculus identity.

\paragraph{Worked example: a Cobb-Douglas cost surface.} A
production model with two inputs, labour $L$ and capital $K$, often
takes the Cobb-Douglas form $Q(L, K) = A L^{\alpha} K^{\beta}$ with
$A, \alpha, \beta > 0$. The marginal products are the two partials,
\[
\pdv{Q}{L} = \alpha A L^{\alpha - 1} K^{\beta}
= \alpha\,\frac{Q}{L}, \qquad
\pdv{Q}{K} = \beta A L^{\alpha} K^{\beta - 1}
= \beta\,\frac{Q}{K}.
\]
The compact forms $\alpha Q/L$ and $\beta Q/K$ are the marginal
products expressed as the output elasticities $\alpha$ and $\beta$
times the average product. The ratio of the two marginal products,
$(\alpha K)/(\beta L)$, is the marginal rate of technical
substitution, the slope of the level set (the isoquant $Q = c$) that
the Lagrange analysis of section 11.4 reads as the optimal input
mix. The partial derivative is the bridge from a production formula
to an optimisation decision.

\subsection{Higher-order partials and Clairaut's theorem}

A function of several variables has several second derivatives. For
$f(x, y)$, the four second partials are $f_{xx}, f_{xy}, f_{yx},
f_{yy}$. The mixed partials $f_{xy}$ and $f_{yx}$ measure the rate
at which $f_{x}$ changes with $y$ and the rate at which $f_{y}$
changes with $x$, and they are usually equal.

\begin{theorem}[Clairaut, equality of mixed partials]
If $f$ has continuous second partial derivatives on an open set,
then on that set
\[
\pdv{}{y}\!\left(\pdv{f}{x}\right) =
\pdv{}{x}\!\left(\pdv{f}{y}\right).
\]
\end{theorem}

The hypothesis of continuity is non-trivial: there are pathological
functions for which $f_{xy} \neq f_{yx}$ at the origin, but every
function the engineer meets in practice satisfies the hypothesis.
The working consequence is that the order of partial differentiation
does not matter, and the second-derivative information of $f$ at a
point lives in a symmetric matrix.

\paragraph{Worked example: verifying Clairaut.} For $f(x, y) =
x^{3} y^{2} + e^{x} \sin y$, the first partials are $f_{x} = 3 x^{2}
y^{2} + e^{x}\sin y$ and $f_{y} = 2 x^{3} y + e^{x}\cos y$. The two
mixed partials,
\[
f_{xy} = \pdv{}{y}(3 x^{2} y^{2} + e^{x}\sin y) = 6 x^{2} y +
e^{x}\cos y, \qquad
f_{yx} = \pdv{}{x}(2 x^{3} y + e^{x}\cos y) = 6 x^{2} y +
e^{x}\cos y,
\]
agree, as Clairaut promises for a function with continuous second
partials everywhere. The reader uses the equality as a cheap audit:
a hand computation in which $f_{xy} \neq f_{yx}$ contains an
arithmetic error, because the functions of engineering practice
never violate the hypothesis. The equality of mixed partials is the
single most useful check on a multivariable differentiation.

\subsection{The gradient}

\begin{definition}[Gradient]
The \engterm{gradient} of $f : \R^{n} \to \R$ at $\vect{x}$ is the
vector of partial derivatives,
\[
\grad f(\vect{x}) = \left(\pdv{f}{x_{1}}, \pdv{f}{x_{2}}, \ldots,
\pdv{f}{x_{n}}\right)^{\mathsf{T}}.
\]
\end{definition}

Volume II Chapter 8 introduced $\grad\Phi$ as a vector field on
$\R^{3}$. The definition above generalises to $\R^{n}$: the gradient
is the column vector of all $n$ partial derivatives. For
$f(x, y, z) = x^{2} + 2 y^{2} + 3 z^{2}$, the gradient is $\grad f =
(2x, 4y, 6z)^{\mathsf{T}}$.

\subsection{The directional derivative and the steepest-ascent
theorem}

The partial derivatives report the rate of change along the
coordinate axes. The reader more often wants the rate along an
arbitrary direction. Let $\unitvec{u}$ be a unit vector in $\R^{n}$.
The \engterm{directional derivative} of $f$ at $\vect{x}$ in the
direction $\unitvec{u}$ is
\[
D_{\unitvec{u}} f(\vect{x}) = \lim_{h \to 0}
\frac{f(\vect{x} + h\unitvec{u}) - f(\vect{x})}{h}.
\]

\begin{theorem}[Gradient and directional derivative]
If $f$ is differentiable at $\vect{x}$, then for every unit vector
$\unitvec{u}$,
\[
D_{\unitvec{u}} f(\vect{x}) = \grad f(\vect{x}) \cdot \unitvec{u}.
\]
\end{theorem}

\begin{proof}[Sketch]
Write $g(h) = f(\vect{x} + h \unitvec{u})$, a function of one
variable, and apply the multivariable chain rule of section 11.2.5
to $g'(0)$.
\end{proof}

The geometric consequence is the steepest-ascent theorem.

\begin{principle}[Gradient is direction of steepest ascent]
At a point where $\grad f(\vect{x}) \neq \vect{0}$, the directional
derivative $D_{\unitvec{u}} f(\vect{x})$ is maximised by the choice
$\unitvec{u} = \grad f / \norm{\grad f}$, with maximum value
$\norm{\grad f(\vect{x})}$. In words: the gradient points in the
direction of fastest increase, with magnitude equal to that fastest
rate; the negative gradient points in the direction of fastest
decrease.
\end{principle}

\begin{proof}
$D_{\unitvec{u}} f = \grad f \cdot \unitvec{u} = \norm{\grad f}
\cos\theta$ where $\theta$ is the angle between $\grad f$ and
$\unitvec{u}$. The right-hand side is maximised at $\cos\theta = 1$,
which means $\unitvec{u}$ aligned with $\grad f$. The maximum value
is $\norm{\grad f}$.
\end{proof}

The steepest-ascent theorem is the foundation of every iterative
optimisation method we meet in this volume and the next. Gradient
descent, introduced briefly in section 11.5 and developed in
Volume II Chapter 15, takes a step in the direction $-\grad f$ at
each iterate. The same theorem motivates the gradient-flow
trajectory of a marble released on a hilly surface and the heat-
diffusion direction in a thermal gradient (Volume II Chapter 12).

\paragraph{Worked example: a directional derivative on a hill.} A
terrain model gives elevation $f(x, y) = 100 - 0.01 x^{2} - 0.02
y^{2}$ in metres, with $x, y$ in metres. A hiker stands at $(30,
20)$. The gradient is $\grad f = (-0.02 x, -0.04 y) = (-0.6, -0.8)$,
of magnitude $\norm{\grad f} = \sqrt{0.6^{2} + 0.8^{2}} = 1.0$. The
hiker walks due north-east, $\unitvec{u} = (1, 1)/\sqrt{2}$. The
rate of elevation change in that direction is
\[
D_{\unitvec{u}} f = \grad f \cdot \unitvec{u} =
\frac{-0.6 - 0.8}{\sqrt{2}} = \frac{-1.4}{\sqrt{2}}
\approx -0.99\,\frac{\text{m elevation}}{\text{m travelled}},
\]
a descent of almost one metre of elevation per metre walked. The
steepest descent is along $-\grad f / \norm{\grad f} = (0.6, 0.8)$,
giving the maximal rate $-\norm{\grad f} = -1.0$. The chosen
north-east heading happens to be close to the fall line, which is
why its rate ($-0.99$) is nearly the steepest possible ($-1.0$); the
two would coincide exactly only if the hiker walked along $(0.6,
0.8)$. A contour line through the hiker's position runs perpendicular
to $\grad f$, along the direction $(0.8, -0.6)$; walking that
heading keeps the elevation constant, $D_{\unitvec{u}} f = 0$, the
level traverse a tired hiker prefers.

\paragraph{Worked example: gradient and the chemical-concentration
field.} A pollutant concentration in a still pond is $C(x, y) =
C_{0}\,e^{-(x^{2} + y^{2})/L^{2}}$, with $L$ a length scale. The
gradient is
\[
\grad C = C_{0}\,e^{-(x^{2} + y^{2})/L^{2}}\cdot
\left(-\frac{2x}{L^{2}}, -\frac{2y}{L^{2}}\right)
= -\frac{2}{L^{2}}\,C\,(x, y).
\]
It points radially inward, toward the source at the origin, with
magnitude $\norm{\grad C} = (2/L^{2})\,C\,\sqrt{x^{2} + y^{2})}$ that
vanishes at the source (a maximum) and in the far field (where $C$
itself vanishes), peaking at intermediate radius. Fick's law of
diffusion sends the pollutant flux down the gradient, $\vect{J} =
-D\,\grad C$, that is, radially outward from the source; the
diffusive spreading of a contaminant is the negative-gradient flow
the steepest-descent theorem describes. Volume IV develops the
diffusion equation that governs how $C$ evolves; this chapter
supplies the gradient that drives it.

\paragraph{Worked example: one step of gradient descent.} Minimise
$f(x, y) = x^{2} + 4 y^{2}$ by gradient descent from the start
$(x_{0}, y_{0}) = (4, 2)$ with step size $\eta = 0.1$. The gradient
is $\grad f = (2x, 8y)$, so at the start $\grad f(4, 2) = (8, 16)$.
The update $\vect{x}_{1} = \vect{x}_{0} - \eta\,\grad f$ gives
\[
\vect{x}_{1} = (4, 2) - 0.1\,(8, 16) = (4 - 0.8, 2 - 1.6)
= (3.2, 0.4).
\]
The objective drops from $f(4, 2) = 16 + 16 = 32$ to $f(3.2, 0.4) =
10.24 + 0.64 = 10.88$, a large first step toward the minimum at the
origin. The asymmetry of the descent is visible: the $y$-coordinate
moved much more than the $x$-coordinate (the step in $y$ was $1.6$
versus $0.8$ in $x$) because the $4y^{2}$ term makes the function
four times as steep in $y$. A single step size $\eta$ applied to both
coordinates over-corrects the steep direction and under-corrects the
gentle one, the condition-number pathology in miniature. The reader
who tracks a few more steps sees the iterate zig-zag toward the
origin, the behaviour the Rosenbrock simulation exercise studies at
its worst.

% Figure: gradient descent on an anisotropic quadratic bowl, showing the
% zig-zag path that the condition-number mismatch produces. Contours are the
% ellipses of f = x^2 + 4 y^2; the iterate crosses the steep valley at every
% step instead of walking down it.
\begin{figure}[htbp]
\centering
\begin{tikzpicture}[scale=1.0,
  contour/.style={engblue, thick},
  gdstep/.style={engred, very thick, -{Latex[length=2.2mm]}},
  dot/.style={engred, fill=engred},
]
  % Elliptical contours of f = x^2 + 4 y^2 = c, so x = sqrt(c) cos t,
  % y = sqrt(c)/2 sin t. Drawn for c = 1, 4, 9, 16.
  \foreach \c in {1, 4, 9, 16} {
    \draw[contour, domain=0:360, samples=80, smooth]
      plot ({sqrt(\c)*cos(\x)}, {0.5*sqrt(\c)*sin(\x)});
  }
  % Axes
  \draw[engblack!45, -{Latex[length=1.6mm]}] (-4.6, 0) -- (4.6, 0) node[right]{$x$};
  \draw[engblack!45, -{Latex[length=1.6mm]}] (0, -2.4) -- (0, 2.4) node[above]{$y$};
  % Iterates of gradient descent from (4,2) with eta = 0.1 on grad = (2x, 8y).
  % x_{k+1} = 0.8 x_k, y_{k+1} = 0.2 y_k. Computed offline.
  \coordinate (p0) at (4.000, 2.000);
  \coordinate (p1) at (3.200, 0.400);
  \coordinate (p2) at (2.560, 0.080);
  \coordinate (p3) at (2.048, 0.016);
  \coordinate (p4) at (1.638, 0.003);
  \draw[gdstep] (p0) -- (p1);
  \draw[gdstep] (p1) -- (p2);
  \draw[gdstep] (p2) -- (p3);
  \draw[gdstep] (p3) -- (p4);
  \foreach \p in {p0, p1, p2, p3, p4} {
    \filldraw[dot] (\p) circle (0.055);
  }
  \filldraw[englime] (0,0) circle (0.07);
  \node[engred, anchor=south west] at (p0) {$\vect{x}_{0}$};
  \node[englime, anchor=north east] at (-0.1, -0.1) {min};
\end{tikzpicture}
\caption[Gradient descent zig-zag on an anisotropic bowl.]{Gradient
descent on $f(x, y) = x^{2} + 4 y^{2}$ from the start $\vect{x}_{0} =
(4, 2)$ with step size $\eta = 0.1$. The level sets are ellipses four
times as steep in $y$ as in $x$, so a single step size overshoots the
steep $y$-direction and undershoots the gentle $x$-direction. The
iterate collapses the $y$-coordinate almost at once and then crawls
along the $x$-axis toward the minimum (green dot at the origin), the
condition-number pathology that Newton's method and adaptive methods
correct. The step vectors are drawn to scale; each halves-and-more the
distance to the axis in $y$ while shrinking it only to four-fifths in
$x$.}
\label{fig:vol02:ch11:descent-zigzag}
\end{figure}


\Cref{fig:vol02:ch11:descent-zigzag} tracks the full run of the same
descent for several steps. The iterate flattens the $y$-coordinate in a
single step and then walks the $x$-axis toward the origin in ever
smaller steps, tracing the zig-zag characteristic of a mismatched step
size on an anisotropic bowl.

\paragraph{Worked example: the directional-derivative rose.} Fix a point
where $\grad f = (3, 0)^{\mathsf{T}}$ and read the rate of change in
every heading at once. By the gradient-and-directional-derivative
theorem, $D_{\unitvec{u}} f = \grad f \cdot \unitvec{u} = 3\cos\theta$,
where $\theta$ is the angle from the positive $x$-axis. The rate is
$+3$ due east (the fall line of ascent), $0$ due north and due south
(along the contour), and $-3$ due west (steepest descent). At the
compass headings $\theta = 60^{\circ}$ and $\theta = 120^{\circ}$ the
rate is $3\cos 60^{\circ} = 1.5$ and $3\cos 120^{\circ} = -1.5$, half
the peak in each sense. \Cref{fig:vol02:ch11:directional-rose} draws the
signed rate as an arrow in each heading; the outline of the arrow tips
traces the cosine lobe that the steepest-ascent theorem predicts. The
single number $\norm{\grad f} = 3$ and the single direction of
$\grad f$ fix the rate in every other direction, the reason the gradient,
not the twelve compass rates, is the object the engineer records.

% Figure: the directional derivative D_u f = |grad f| cos theta drawn as a
% rose of arrows around a point. The arrow length is proportional to the
% signed rate of change in each heading; it is largest along grad f, zero
% along the two contour-tangent headings, most negative opposite grad f.
\begin{figure}[htbp]
\centering
\begin{tikzpicture}[scale=1.0,
  grad/.style={engred, very thick, -{Latex[length=2.6mm]}},
  ray/.style={engblue, thick, -{Latex[length=1.8mm]}},
  zero/.style={enggray, thick, dashed},
]
  % grad f points along +x here (heading 0), magnitude 1 (drawn at 2.6 units).
  \def\G{2.6}
  % Rose: for heading angle a (degrees), D_u = |grad f| cos a. Draw an arrow
  % of length \G*cos(a) when positive (blue outward) and negative headings
  % as inward blue arrows of length \G*|cos(a)|.
  \foreach \a in {0,30,60,90,120,150,180,210,240,270,300,330} {
    \pgfmathsetmacro{\len}{\G*cos(\a)}
    \pgfmathsetmacro{\alen}{abs(\len)}
    \ifdim\alen pt>0.05pt
      \draw[ray] (0,0) -- ({\alen*cos(\a)}, {\alen*sin(\a)});
    \fi
  }
  % The gradient direction highlighted.
  \draw[grad] (0,0) -- ({\G*cos(0)}, {\G*sin(0)})
    node[anchor=west]{$\grad f$};
  % Contour tangent: headings 90 and 270 give D_u = 0 (dashed line).
  \draw[zero] (0, -\G) -- (0, \G);
  \node[enggray, anchor=south east] at (0, \G) {contour tangent};
  \node[engblue, anchor=north] at (0, -\G) {$D_{\unitvec{u}} f = 0$};
  \filldraw[engblack] (0,0) circle (0.05);
  \node[engblack, anchor=north west] at (0.05, -0.05) {$\vect{x}$};
\end{tikzpicture}
\caption[Directional derivative as a rose of rates.]{The directional
derivative $D_{\unitvec{u}} f = \norm{\grad f}\cos\theta$ read as a rose
of signed rates around the point $\vect{x}$, with $\grad f$ pointing
right (red). Each blue arrow's length is the rate of change of $f$ in
that heading. The rate is largest along $\grad f$ (steepest ascent),
falls off as the cosine of the angle away from it, and reaches zero
along the two headings tangent to the contour through $\vect{x}$ (grey
dashed). Headings into the left half-plane give negative rates (descent)
and are drawn as the mirror lengths; the single most negative heading is
directly opposite $\grad f$. The rose is the geometric statement of the
steepest-ascent theorem.}
\label{fig:vol02:ch11:directional-rose}
\end{figure}


\subsection{The total differential and the chain rule}

For $f : \R^{n} \to \R$, the local linear approximation around a
point $\vect{a}$ is
\[
f(\vect{a} + \dd\vect{x}) \approx f(\vect{a}) + \grad f(\vect{a})
\cdot \dd\vect{x}.
\]
The right-hand side is the \engterm{total differential} of $f$ at
$\vect{a}$. It generalises the one-variable linearisation $f(a +
\dd x) \approx f(a) + f'(a)\,\dd x$. The accuracy of the
approximation is to first order in $\norm{\dd\vect{x}}$, with error
of order $\norm{\dd\vect{x}}^{2}$ controlled by the Hessian of
section 11.3.

The multivariable chain rule generalises the one-variable rule. If
$\vect{x}(t) : \R \to \R^{n}$ is a path and $f : \R^{n} \to \R$, then
$g(t) = f(\vect{x}(t))$ satisfies
\[
\dv{g}{t} = \grad f(\vect{x}(t)) \cdot \dv{\vect{x}}{t}.
\]
For two variables: $\dv{g}{t} = f_{x}(x, y)\,\dv{x}{t} +
f_{y}(x, y)\,\dv{y}{t}$. The reader has already used the chain rule
implicitly in vector calculus (Volume II Chapter 8) when computing
$\dv{}{t}\Phi(\vect{r}(t)) = \grad\Phi \cdot \vect{v}$.

\paragraph{Worked example: the chain rule along a path.} A
temperature field is $T(x, y) = x^{2} - y^{2}$, and a sensor moves
along the path $x(t) = \cos t$, $y(t) = \sin t$ (a unit circle). The
rate of temperature change the sensor records is, by the chain rule,
\[
\dv{T}{t} = T_{x}\,\dot{x} + T_{y}\,\dot{y}
= (2x)(-\sin t) + (-2y)(\cos t)
= -2\cos t\sin t - 2\sin t\cos t = -4\sin t\cos t = -2\sin 2t.
\]
The reader can check this against the direct substitution $T(t) =
\cos^{2}t - \sin^{2}t = \cos 2t$, whose derivative is $-2\sin 2t$,
agreeing. The chain rule gives the answer without ever forming the
composite $\cos 2t$ explicitly, which matters when $T$ is known only
as data or as a black-box function. The recorded rate is the
gradient dotted with the velocity, $\grad T \cdot \vect{v}$, the
same object the steepest-ascent theorem reads.

\paragraph{Worked example: a chain rule with two intermediate
variables.} Let $w = f(u, v)$ with $u = u(s, t)$ and $v = v(s, t)$.
The chain rule with two intermediate and two independent variables
reads
\[
\pdv{w}{s} = \pdv{f}{u}\pdv{u}{s} + \pdv{f}{v}\pdv{v}{s},
\qquad
\pdv{w}{t} = \pdv{f}{u}\pdv{u}{t} + \pdv{f}{v}\pdv{v}{t}.
\]
Concretely, for $w = u^{2} + v^{2}$ with $u = s + t$, $v = s - t$,
\[
\pdv{w}{s} = 2u(1) + 2v(1) = 2(s + t) + 2(s - t) = 4s,
\qquad
\pdv{w}{t} = 2u(1) + 2v(-1) = 2(s + t) - 2(s - t) = 4t.
\]
The direct route confirms it: $w = (s + t)^{2} + (s - t)^{2} =
2s^{2} + 2t^{2}$, so $w_{s} = 4s$ and $w_{t} = 4t$. The bookkeeping
that the chain rule organises (one term per path from the output
through each intermediate variable to each input) is exactly the
matrix product $J^{\mathsf{T}}\grad_{\vect{u}} f$ written component
by component; the reader who draws the dependency as a tree and sums
over root-to-leaf paths is computing the matrix product by hand.

If instead $\vect{x}(\vect{u})$ is a function $\R^{m} \to \R^{n}$
and $f : \R^{n} \to \R$, then the composition $h(\vect{u}) =
f(\vect{x}(\vect{u}))$ has gradient
\[
\pdv{h}{u_{j}} = \sum_{i = 1}^{n} \pdv{f}{x_{i}}\,\pdv{x_{i}}{u_{j}},
\]
which can be written as a matrix-vector product $\grad_{\vect{u}} h
= J^{\mathsf{T}} \grad_{\vect{x}} f$, where $J$ is the $n \times m$
\engterm{Jacobian matrix} $J_{ij} = \pdv{x_{i}}{u_{j}}$. Section 11.7
takes up the Jacobian as the change-of-variable factor for multiple
integrals.

\paragraph{Worked example: the chain rule converts a Cartesian
derivative to polar.} A field known in Cartesian form is often sampled or
solved in polar coordinates, and the chain rule is the bridge. Let
$u(x, y)$ be a scalar field and write it in polar coordinates through
$x = r\cos\theta$, $y = r\sin\theta$. The chain rule gives the radial
derivative as
\[
\pdv{u}{r} = \pdv{u}{x}\pdv{x}{r} + \pdv{u}{y}\pdv{y}{r}
= u_{x}\cos\theta + u_{y}\sin\theta = \grad u \cdot \unitvec{r},
\]
which is the directional derivative along the outward radial unit vector
$\unitvec{r} = (\cos\theta, \sin\theta)$, as it must be. The angular
derivative is
\[
\pdv{u}{\theta} = u_{x}(-r\sin\theta) + u_{y}(r\cos\theta)
= r\,\grad u\cdot\unitvec{\theta},
\]
with $\unitvec{\theta} = (-\sin\theta, \cos\theta)$ the unit vector in the
direction of increasing angle, and the extra factor $r$ the arc-length
scaling of an angular step. Solving the two relations for the Cartesian
partials expresses $u_{x}$ and $u_{y}$ in terms of $u_{r}$ and
$u_{\theta}$,
\[
u_{x} = u_{r}\cos\theta - \frac{u_{\theta}}{r}\sin\theta,
\qquad
u_{y} = u_{r}\sin\theta + \frac{u_{\theta}}{r}\cos\theta,
\]
the transformation that rewrites any Cartesian first-order operator in
polar form. Applying it twice, and collecting terms, converts the
Cartesian Laplacian $u_{xx} + u_{yy}$ into its polar form
$u_{rr} + \tfrac{1}{r}u_{r} + \tfrac{1}{r^{2}}u_{\theta\theta}$, the
identity the heat and wave equations of Volume IV use whenever the
geometry is circular. The reader need not reproduce the two-step algebra
here; the point is that the chain rule alone, applied mechanically,
carries a differential operator from one coordinate system to another,
and the geometric readings ($u_{r} = \grad u\cdot\unitvec{r}$,
$u_{\theta}/r = \grad u\cdot\unitvec{\theta}$) are the audit that the
bookkeeping stayed correct.

\subsection{The tangent plane and linearisation}

The total differential has a geometric face. The graph $z = f(x, y)$
is a surface in $\R^{3}$; at a point $(a, b, f(a, b))$ the surface
has a \engterm{tangent plane}, the best linear approximation to the
surface near that point. The tangent plane is
\[
z = f(a, b) + f_{x}(a, b)\,(x - a) + f_{y}(a, b)\,(y - b),
\]
the affine function whose value and both first partials match $f$ at
$(a, b)$. \Cref{fig:vol02:ch11:tangent-plane} shows the construction:
the plane touches the surface at one point and departs from it by an
amount controlled, to leading order, by the Hessian of section 11.3.
The reader recognises the right-hand side as $f(a, b) + \grad f(a, b)
\cdot (\vect{x} - \vect{a})$, the total differential rewritten with
$\vect{a} = (a, b)$.

% Figure: the tangent plane to a surface z = f(x,y) at a point. A
% saddle-free upward-curving surface (paraboloid patch) with the
% flat tangent plane touching it at one point, illustrating the
% linearisation of section 11.2.
\begin{figure}[htbp]
\centering
\begin{tikzpicture}[scale=1.0,
  surf/.style={engblue, thick},
  plane/.style={engred, thick, fill=engred!10, fill opacity=0.5},
  axis/.style={engblack, thick, ->},
  label/.style={font=\small},
]
% Oblique projection axes (simple 3D-look in 2D)
\draw[axis] (0, 0) -- (4.6, 0) node[right, label] {$x$};
\draw[axis] (0, 0) -- ({1.4*cos(35)}, {1.4*sin(35)}) coordinate (yend);
\node[label] at ($(yend) + (0.15, 0.05)$) {$y$};
\draw[axis] (0, 0) -- (0, 3.6) node[above, label] {$z$};

% A few surface "ribs": cross-sections of an upward bowl
% z increases away from a central touch point at screen (2.2, 1.0)
\foreach \s in {0.0, 0.4, 0.8} {
  \draw[surf, opacity=0.55]
    plot[domain=0:3.6, samples=40]
    ({\x + \s*cos(35)}, {1.1 + 0.30*(\x-1.8)*(\x-1.8) + \s*sin(35)});
}

% Tangent plane: a parallelogram tangent at the bowl's low point.
% Touch point at screen coordinate roughly (1.8, 1.1).
\coordinate (T) at (1.8, 1.1);
\filldraw[engblack] (T) circle (0.05);
\node[label, anchor=north] at ($(T) + (0.0, -0.08)$) {$(a,b,f(a,b))$};

% Draw the flat tangent plane as a slanted parallelogram around T.
\draw[plane]
  ($(T) + (-1.2, -0.15)$) --
  ($(T) + (1.2, 0.15)$) --
  ($(T) + ({1.2 + 1.0*cos(35)}, {0.15 + 1.0*sin(35)})$) --
  ($(T) + ({-1.2 + 1.0*cos(35)}, {-0.15 + 1.0*sin(35)})$) -- cycle;
\node[label, engred, anchor=west] at ($(T) + (1.3, 0.55)$) {tangent plane};

\node[label, engblue, anchor=west] at (3.5, 2.9) {$z = f(x,y)$};

\end{tikzpicture}
\caption[Tangent plane to a surface.]{The tangent plane to the
surface $z = f(x, y)$ (blue) at the point $(a, b, f(a, b))$ (black
dot). The plane $z = f(a, b) + f_{x}(a, b)(x - a) + f_{y}(a, b)
(y - b)$ matches the surface's value and both first partials at the
touch point and is the best linear approximation nearby. The gap
between plane and surface grows quadratically in the displacement,
controlled by the Hessian; the linearisation is accurate to first
order in $\norm{\dd\vect{x}}$.}
\label{fig:vol02:ch11:tangent-plane}
\end{figure}


\paragraph{Worked example: linearising a square-root surface.}
Estimate $\sqrt{(3.02)^{2} + (3.97)^{2}}$ without a calculator. Let
$f(x, y) = \sqrt{x^{2} + y^{2}}$, and linearise about $(a, b) =
(3, 4)$ where $f(3, 4) = 5$ exactly. The partials are $f_{x} =
x/\sqrt{x^{2} + y^{2}}$ and $f_{y} = y/\sqrt{x^{2} + y^{2}}$, so at
$(3, 4)$ they are $f_{x} = 3/5 = 0.6$ and $f_{y} = 4/5 = 0.8$. The
linear approximation,
\[
f(3.02, 3.97) \approx 5 + 0.6\,(0.02) + 0.8\,(-0.03)
= 5 + 0.012 - 0.024 = 4.988,
\]
agrees with the exact value $4.98823\ldots$ to four significant
figures. The error of the linearisation, of order $\norm{\dd\vect{x}}
^{2} \approx (0.036)^{2} \approx 1.3 \times 10^{-3}$ scaled by the
surface curvature, is comfortably below the last reported digit.
This is the multivariable form of the engineer's habit of
estimating a small change by a derivative times a step, and it is
the foundation of error propagation in measurement (Volume I).

\paragraph{Worked example: error propagation through a formula.} A
resistor network gives total resistance $R(R_{1}, R_{2}) = R_{1}
R_{2} / (R_{1} + R_{2})$ for two resistors in parallel. With $R_{1}
= 100\,\Omega \pm 2\,\Omega$ and $R_{2} = 200\,\Omega \pm
4\,\Omega$, the nominal total is $R = (100)(200)/300 = 66.7\,\Omega$.
The partials, evaluated at the nominal point, give the sensitivity
of $R$ to each input. Differentiating,
\[
\pdv{R}{R_{1}} = \frac{R_{2}^{2}}{(R_{1} + R_{2})^{2}}
= \frac{200^{2}}{300^{2}} = 0.444, \qquad
\pdv{R}{R_{2}} = \frac{R_{1}^{2}}{(R_{1} + R_{2})^{2}}
= \frac{100^{2}}{300^{2}} = 0.111.
\]
Treating the tolerances as independent and summing in quadrature
(Volume I error-propagation rule), the uncertainty in $R$ is
\[
\delta R = \sqrt{\left(\pdv{R}{R_{1}}\delta R_{1}\right)^{2}
+ \left(\pdv{R}{R_{2}}\delta R_{2}\right)^{2}}
= \sqrt{(0.444 \cdot 2)^{2} + (0.111 \cdot 4)^{2}}
\approx 0.99\,\Omega.
\]
The total resistance is $66.7\,\Omega \pm 1.0\,\Omega$. The gradient
of the formula is the engine of the error budget: the partial with
respect to $R_{1}$ is four times the partial with respect to $R_{2}$,
so the smaller resistor, despite its smaller absolute tolerance,
dominates the uncertainty. The reader who wants to tighten the
output tolerance buys a better $R_{1}$ first.

\paragraph{Worked example: implicit differentiation of a level curve.}
A curve is given implicitly by $F(x, y) = x^{3} + y^{3} - 6 x y = 0$
(the folium of Descartes). The slope $\dv{y}{x}$ of the curve at a point
is available without solving for $y$, through the gradient. Differentiate
the constraint $F = 0$ along the curve: by the chain rule,
$F_{x} + F_{y}\dv{y}{x} = 0$, so
\[
\dv{y}{x} = -\frac{F_{x}}{F_{y}}
= -\frac{3 x^{2} - 6 y}{3 y^{2} - 6 x}
= -\frac{x^{2} - 2 y}{y^{2} - 2 x},
\]
valid wherever $F_{y} = 3 y^{2} - 6 x \neq 0$. At the point $(3, 3)$,
which lies on the curve ($27 + 27 - 54 = 0$), the slope is
$-(9 - 6)/(9 - 6) = -1$: the tangent runs at $45^{\circ}$ down to the
right. The rule $\dv{y}{x} = -F_{x}/F_{y}$ is the implicit-function
theorem of the Lagrange mastery box in miniature, and it says the curve's
tangent is perpendicular to $\grad F = (F_{x}, F_{y})$, because a level
curve of $F$ is everywhere orthogonal to $\grad F$. The denominator's
vanishing, $F_{y} = 0$, is exactly where the tangent turns vertical and
the local description $y = y(x)$ breaks down, the point at which the
roles of $x$ and $y$ must swap. The engineer who tracks a design
constraint as a level curve reads its local trade-off slope directly
from the ratio of partials, never solving the implicit relation.

\begin{archetype}[Optimisation, multivariable form]
The single-variable optimisation archetype of Volume II Chapter 5
asks for points where $f'(x) = 0$. The multivariable form asks for
points where $\grad f(\vect{x}) = \vect{0}$, a system of $n$ scalar
equations in $n$ unknowns. The condition is necessary for an
unconstrained interior maximum or minimum (the analogue of Fermat's
condition), and the second-derivative test of section 11.3 supplies
sufficiency. The archetype recurs across Volume III (mechanical
equilibria as critical points of potential energy), Volume IV
(thermodynamic equilibria as critical points of free energy),
Volume VII (loss-function minima in machine learning), and
Volume IX (control objectives and minimum-cost design).
\end{archetype}

\begin{estimation}
\textbf{Question.} A rectangular field is to be enclosed by a fence
on three sides only (a river forms the fourth side). The total
length of fence is $L = 200\,\si{\meter}$. Estimate the dimensions
that maximise the area, then compute exactly.

\textbf{Estimate.} The area $A = x y$ where $x$ is the side parallel
to the river and $y$ is each of the two perpendicular sides; the
fence constraint is $x + 2 y = L$. A square (i.e.\ $x = y$) gives
$3 x = 200$, $x = y \approx 67\,\si{\meter}$, $A \approx
4500\,\si{\meter\squared}$. The reader's intuition for closed
rectangles is that the square is optimal; for the open-on-one-side
case the optimum is biased toward a longer river-side. Best estimate
before calculation: $x \approx 100\,\si{\meter}$, $y \approx
50\,\si{\meter}$, $A \approx 5000\,\si{\meter\squared}$.

\textbf{Calculate.} Eliminate the constraint: $x = L - 2 y$, so
$A(y) = (L - 2 y) y = L y - 2 y^{2}$. The critical point satisfies
$A'(y) = L - 4 y = 0$, giving $y^{*} = L/4 = 50\,\si{\meter}$,
$x^{*} = L/2 = 100\,\si{\meter}$, $A^{*} = L^{2}/8 =
5000\,\si{\meter\squared}$. The estimate landed on the answer. The
multivariable form of this problem (without elimination) is the
worked example of the next section.
\end{estimation}

\section{Critical points: maxima, minima, saddles}\pathtag{core}

A point where $\grad f = \vect{0}$ is a \engterm{critical point} of
$f$. As in one variable, the critical-point condition is necessary
but not sufficient: a critical point may be a local maximum, a
local minimum, a saddle, or a degenerate point at which the second-
derivative test cannot decide.

\subsection{The Hessian}

\begin{definition}[Hessian]
The \engterm{Hessian} of $f : \R^{n} \to \R$ at $\vect{x}$ is the
$n \times n$ matrix of second partial derivatives,
\[
H(\vect{x}) = \begin{pmatrix}
f_{x_{1} x_{1}} & f_{x_{1} x_{2}} & \cdots & f_{x_{1} x_{n}} \\
f_{x_{2} x_{1}} & f_{x_{2} x_{2}} & \cdots & f_{x_{2} x_{n}} \\
\vdots & \vdots & \ddots & \vdots \\
f_{x_{n} x_{1}} & f_{x_{n} x_{2}} & \cdots & f_{x_{n} x_{n}}
\end{pmatrix}.
\]
\end{definition}

By Clairaut's theorem, $H$ is symmetric whenever $f$ has continuous
second partials. The Hessian is the multivariable generalisation of
the second derivative $f''$.

\subsection{The second-order Taylor expansion}

Around a critical point $\vect{x}^{*}$, the second-order Taylor
expansion of $f$ is
\[
f(\vect{x}^{*} + \dd\vect{x}) = f(\vect{x}^{*}) +
\tfrac{1}{2}\,\dd\vect{x}^{\mathsf{T}} H(\vect{x}^{*})\,\dd\vect{x}
+ \oom{\norm{\dd\vect{x}}^{3}},
\]
because the linear term $\grad f(\vect{x}^{*}) \cdot \dd\vect{x}$
vanishes at a critical point. The behaviour of $f$ near $\vect{x}^{*}$
is therefore controlled, to leading order, by the quadratic form
$\dd\vect{x}^{\mathsf{T}} H(\vect{x}^{*})\,\dd\vect{x}$.

A symmetric matrix $H$ has $n$ real eigenvalues $\lambda_{1}, \ldots,
\lambda_{n}$ (Volume II Chapter 10). The signs of those eigenvalues
classify the critical point.

\begin{theorem}[Second-derivative test, multivariable]
Let $f$ have continuous second partials and let $\vect{x}^{*}$ be a
critical point of $f$. Let $H = H(\vect{x}^{*})$.
\begin{itemize}[leftmargin=*]
  \item If every eigenvalue of $H$ is positive, $\vect{x}^{*}$ is a
    local minimum.
  \item If every eigenvalue of $H$ is negative, $\vect{x}^{*}$ is a
    local maximum.
  \item If $H$ has at least one positive and at least one negative
    eigenvalue, $\vect{x}^{*}$ is a saddle point.
  \item If $H$ has any zero eigenvalue and the others are all of
    one sign, the test is inconclusive: the next-higher-order Taylor
    term must be inspected.
\end{itemize}
\end{theorem}

For $n = 2$ the test takes a more memorable form. With $H =
\begin{pmatrix} A & B \\ B & C \end{pmatrix}$, the determinant is
$\det H = AC - B^{2}$. The eigenvalues are positive when $\det H > 0$
and $A > 0$, negative when $\det H > 0$ and $A < 0$, and of mixed
sign when $\det H < 0$.

\begin{itemize}[leftmargin=*]
  \item $\det H > 0$ and $A > 0$: local minimum.
  \item $\det H > 0$ and $A < 0$: local maximum.
  \item $\det H < 0$: saddle.
  \item $\det H = 0$: inconclusive.
\end{itemize}

\begin{mastery}
The second-derivative test is the sign reading of the quadratic
form $q(\dd\vect{x}) = \dd\vect{x}^{\mathsf{T}} H\,\dd\vect{x}$, and
the eigenvalue criterion of the theorem is its derivation. Diagonalise
the symmetric Hessian by the spectral theorem of Volume II
Chapter 10: $H = Q\Lambda Q^{\mathsf{T}}$ with $Q$ orthogonal and
$\Lambda = \diag(\lambda_{1}, \ldots, \lambda_{n})$. In the rotated
coordinates $\vect{y} = Q^{\mathsf{T}}\dd\vect{x}$, the form
separates,
\[
q = \vect{y}^{\mathsf{T}}\Lambda\,\vect{y} =
\sum_{i = 1}^{n}\lambda_{i}\,y_{i}^{2}.
\]
If every $\lambda_{i} > 0$, then $q > 0$ for every nonzero
displacement, so $f$ rises in all directions and $\vect{x}^{*}$ is a
strict local minimum. If every $\lambda_{i} < 0$, $q < 0$ in all
directions, a strict local maximum. If the eigenvalues have mixed
sign, $q > 0$ along the eigenvectors of positive eigenvalue and
$q < 0$ along those of negative eigenvalue, a saddle. A zero
eigenvalue leaves $q = 0$ along its eigenvector, and the
second-order term decides nothing in that direction, which is the
degenerate case the failure section takes up. The two-variable
$\det H = AC - B^{2}$ test is the same statement in disguise: for a
$2\times 2$ symmetric matrix $\det H = \lambda_{1}\lambda_{2}$ and
$\tr H = A + C = \lambda_{1} + \lambda_{2}$, so $\det H > 0$ forces
the eigenvalues to share a sign (fixed by that of $A = \tr H -
C$), $\det H < 0$ forces opposite signs, and $\det H = 0$ forces a
zero eigenvalue. The determinant-and-trace shortcut and the
eigenvalue criterion are one test read two ways.
\end{mastery}

\Cref{fig:vol02:ch11:critical-classification} shows the four cases
side by side. The local minimum and the local maximum have nested
closed contours indistinguishable without the sign of $f$; the
saddle has hyperbolic contours that exchange the two principal
directions; the degenerate case has parallel-line contours that
report the loss of one Hessian eigenvalue. The reader learns to
read a contour plot for the type of critical point at a glance.

% Figure: classification of critical points by Hessian eigenvalues.
% Four panels: positive-definite (bowl), negative-definite (dome),
% indefinite (saddle), semi-definite (degenerate trough).
\begin{figure}[htbp]
\centering
\begin{tikzpicture}[scale=0.8,
  contour/.style={engblue, thick},
  panel/.style={draw=engblack, thick},
  caption/.style={font=\small, align=center},
  star/.style={engred},
]

% ---- Panel 1: positive definite (local min): f = x^2 + y^2 ----
\begin{scope}[shift={(0, 0)}]
  \draw[panel] (-1.6, -1.6) rectangle (1.6, 1.6);
  \foreach \r in {0.4, 0.8, 1.2} {
    \draw[contour] (0, 0) circle (\r);
  }
  \filldraw[star] (0,0) circle (0.06);
  \node[caption] at (0, -2.05) {min\\$\lambda_{1}, \lambda_{2} > 0$};
\end{scope}

% ---- Panel 2: negative definite (local max): f = -x^2 - y^2 ----
\begin{scope}[shift={(4, 0)}]
  \draw[panel] (-1.6, -1.6) rectangle (1.6, 1.6);
  \foreach \r in {0.4, 0.8, 1.2} {
    \draw[contour] (0, 0) circle (\r);
  }
  \filldraw[star] (0,0) circle (0.06);
  \node[caption] at (0, -2.05) {max\\$\lambda_{1}, \lambda_{2} < 0$};
\end{scope}

% ---- Panel 3: indefinite (saddle): f = x^2 - y^2 ----
\begin{scope}[shift={(8, 0)}]
  \draw[panel] (-1.6, -1.6) rectangle (1.6, 1.6);
  % Hyperbolae x^2 - y^2 = c for c = +/- 0.25, +/- 0.81, +/- 1.44
  % Restrict domain to real-valued branches: |x| >= sqrt(c) for the y(x) branches
  \foreach \c [evaluate=\c as \r using {sqrt(\c)}] in {0.25, 0.81, 1.44} {
    \draw[contour, domain=\r:1.55, samples=40]
      plot ({\x}, {sqrt(\x*\x - \c)});
    \draw[contour, domain=\r:1.55, samples=40]
      plot ({\x}, {-sqrt(\x*\x - \c)});
    \draw[contour, domain=-1.55:-\r, samples=40]
      plot ({\x}, {sqrt(\x*\x - \c)});
    \draw[contour, domain=-1.55:-\r, samples=40]
      plot ({\x}, {-sqrt(\x*\x - \c)});
    \draw[contour, domain=-1.55:1.55, samples=80]
      plot ({sqrt(\x*\x + \c)}, {\x});
    \draw[contour, domain=-1.55:1.55, samples=80]
      plot ({-sqrt(\x*\x + \c)}, {\x});
  }
  % Asymptotes
  \draw[engblack!50] (-1.5, -1.5) -- (1.5, 1.5);
  \draw[engblack!50] (-1.5, 1.5) -- (1.5, -1.5);
  \filldraw[star] (0,0) circle (0.06);
  \node[caption] at (0, -2.05) {saddle\\$\lambda_{1} > 0,\ \lambda_{2} < 0$};
\end{scope}

% ---- Panel 4: semi-definite (degenerate trough): f = x^2 ----
\begin{scope}[shift={(12, 0)}]
  \draw[panel] (-1.6, -1.6) rectangle (1.6, 1.6);
  % Vertical lines x = +/- sqrt(c) for c = 0.25, 1, 1.96 (but 1.96 > 1.6 so skip)
  \foreach \c in {0.25, 0.81, 1.44} {
    \draw[contour] ({sqrt(\c)}, -1.55) -- ({sqrt(\c)}, 1.55);
    \draw[contour] ({-sqrt(\c)}, -1.55) -- ({-sqrt(\c)}, 1.55);
  }
  \draw[engblack!60, thick, dashed] (0, -1.55) -- (0, 1.55);
  \filldraw[star] (0,0) circle (0.06);
  \node[caption] at (0, -2.05) {degenerate\\$\lambda_{1} > 0,\ \lambda_{2} = 0$};
\end{scope}

\end{tikzpicture}
\caption[Critical points classified by the Hessian
eigenvalues.]{Four canonical critical-point types in two variables,
shown as contour plots near the origin. From left: a local minimum
(both Hessian eigenvalues positive, contours nested closed curves
around the critical point); a local maximum (both eigenvalues
negative, geometrically indistinguishable from the minimum without
the sign of the function values); a saddle (one eigenvalue of each
sign, contours hyperbolae with the critical point at the asymptote
intersection); a degenerate critical point (one zero eigenvalue,
contours degenerate to parallel lines, the second-derivative test
inconclusive). The red dot marks the critical point in each panel.}
\label{fig:vol02:ch11:critical-classification}
\end{figure}


\subsection{Worked example: the saddle}

Let $f(x, y) = x^{2} - y^{2}$. Then $\grad f = (2x, -2y)^{\mathsf{T}}$
and the only critical point is the origin. The Hessian is $H =
\diag(2, -2)$, with eigenvalues $2$ and $-2$. The origin is a
saddle: $f$ increases along the $x$-direction and decreases along
the $y$-direction. The reader has met this surface as a horse-saddle
or a Pringles potato chip.

\subsection{Worked example: three critical points and the full
test}

Let $f(x, y) = x^{3} - 3 x + y^{2}$. The gradient is $\grad f =
(3 x^{2} - 3, 2 y)^{\mathsf{T}}$. The condition $\grad f = \vect{0}$
gives $x^{2} = 1$ and $y = 0$, so two critical points, $(1, 0)$ and
$(-1, 0)$. The Hessian is
\[
H = \begin{pmatrix} 6 x & 0 \\ 0 & 2 \end{pmatrix}.
\]
At $(1, 0)$, $H = \diag(6, 2)$: both eigenvalues positive,
$\det H = 12 > 0$ with $A = 6 > 0$, a local minimum. At $(-1, 0)$,
$H = \diag(-6, 2)$: eigenvalues of opposite sign, $\det H = -12 <
0$, a saddle. The reader reads the geometry directly from the
function. The $x^{3} - 3x$ term is a one-variable cubic with a local
minimum at $x = 1$ and a local maximum at $x = -1$; the $+y^{2}$ term
adds an upward bowl in the $y$-direction. At $x = 1$ the function
curves up in both directions (a minimum); at $x = -1$ it curves down
in $x$ and up in $y$ (a saddle). The separable structure makes the
Hessian diagonal and the classification immediate.

\subsection{Worked example: the second-order Taylor expansion in
use}

The classification theorem is the sign reading of the quadratic form
$\dd\vect{x}^{\mathsf{T}} H\,\dd\vect{x}$. Make the form explicit for
$f(x, y) = 2 x^{2} + 2 x y + 3 y^{2}$ at its critical point. The
gradient $\grad f = (4 x + 2 y, 2 x + 6 y)^{\mathsf{T}}$ vanishes
only at the origin. The Hessian is constant,
\[
H = \begin{pmatrix} 4 & 2 \\ 2 & 6 \end{pmatrix},
\]
with $\det H = 24 - 4 = 20 > 0$ and $A = 4 > 0$, a local minimum.
Because $f$ is itself a pure quadratic, the second-order Taylor
expansion about the origin is exact: $f(\dd\vect{x}) = \tfrac{1}{2}
\dd\vect{x}^{\mathsf{T}} H\,\dd\vect{x}$, and the reader can verify
$\tfrac{1}{2}(4 x^{2} + 2 \cdot 2 x y + 6 y^{2}) = 2 x^{2} + 2 x y +
3 y^{2}$, recovering $f$. The eigenvalues of $H$ are $\lambda =
5 \pm \sqrt{1 + 4} = 5 \pm \sqrt{5}$, both positive (approximately
$7.24$ and $2.76$), confirming the minimum and reporting the
curvatures along the two principal axes. The level sets $f = c$ are
ellipses whose axes align with the Hessian's eigenvectors and whose
axis lengths scale as $1/\sqrt{\lambda}$: the steep direction
(large $\lambda$) gives a short axis, the gentle direction a long
one. The reader who has read Volume II Chapter 10 recognises the
principal-axis theorem for quadratic forms; the Hessian's
eigendecomposition is the principal-axis decomposition of the local
shape of $f$.

\begin{mastery}
The condition number $\kappa = \lambda_{\max}/\lambda_{\min}$ of the
Hessian at a minimum measures the eccentricity of the local level-set
ellipses and predicts the difficulty of gradient descent. For the
example above, $\kappa = (5 + \sqrt 5)/(5 - \sqrt 5) \approx 2.62$, a
well-conditioned bowl. The Rosenbrock function of the simulation
exercises has $\kappa \approx 172$ at its minimum, a long thin
valley that defeats naive descent. A linear change of variables that
diagonalises $H$ and rescales each axis by $\sqrt{\lambda}$ turns
any non-degenerate minimum into a perfect circular bowl with
$\kappa = 1$; this is the geometric content of \engterm{Newton's
method}, which premultiplies the gradient by $H^{-1}$ to take the
ill-conditioning out of the step. Volume II Chapter 15 develops the
algorithm; the geometry is that Newton's method works in the
coordinate system where every minimum looks the same.
\end{mastery}

\subsection{Worked example: one Newton step against one gradient step}

The condition-number remark above is worth making concrete with a
single step of each method on a bowl steep in one direction and
gentle in the other. Take $f(x, y) = \tfrac{1}{2}x^{2} +
\tfrac{5}{2}y^{2}$, whose minimum is the origin, with gradient
$\grad f = (x, 5 y)^{\mathsf{T}}$ and constant Hessian $H =
\diag(1, 5)$. Start at $(4, 1)$, where $\grad f = (4, 5)^{\mathsf{T}}$.

Gradient descent moves along $-\grad f$. A step size $\eta$ small
enough to keep the steep $y$-direction stable, $\eta < 2/\lambda_{
\max} = 2/5 = 0.4$, gives at $\eta = 0.18$ the update $(4, 1) -
0.18\,(4, 5) = (3.28, 0.10)$: the $y$-coordinate collapses almost to
zero while $x$ barely moves, and the iterate must then crawl along
the $x$-axis in many further steps. Newton's method instead
premultiplies the gradient by $H^{-1} = \diag(1, \tfrac{1}{5})$,
\[
\vect{x}_{1} = \vect{x}_{0} - H^{-1}\grad f
= (4, 1) - \Big(1\cdot 4,\ \tfrac{1}{5}\cdot 5\Big)
= (4, 1) - (4, 1) = (0, 0),
\]
landing on the minimum in one move. The reason is exact for a
quadratic: the second-order Taylor model $f(\vect{x}_{0} +
\dd\vect{x}) = f(\vect{x}_{0}) + \grad f\cdot\dd\vect{x} +
\tfrac{1}{2}\dd\vect{x}^{\mathsf{T}} H\,\dd\vect{x}$ is the whole
function, so the step that zeroes the model's gradient, $\dd\vect{x}
= -H^{-1}\grad f$, zeroes the true gradient too. The inverse Hessian
rescales each axis by its own curvature, undoing the anisotropy that
the single scalar $\eta$ of gradient descent cannot.
\Cref{fig:vol02:ch11:newton-step} draws the two steps side by side.
The price of the Newton step is the $H^{-1}$ solve, $\oom{n^{3}}$ per
step in $n$ variables, against the $\oom{n}$ of a gradient step; the
working trade between the two, and the quasi-Newton methods that
approximate $H^{-1}$ cheaply, is the subject of Volume II Chapter 15.

% Figure: one gradient-descent step versus one Newton step on an
% anisotropic quadratic bowl f(x,y) = x^2/2 + 5 y^2/2, whose minimum is the
% origin. From the start (4,1) the negative gradient (-4,-5) points across
% the valley, so a gradient step overshoots the steep y-direction and
% undershoots the gentle x-direction; the Newton step -H^{-1} grad f lands
% on the minimum in one move because H^{-1} rescales each axis by its
% curvature.
\begin{figure}[htbp]
\centering
\begin{tikzpicture}[scale=1.0,
  contour/.style={enggray, thin},
  gdstep/.style={engred, very thick, -{Stealth[length=3mm]}},
  newton/.style={engblue, very thick, -{Stealth[length=3mm]}},
  pt/.style={circle, fill=engblack, inner sep=1.4pt},
]
  \begin{axis}[
    width=9.6cm, height=6.6cm,
    xlabel={$x$}, ylabel={$y$},
    xmin=-1.2, xmax=5, ymin=-1.6, ymax=1.8,
    axis lines=middle,
    xtick={-1,0,...,5}, ytick={-1,0,1},
    tick label style={font=\small},
    label style={font=\small},
    clip=false,
  ]
    % Elliptical contours of f = x^2/2 + 5 y^2/2 for a few levels.
    \foreach \c in {1,4,9} {
      \addplot[contour, domain=0:360, samples=80]
        ({sqrt(2*\c)*cos(x)}, {sqrt(2*\c/5)*sin(x)});
    }
    % Start point (4,1).
    \node[pt] at (axis cs:4,1) {};
    \node[anchor=south west, font=\small] at (axis cs:4,1) {start $(4,1)$};
    % Gradient descent step with eta = 0.18: grad = (4,5), step = (-0.72,-0.9)
    % lands at (3.28, 0.1). Draw the arrow.
    \draw[gdstep] (axis cs:4,1) -- (axis cs:3.28,0.1);
    \node[engred, anchor=west, font=\small] at (axis cs:3.35,0.55)
      {one GD step};
    % Newton step lands exactly at origin.
    \draw[newton] (axis cs:4,1) -- (axis cs:0,0);
    \node[engblue, anchor=south, font=\small] at (axis cs:2.0,0.62)
      {one Newton step};
    % Minimum marker.
    \node[pt] at (axis cs:0,0) {};
    \node[anchor=north east, font=\small] at (axis cs:0,-0.05) {minimum};
  \end{axis}
\end{tikzpicture}
\caption[One gradient step versus one Newton step on an anisotropic
bowl.]{The bowl $f(x, y) = \tfrac{1}{2}x^{2} + \tfrac{5}{2}y^{2}$ has
elliptical contours, five times steeper along $y$ than along $x$. From the
start $(4, 1)$ the negative gradient $(-4, -5)$ points mostly across the
valley, so a single gradient-descent step (red) drops the $y$-coordinate
almost to zero while barely reducing $x$, and the iterate must then crawl
along the $x$-axis. The Newton step $-H^{-1}\grad f$ (blue) premultiplies
the gradient by the inverse curvature and lands on the minimum in one move,
because $H^{-1}$ rescales each axis by its own steepness. The picture is the
geometry behind the condition-number remark: Newton's method works in the
coordinate system where every bowl is round.}
\label{fig:vol02:ch11:newton-step}
\end{figure}


\subsection{Worked example: a Taylor expansion of a transcendental
surface}

The pure-quadratic example above had an exact second-order
expansion because the function carried no higher terms. The
generic case keeps a remainder, and the worked discipline is to
form the gradient, the Hessian, and the quadratic correction in
order. Expand $f(x, y) = e^{x}\cos y$ about the point $\vect{a} =
(0, 0)$ to second order. The function value is $f(0, 0) = 1$. The
first partials are $f_{x} = e^{x}\cos y$ and $f_{y} = -e^{x}\sin y$,
so $\grad f(0, 0) = (1, 0)^{\mathsf{T}}$. The second partials are
\[
f_{xx} = e^{x}\cos y, \quad
f_{xy} = -e^{x}\sin y, \quad
f_{yy} = -e^{x}\cos y,
\]
which at the origin give the Hessian
\[
H(0, 0) = \begin{pmatrix} 1 & 0 \\ 0 & -1 \end{pmatrix}.
\]
The mixed partial $f_{xy}(0,0) = 0$ confirms Clairaut and tells the
reader the two coordinate directions decouple at this point. The
second-order Taylor polynomial is
\[
f(x, y) \approx f(0, 0) + \grad f(0, 0)\cdot(x, y) +
\tfrac{1}{2}(x, y)\,H(0, 0)\,(x, y)^{\mathsf{T}}
= 1 + x + \tfrac{1}{2}(x^{2} - y^{2}).
\]
The same polynomial follows from multiplying the one-variable
series $e^{x} = 1 + x + \tfrac{x^{2}}{2} + \cdots$ and $\cos y =
1 - \tfrac{y^{2}}{2} + \cdots$ and keeping terms through total
degree two: $(1 + x + \tfrac{x^{2}}{2})(1 - \tfrac{y^{2}}{2}) =
1 + x + \tfrac{x^{2}}{2} - \tfrac{y^{2}}{2}$ to second order, the
same result. The two routes agreeing is the multivariable analogue
of the one-variable consistency check. The origin is not a critical
point of $f$ (the gradient $(1, 0)$ is nonzero), so the Hessian's
mixed-sign eigenvalues $\pm 1$ do not classify a saddle here; they
report the local curvature of the surface, upward along $x$ and
downward along $y$, around a sloping tangent plane. A numerical
spot-check at $(0.1, 0.1)$: the exact value is $e^{0.1}\cos(0.1) =
1.1052 \times 0.9950 = 1.0997$, and the quadratic model gives
$1 + 0.1 + \tfrac{1}{2}(0.01 - 0.01) = 1.1000$, agreeing to four
significant figures, with the small discrepancy carried by the
third-order term the model dropped.

\subsection{Worked example: the constrained-fence problem
revisited}

Let $f(x, y) = xy + \lambda(x + 2y - L)$ where $\lambda$ is a
parameter to be set. Setting $\grad_{(x, y)} f = 0$ gives
$y + \lambda = 0$ and $x + 2\lambda = 0$, so $x = -2\lambda$ and
$y = -\lambda$. The constraint $x + 2 y = L$ then forces
$-2\lambda - 2\lambda = L$, $\lambda = -L/4$, $x = L/2$, $y = L/4$,
the same answer as the eliminated form of section 11.2.6. The
extra structure of section 11.4 makes this trick general.

\subsection{Worked example: a critical point in three variables}

The two-variable determinant test does not extend past $n = 2$; in three
or more variables the classification returns to the eigenvalue signs of
the theorem. Classify the critical points of
\[
f(x, y, z) = x^{2} + y^{2} + z^{2} - x y - y z + x - z.
\]
The gradient is
\[
\grad f = (2 x - y + 1,\; 2 y - x - z,\; 2 z - y - 1)^{\mathsf{T}}.
\]
Setting each component to zero gives the linear system $2x - y = -1$,
$-x + 2y - z = 0$, $-y + 2z = 1$. Adding the first and third,
$2x + 2z - 2y = 0$, so $x + z = y$; substituting into the middle
equation, $-x + 2y - z = -y + 2y = y = 0$. Hence $y = 0$, and the outer
equations give $x = -\tfrac{1}{2}$, $z = \tfrac{1}{2}$. The single
critical point is $(-\tfrac{1}{2}, 0, \tfrac{1}{2})$. The Hessian is
constant,
\[
H = \begin{pmatrix} 2 & -1 & 0 \\ -1 & 2 & -1 \\ 0 & -1 & 2
\end{pmatrix}.
\]
This is the tridiagonal matrix of the discrete second-difference
operator, whose eigenvalues are known in closed form,
$\lambda_{k} = 2 - 2\cos\!\big(k\pi/4\big)$ for $k = 1, 2, 3$, giving
$\lambda_{1} = 2 - \sqrt{2} \approx 0.59$, $\lambda_{2} = 2$, and
$\lambda_{3} = 2 + \sqrt{2} \approx 3.41$. All three eigenvalues are
positive, so the critical point is a local (and, since $f$ is a convex
quadratic, global) minimum. The reader who prefers not to quote the
eigenvalue formula reaches the same verdict from the leading principal
minors, $2 > 0$, $\det\begin{psmallmatrix} 2 & -1 \\ -1 & 2
\end{psmallmatrix} = 3 > 0$, and $\det H = 2(4 - 1) - (-1)(-2 - 0) =
6 - 2 = 4 > 0$: all positive is Sylvester's criterion for
positive-definiteness, the $n$-variable replacement for the
determinant-and-trace shortcut of two variables. The condition number
$\lambda_{3}/\lambda_{1} = (2 + \sqrt 2)/(2 - \sqrt 2) \approx 5.8$
reports a mildly elongated bowl, and the coupling terms $-xy$ and $-yz$
are what tilt its principal axes away from the coordinate directions.

\subsection{Worked example: a non-diagonal Hessian and the full test}

The diagonal Hessians of the separable examples above hand the
classification over for free. The generic case has off-diagonal terms,
and the $\det H = AC - B^{2}$ test earns its keep. Classify the
critical points of $f(x, y) = x^{3} + y^{3} - 3 x y$. The gradient is
$\grad f = (3 x^{2} - 3 y, 3 y^{2} - 3 x)^{\mathsf{T}}$, and setting
both components to zero gives $y = x^{2}$ and $x = y^{2}$. Substituting
the first into the second, $x = x^{4}$, so $x(x^{3} - 1) = 0$ and the
real roots are $x = 0$ and $x = 1$, giving the two critical points
$(0, 0)$ and $(1, 1)$. The Hessian carries a nonzero off-diagonal,
\[
H = \begin{pmatrix} 6 x & -3 \\ -3 & 6 y \end{pmatrix},
\qquad \det H = 36 x y - 9.
\]
At $(0, 0)$, $H = \begin{psmallmatrix} 0 & -3 \\ -3 & 0
\end{psmallmatrix}$ with $\det H = -9 < 0$: a saddle, eigenvalues
$\pm 3$. The off-diagonal alone decides it; a reader who glanced only
at the zero diagonal and guessed "degenerate" would be wrong, because
the coupling term supplies the mixed-sign curvature. At $(1, 1)$,
$H = \begin{psmallmatrix} 6 & -3 \\ -3 & 6 \end{psmallmatrix}$ with
$\det H = 36 - 9 = 27 > 0$ and $A = 6 > 0$: a local minimum,
eigenvalues $6 \pm 3 = 9$ and $3$. The two eigenvectors, along $(1, 1)$
and $(1, -1)$, are the principal axes of the local bowl, rotated
$45^{\circ}$ from the coordinate axes by the $-3xy$ coupling exactly as
the three-variable example's tilt arose from its coupling. The lesson:
the off-diagonal entry is not a nuisance term to be dropped but the
carrier of the interaction between the two variables, and only the
determinant, which weighs it against the diagonal product, reads the
curvature correctly.

\subsection{Why saddle points matter for engineering}

A saddle point is the structural feature that defeats naive
optimisation. A gradient-descent iterate at exactly a saddle is
stationary; a gradient-descent iterate near a saddle is misdirected.
In low dimensions ($n \leq 3$) saddles are rare in practice; in
high dimensions they are not. In the loss landscapes of large
machine-learning models, the critical points encountered by training
are overwhelmingly saddle points rather than local minima
\cite{paper:v2ch11-dauphin-saddle-2014}. The working consequence:
an optimiser that does not handle saddle points gracefully (by
adding noise, by following negative-curvature directions, by using
second-order information) gets stuck. Volume II Chapter 15 develops
the algorithmic responses; this chapter supplies the geometry.

\section{Lagrange multipliers}\pathtag{core}

The unconstrained critical-point condition $\grad f = \vect{0}$
solves the problem of finding extrema of $f$ on an open region. Most
engineering problems sit on a constraint: maximise the volume
subject to a fixed surface area, minimise the cost subject to a
production target, design a profile subject to a shape requirement.
The constraint reduces the dimension of the feasible set and changes
the critical-point condition.

\subsection{The geometric idea}

We seek the extrema of $f(\vect{x})$ subject to one equality
constraint $g(\vect{x}) = 0$. The constraint defines a surface (or
curve, or hypersurface) in $\R^{n}$. At an extremum of $f$ on this
surface, $f$ cannot increase or decrease in any direction tangent
to the surface, otherwise we could move along the surface to a
higher (or lower) value. The directions tangent to $g = 0$ at a
point $\vect{x}^{*}$ are exactly the directions $\dd\vect{x}$ that
satisfy $\grad g(\vect{x}^{*}) \cdot \dd\vect{x} = 0$, because
$g(\vect{x}^{*} + \dd\vect{x}) \approx g(\vect{x}^{*}) +
\grad g \cdot \dd\vect{x} = 0$ to first order.

The directional derivative of $f$ in any tangent direction must
vanish: $\grad f \cdot \dd\vect{x} = 0$ for every $\dd\vect{x}$ with
$\grad g \cdot \dd\vect{x} = 0$. The two gradients $\grad f$ and
$\grad g$ at the constrained extremum are therefore parallel: there
is a scalar $\lambda$ such that
\[
\grad f(\vect{x}^{*}) = \lambda\,\grad g(\vect{x}^{*}).
\]
The scalar $\lambda$ is the \engterm{Lagrange multiplier}.
\Cref{fig:vol02:ch11:lagrange-tangency} shows the geometry on the
problem of maximising $f(x, y) = x + y$ on the circle $x^{2} + y^{2}
= 4$: the optimum is the point where a level line of $f$ (a line
of slope $-1$) is tangent to the circle, the two gradients pointing
in the same direction.

% Figure: Lagrange-multiplier geometry. Level curves of f(x,y) and
% the constraint curve g(x,y)=0; at the constrained extremum the
% level curve is tangent to the constraint, i.e. grad f parallel to
% grad g.
\begin{figure}[htbp]
\centering
\begin{tikzpicture}[scale=1.0,
  fcurve/.style={engblue, thick},
  gcurve/.style={engred, very thick},
  gradf/.style={->, engblue, very thick},
  gradg/.style={->, engred, very thick},
  axis/.style={engblack, thick, ->},
  label/.style={font=\small},
]
% Axes
\draw[axis] (-0.3, 0) -- (4.3, 0) node[right, label] {$x$};
\draw[axis] (0, -0.3) -- (0, 3.3) node[above, label] {$y$};

% Level curves of f(x,y) = x + y = c, three values
\foreach \c in {1.0, 2.0, 2.8, 3.6} {
  \draw[fcurve, opacity=0.7] (\c, 0) -- (0, \c);
}
\node[label, engblue] at (3.5, 0.18) {$f = c$};

% Constraint curve g(x,y) = x^2 + y^2 - 4 = 0  (circle radius 2)
% Show the first-quadrant arc
\draw[gcurve, domain=0:90, samples=50] plot ({2*cos(\x)}, {2*sin(\x)});
\node[label, engred] at (1.6, 1.95) {$g = 0$};

% Optimum of f = x+y on the circle: x = y = sqrt(2), f = 2 sqrt(2) ~ 2.828
\coordinate (Pstar) at (1.4142, 1.4142);
\filldraw[engblack] (Pstar) circle (0.06);
\node[label, anchor=south west] at (Pstar) {$\vect{x}^{*}$};

% Gradient of f at the optimum: (1, 1), drawn scaled
\draw[gradf] (Pstar) -- ++(0.8, 0.8) node[anchor=south, label] {$\grad f$};
% Gradient of g at the optimum: (2x, 2y) = (2.83, 2.83), parallel to (1,1)
\draw[gradg] (Pstar) -- ++(1.1, 1.1) node[anchor=south west, label, xshift=4pt] {$\grad g$};

% A "wrong" point on the constraint where the gradients are not parallel
\coordinate (Q) at (1.732, 1.0);  % angle 30 degrees on the circle
\filldraw[engblack!50] (Q) circle (0.04);
\draw[gradf, opacity=0.6] (Q) -- ++(0.6, 0.6);
\draw[gradg, opacity=0.6] (Q) -- ++({1.732*0.4}, {1.0*0.4});
\node[label, anchor=west] at (Q) {\small not optimal};

\end{tikzpicture}
\caption[Lagrange-multiplier tangency.]{The constrained optimum of
$f$ on the curve $g = 0$ is the point where a level curve of $f$ is
tangent to the constraint (heavy black dot, with parallel gradients
$\grad f$ in blue and $\grad g$ in red). At any other point on the
constraint (grey, lower right), the two gradients are not parallel:
moving along the constraint in the direction that decreases the
angle between them changes $f$. The Lagrange-multiplier condition
$\grad f = \lambda \grad g$ encodes the parallelism algebraically,
with $\lambda$ the proportionality factor.}
\label{fig:vol02:ch11:lagrange-tangency}
\end{figure}


\begin{mastery}
The phrase "the constraint defines a surface" rests on the
\engterm{implicit function theorem}, the result that justifies
treating a level set $g(\vect{x}) = 0$ as a smooth manifold. The
theorem states that if $g : \R^{n} \to \R$ is continuously
differentiable near $\vect{x}^{*}$, if $g(\vect{x}^{*}) = 0$, and if
some partial derivative is nonzero there, say $\partial g/\partial
x_{n}(\vect{x}^{*}) \neq 0$, then near $\vect{x}^{*}$ the equation
$g = 0$ can be solved for that variable as a smooth function of the
others, $x_{n} = h(x_{1}, \ldots, x_{n - 1})$, with the solved form's
derivatives given by implicit differentiation,
\[
\pdv{h}{x_{i}} = -\frac{\partial g/\partial x_{i}}
{\partial g/\partial x_{n}}.
\]
\Cref{fig:vol02:ch11:implicit-function} shows the geometry on the
unit circle $g = x^{2} + y^{2} - 1 = 0$: at a point where
$\partial g/\partial y \neq 0$, the curve is locally the graph
$y = h(x)$; at the point $(1, 0)$, where $\partial g/\partial y = 0$
and the tangent is vertical, the local-graph description in $y$
fails and the roles of $x$ and $y$ swap. The nonzero-gradient
condition $\grad g(\vect{x}^{*}) \neq \vect{0}$ is exactly the
constraint qualification that the Lagrange derivation (and the
derivation exercise of this chapter) assumes: it guarantees the
constraint set is a genuine $(n - 1)$-dimensional surface with a
well-defined tangent space, rather than a point, a crossing, or a
cusp where the parallel-gradient argument has no surface to be
tangent to. The vector form, with $g : \R^{n} \to \R^{k}$ a system
of $k$ constraints, requires the constraint Jacobian to have full
row rank $k$ at $\vect{x}^{*}$, which is the multi-constraint
qualification of section 11.4.3.
\end{mastery}

% Figure: the implicit function theorem. A level curve g(x,y)=0 with a
% point where g_y != 0, so the curve is locally the graph of y = h(x);
% and a point where g_y = 0 (vertical tangent), where the local-graph
% description in y fails and x = k(y) must be used instead.
\begin{figure}[htbp]
\centering
\begin{tikzpicture}[scale=2.0,
  curve/.style={engblue, very thick},
  tang/.style={engred, thick},
  axis/.style={engblack, thick, ->},
  label/.style={font=\small},
]
% Axes.
\draw[axis] (-1.55, 0) -- (1.7, 0) node[right, label] {$x$};
\draw[axis] (0, -1.55) -- (0, 1.7) node[above, label] {$y$};

% The level curve g(x,y) = x^2 + y^2 - 1 = 0: the unit circle.
\draw[curve] (0,0) circle (1.0);
\node[label, engblue, anchor=west] at ({cos(40)+0.05}, {sin(40)+0.18})
  {$g(x,y)=0$};

% Regular point P where g_y != 0: locally y = h(x).
\coordinate (P) at ({cos(55)}, {sin(55)});
\filldraw[engblack] (P) circle (0.03);
\node[label, anchor=south west] at (P) {$P$};
% Tangent at P has slope -g_x/g_y = -x/y; draw a short tangent segment.
\draw[tang]
  ($(P) + ({-0.55}, {0.55*cos(55)/sin(55)})$) --
  ($(P) + ({0.55}, {-0.55*cos(55)/sin(55)})$);
\node[label, engred, anchor=west] at ($(P) + (0.45, 0.25)$)
  {$y=h(x)$ here};

% Critical point Q where g_y = 0 (vertical tangent at (1,0)).
\coordinate (Q) at (1.0, 0);
\filldraw[engblack] (Q) circle (0.03);
\node[label, anchor=south west] at ($(Q) + (0.02, 0.02)$) {$Q$};
\draw[tang] ($(Q) + (0, -0.55)$) -- ($(Q) + (0, 0.55)$);
\node[label, engred, anchor=west] at ($(Q) + (0.12, -0.35)$)
  {$g_y=0$: use $x=k(y)$};

\end{tikzpicture}
\caption[The implicit function theorem on a level curve.]{The level
set $g(x, y) = 0$ (here the unit circle) is locally the graph of a
function. At a point $P$ where $\partial g/\partial y \neq 0$, the
curve is locally $y = h(x)$ with $h'(x) = -g_{x}/g_{y}$. At the point
$Q = (1, 0)$ the partial $\partial g/\partial y$ vanishes (the tangent
is vertical) and no single-valued $y = h(x)$ describes the curve there;
the roles of $x$ and $y$ swap, and the curve is locally $x = k(y)$
because $\partial g/\partial x \neq 0$ at $Q$. The implicit function
theorem guarantees a local solved form wherever the relevant partial
is nonzero, which is the same nondegeneracy the Lagrange constraint
qualification requires.}
\label{fig:vol02:ch11:implicit-function}
\end{figure}


\subsection{The Lagrangian}

\begin{definition}[Lagrangian]
For an optimisation problem
\[
\text{minimise (or maximise)}\quad f(\vect{x})\quad
\text{subject to}\quad g(\vect{x}) = 0,
\]
the \engterm{Lagrangian} is
\[
\mathcal{L}(\vect{x}, \lambda) = f(\vect{x}) - \lambda\,g(\vect{x}).
\]
\end{definition}

The first-order conditions for a constrained extremum are obtained
by setting all partial derivatives of $\mathcal{L}$ to zero:
\begin{align*}
\grad_{\vect{x}} \mathcal{L} &= \grad f(\vect{x}) - \lambda\,
\grad g(\vect{x}) = \vect{0}, \\
\pdv{\mathcal{L}}{\lambda} &= -g(\vect{x}) = 0.
\end{align*}
The first equation is the parallel-gradient condition; the second is
the constraint itself. The system is $n + 1$ equations in $n + 1$
unknowns ($x_{1}, \ldots, x_{n}, \lambda$). The reader solves it by
elimination, by substitution, or by the structure of the specific
problem.

\subsection{Worked example: both extrema on a closed curve}

A single equality constraint over a closed curve usually yields two
optima, a constrained maximum and a constrained minimum, and the
Lagrange condition finds both at once. Extremise $f(x, y) = x + 2y$
on the ellipse $g(x, y) = x^{2}/4 + y^{2} - 1 = 0$. The gradients are
$\grad f = (1, 2)^{\mathsf{T}}$ and $\grad g = (x/2, 2y)^{\mathsf{T}}$,
so $\grad f = \lambda\grad g$ reads
\[
1 = \lambda\,\frac{x}{2}, \qquad 2 = \lambda\,(2y),
\]
giving $x = 2/\lambda$ and $y = 1/\lambda$. Substituting into the
constraint,
\[
\frac{(2/\lambda)^{2}}{4} + \Big(\frac{1}{\lambda}\Big)^{2} = 1
\quad\Longrightarrow\quad \frac{1}{\lambda^{2}} + \frac{1}{\lambda^{2}}
= 1 \quad\Longrightarrow\quad \lambda^{2} = 2, \quad \lambda =
\pm\sqrt{2}.
\]
The two multipliers give the two optima. With $\lambda = +\sqrt{2}$,
$(x, y) = (\sqrt{2}, 1/\sqrt{2})$ and $f = \sqrt{2} + 2/\sqrt{2} =
2\sqrt{2} \approx 2.83$, the constrained maximum. With $\lambda =
-\sqrt{2}$, the antipodal point $(-\sqrt{2}, -1/\sqrt{2})$ gives
$f = -2\sqrt{2}$, the constrained minimum. The sign of the multiplier
distinguishes the two: at the maximum $\grad f$ and $\grad g$ point
the same way, at the minimum they oppose. The reader confirms the
answer against the parametrisation $x = 2\cos t$, $y = \sin t$, which
turns $f$ into $2\cos t + 2\sin t = 2\sqrt{2}\sin(t + 45^{\circ})$,
extremised at $\pm 2\sqrt{2}$, agreeing.
\Cref{fig:vol02:ch11:ellipse-extrema} shows the two tangencies. The
lesson carries to any objective on a bounded constraint set: solving
$\grad f = \lambda\grad g$ enumerates every candidate, and the
objective value, not the method, sorts maxima from minima. On an
unbounded constraint curve one or both extrema can fail to exist, and
the multiplier equation then has no real solution in the runaway
direction, the algebraic signature of an objective with no bound on
that constraint.

% Figure: constrained maximum and minimum of a linear objective on an
% ellipse. Maximise and minimise f(x,y) = x + 2y on the ellipse
% x^2/4 + y^2 = 1. The two optima are the points where a level line of f
% (slope -1/2) is tangent to the ellipse; the far tangent is the maximum,
% the near tangent the minimum, and the two constraint gradients point in
% opposite senses, so the multiplier changes sign between them.
\begin{figure}[htbp]
\centering
\begin{tikzpicture}[scale=1.0,
  ell/.style={engblue, very thick},
  lev/.style={enggray, thin},
  tang/.style={engred, thick},
  gvec/.style={engorange, very thick, -{Stealth[length=2.6mm]}},
  pt/.style={circle, fill=engblack, inner sep=1.5pt},
]
  \begin{axis}[
    width=9.6cm, height=6.8cm,
    axis equal image,
    xlabel={$x$}, ylabel={$y$},
    xmin=-3.4, xmax=3.6, ymin=-1.9, ymax=2.0,
    axis lines=middle,
    xtick={-3,-2,-1,1,2,3}, ytick={-1,1},
    tick label style={font=\small},
    label style={font=\small},
    clip=false,
  ]
    % The ellipse x^2/4 + y^2 = 1.
    \addplot[ell, domain=0:360, samples=100] ({2*cos(x)}, {sin(x)});
    % A few level lines f = x + 2y = c, slope -1/2. Non-tangent guides.
    \foreach \c in {-2,0} {
      \addplot[lev, domain=-3.4:3.6, samples=2] {(\c - x)/2};
    }
    % Optimum: grad f = (1,2). f max = sqrt(1+16)= sqrt(17) approx 4.123 at
    % (x,y) = (2*1/sqrt(17*... )). Solve: on ellipse, extremum of x+2y.
    % param x=2cos t, y=sin t; f = 2cos t + 2 sin t = 2 sqrt2 sin(t+45).
    % max at t=45deg: x=2/sqrt2=1.414, y=1/sqrt2=0.707, f=2 sqrt2=2.828.
    % min at t=225deg: x=-1.414, y=-0.707, f=-2.828.
    % Tangent lines through those points, slope -1/2.
    \addplot[tang, domain=-1.2:3.6, samples=2] {(2.828 - x)/2};
    \addplot[tang, domain=-3.4:1.2, samples=2] {(-2.828 - x)/2};
    \node[pt] at (axis cs:1.414,0.707) {};
    \node[pt] at (axis cs:-1.414,-0.707) {};
    \node[engred, anchor=south west, font=\small] at (axis cs:1.5,0.78)
      {max};
    \node[engred, anchor=north east, font=\small] at (axis cs:-1.5,-0.78)
      {min};
    % Constraint gradient (outward normal) at the max point ~ direction (x/2, 2y).
    \draw[gvec] (axis cs:1.414,0.707) -- (axis cs:2.05,2.12);
    \node[engorange, anchor=west, font=\small] at (axis cs:2.0,1.6)
      {$\nabla g$};
  \end{axis}
\end{tikzpicture}
\caption[Constrained maximum and minimum of a linear objective on an
ellipse.]{The objective $f(x, y) = x + 2y$ has straight level lines of slope
$-\tfrac{1}{2}$ (grey). On the ellipse $x^{2}/4 + y^{2} = 1$ (blue) the
objective is extremised where a level line runs tangent to the curve (red).
The far tangent gives the constrained maximum $f = 2\sqrt{2}$ at
$(\sqrt{2}, \tfrac{1}{\sqrt{2}})$; the near tangent gives the constrained
minimum $f = -2\sqrt{2}$ at the antipodal point. At both, the objective
gradient $(1, 2)$ is parallel to the constraint gradient (orange), which is
the Lagrange condition; the multiplier is positive at the maximum and
negative at the minimum, because the two constraint normals point in
opposite senses relative to $\grad f$. One constraint yields two optima, the
generic feature of extremising over a closed curve.}
\label{fig:vol02:ch11:ellipse-extrema}
\end{figure}


\subsection{Multiple equality constraints}

For $k$ equality constraints $g_{1}, \ldots, g_{k}$, the Lagrangian
is
\[
\mathcal{L}(\vect{x}, \boldsymbol{\lambda}) = f(\vect{x}) -
\sum_{j = 1}^{k} \lambda_{j}\,g_{j}(\vect{x}),
\]
and the first-order conditions are
\[
\grad f(\vect{x}) = \sum_{j = 1}^{k} \lambda_{j}\,\grad g_{j}(
\vect{x}), \qquad g_{j}(\vect{x}) = 0 \text{ for } j = 1, \ldots, k.
\]
Each constraint contributes its own multiplier; the optimisation
gradient is a linear combination of the constraint gradients. The
geometric reading: at a constrained extremum, the cost gradient
$\grad f$ lies in the linear span of the constraint gradients
$\grad g_{j}$. The system is now $n + k$ equations in $n + k$
unknowns.

\paragraph{Worked example: two equality constraints in three
variables.} Minimise $f(x, y, z) = x^{2} + y^{2} + z^{2}$ (squared
distance from the origin) subject to the plane $g_{1} = x + y + z -
1 = 0$ and the plane $g_{2} = x - y - z = 0$. The feasible set is
the line where the two planes meet; the problem asks for the point
of that line nearest the origin. The Lagrange condition is $\grad f
= \lambda_{1}\grad g_{1} + \lambda_{2}\grad g_{2}$, that is,
\[
(2x, 2y, 2z) = \lambda_{1}(1, 1, 1) + \lambda_{2}(1, -1, -1).
\]
Componentwise, $2x = \lambda_{1} + \lambda_{2}$, $2y = \lambda_{1} -
\lambda_{2}$, $2z = \lambda_{1} - \lambda_{2}$, so $y = z$.
Substitute into the constraints. The second, $x - y - z = 0$, gives
$x = 2y$. The first, $x + y + z = 1$, gives $2y + y + y = 1$, so
$y = \tfrac{1}{4}$, $z = \tfrac{1}{4}$, $x = \tfrac{1}{2}$. The
nearest point is $(\tfrac{1}{2}, \tfrac{1}{4}, \tfrac{1}{4})$, at
distance $\sqrt{\tfrac{1}{4} + \tfrac{1}{16} + \tfrac{1}{16}} =
\sqrt{6}/4 \approx 0.61$ from the origin. The feasibility audit
passes: $\tfrac{1}{2} + \tfrac{1}{4} + \tfrac{1}{4} = 1$ and
$\tfrac{1}{2} - \tfrac{1}{4} - \tfrac{1}{4} = 0$. The multipliers
follow from $2z = \lambda_{1} - \lambda_{2} = \tfrac{1}{2}$ and
$2x = \lambda_{1} + \lambda_{2} = 1$, giving $\lambda_{1} =
\tfrac{3}{4}$, $\lambda_{2} = \tfrac{1}{4}$. The discipline the
example teaches: every Lagrange solution must be substituted back
into every constraint, because a failed check means an arithmetic
slip, not a failure of the method. The structural point is geometric:
two constraint gradients span a plane, and at the optimum $\grad f$
lies in that plane, so the line from the origin to the nearest point
is perpendicular to the feasible line.

\subsection{Worked example: the optimal can}

The classic worked example of Lagrange. A closed cylindrical can of
radius $r$ and height $h$ has volume $V = \pi r^{2} h$ and surface
area $S = 2\pi r^{2} + 2\pi r h$. Minimise the surface area at fixed
volume $V_{0}$.

The Lagrangian is $\mathcal{L} = 2\pi r^{2} + 2\pi r h -
\lambda(\pi r^{2} h - V_{0})$. The first-order conditions:
\begin{align*}
\pdv{\mathcal{L}}{r} &= 4\pi r + 2\pi h - 2\pi\lambda r h = 0, \\
\pdv{\mathcal{L}}{h} &= 2\pi r - \pi\lambda r^{2} = 0, \\
\pdv{\mathcal{L}}{\lambda} &= -(\pi r^{2} h - V_{0}) = 0.
\end{align*}
From the second equation, $\lambda = 2/r$. Substituting into the
first: $4\pi r + 2\pi h - 4\pi h = 0$, so $h = 2 r$. The optimal
can is as tall as it is wide: a cube of revolution. The constraint
then sets $r = (V_{0}/(2\pi))^{1/3}$, $h = 2 r$. The reader recognises
this as the rule of thumb for a soup can. Real soup cans are not
quite this proportion because the seamed metal cylinder has different
material costs for the side and the lid, which is the worked example
of an exercise.

\subsection{Worked example: maximum entropy on a probability
simplex}

A second worked Lagrange problem, drawn from information theory and
recurring throughout Volumes VII and IX. Among all probability
distributions $(p_{1}, p_{2}, p_{3})$ on three outcomes, which one
maximises the Shannon entropy $H = -\sum_{i} p_{i}\ln p_{i}$? The
single constraint is that the probabilities sum to one,
$g(\vect{p}) = p_{1} + p_{2} + p_{3} - 1 = 0$. The Lagrangian is
\[
\mathcal{L} = -\sum_{i} p_{i}\ln p_{i} - \lambda\Big(\sum_{i} p_{i}
- 1\Big).
\]
The stationarity condition for each $p_{i}$,
\[
\pdv{\mathcal{L}}{p_{i}} = -\ln p_{i} - 1 - \lambda = 0
\quad\Longrightarrow\quad p_{i} = e^{-1 - \lambda},
\]
gives the same value for every $i$: the optimal distribution is
uniform. The constraint then forces $3\,e^{-1 - \lambda} = 1$, so
$p_{i} = 1/3$ for all $i$. The maximum-entropy distribution on a
finite outcome set with no further constraint is the uniform
distribution, the formal statement of the principle of indifference.
The multiplier $\lambda$ is fixed by the normalisation; its value
$\lambda = -1 - \ln 3$ is the log-partition constant. When a further
constraint is added (a fixed mean, say), the same method delivers
the exponential-family distributions that recur in statistical
mechanics (the Boltzmann distribution, Volume IV) and in machine
learning (the softmax and logistic models, Volume VII). The
structural lesson: a Lagrange multiplier per constraint turns a
constrained entropy maximisation into an unconstrained
stationarity, and the answer is an exponential in the constraint
functions.

\subsection{Worked example: closest point on a plane}

A frequent geometric subproblem is to find the point of a plane
nearest a given point, the foot of the perpendicular. Minimise the
squared distance $f(\vect{x}) = \norm{\vect{x} - \vect{x}_{0}}^{2}$
subject to the plane constraint $g(\vect{x}) = \vect{a} \cdot
\vect{x} - b = 0$. The Lagrange condition $\grad f = \lambda \grad
g$ reads $2(\vect{x} - \vect{x}_{0}) = \lambda\,\vect{a}$, so the
displacement $\vect{x} - \vect{x}_{0}$ is parallel to the plane's
normal $\vect{a}$. Write $\vect{x} = \vect{x}_{0} + t\,\vect{a}$ for
a scalar $t = \lambda/2$ and substitute into the constraint:
$\vect{a} \cdot (\vect{x}_{0} + t\,\vect{a}) = b$, giving
\[
t = \frac{b - \vect{a} \cdot \vect{x}_{0}}{\norm{\vect{a}}^{2}},
\qquad
\vect{x}^{*} = \vect{x}_{0} + \frac{b - \vect{a} \cdot \vect{x}_{0}}
{\norm{\vect{a}}^{2}}\,\vect{a}.
\]
The minimum distance is $\norm{\vect{x}^{*} - \vect{x}_{0}} =
\abs{b - \vect{a}\cdot\vect{x}_{0}}/\norm{\vect{a}}$, the standard
point-to-plane distance formula, here derived as a Lagrange problem
rather than memorised. The geometric content is that the shortest
path from a point to a constraint surface runs along the surface
normal, which is exactly the parallel-gradient condition. The same
structure underlies the projection step in constrained optimisation
algorithms (Volume II Chapter 15) and the least-norm solution of an
underdetermined linear system (Volume II Chapter 9).

\subsection{Worked example: equality-constrained least squares}

The least-squares fit of section 11.9 minimises a quadratic with no
constraint. A recurring variant fixes one linear combination of the
parameters exactly while minimising the residual on the rest: a
calibration curve forced through a known reference point, a mixture
whose fractions must sum to one, a control allocation whose total
effort is pinned. The problem is
\[
\text{minimise}\quad \tfrac{1}{2}\norm{X\boldsymbol{\beta} -
\vect{z}}^{2}\quad\text{subject to}\quad
\vect{a}^{\mathsf{T}}\boldsymbol{\beta} = c,
\]
a quadratic objective with one linear equality. The Lagrangian is
\[
\mathcal{L}(\boldsymbol{\beta}, \lambda) =
\tfrac{1}{2}\norm{X\boldsymbol{\beta} - \vect{z}}^{2} -
\lambda(\vect{a}^{\mathsf{T}}\boldsymbol{\beta} - c).
\]
Stationarity in $\boldsymbol{\beta}$ gives $X^{\mathsf{T}}(X
\boldsymbol{\beta} - \vect{z}) - \lambda\vect{a} = \vect{0}$, that
is, the constrained normal equations
\[
X^{\mathsf{T}}X\,\boldsymbol{\beta} = X^{\mathsf{T}}\vect{z} +
\lambda\vect{a}.
\]
Write $\boldsymbol{\beta}_{0} = (X^{\mathsf{T}}X)^{-1}X^{\mathsf{T}}
\vect{z}$ for the unconstrained least-squares solution and $M =
(X^{\mathsf{T}}X)^{-1}$ for the inverse of the normal matrix. Then
$\boldsymbol{\beta} = \boldsymbol{\beta}_{0} + \lambda M\vect{a}$:
the constrained solution is the unconstrained one plus a correction
along the direction $M\vect{a}$, scaled by the multiplier. Impose
the constraint $\vect{a}^{\mathsf{T}}\boldsymbol{\beta} = c$ to fix
$\lambda$,
\[
\vect{a}^{\mathsf{T}}\boldsymbol{\beta}_{0} + \lambda\,
\vect{a}^{\mathsf{T}} M\vect{a} = c
\quad\Longrightarrow\quad
\lambda = \frac{c - \vect{a}^{\mathsf{T}}\boldsymbol{\beta}_{0}}
{\vect{a}^{\mathsf{T}} M\vect{a}}.
\]
The reading is precise. When the unconstrained fit already satisfies
the constraint ($\vect{a}^{\mathsf{T}}\boldsymbol{\beta}_{0} = c$),
the multiplier is zero and no correction is made; the constraint was
slack. Otherwise the correction is proportional to the constraint
residual $c - \vect{a}^{\mathsf{T}}\boldsymbol{\beta}_{0}$ and runs
along $M\vect{a}$, the direction in which moving $\boldsymbol{\beta}$
changes the constraint value most cheaply in the metric set by the
data. A concrete instance: fit a line $z = \beta_{1} x + \beta_{2}$
to data but force it through the origin, $\beta_{2} = 0$, so $\vect{a}
= (0, 1)^{\mathsf{T}}$ and $c = 0$. The formula sets $\beta_{2}$ to
zero and re-solves for $\beta_{1}$ as the through-origin slope
$\sum x_{i} z_{i}/\sum x_{i}^{2}$, the familiar zero-intercept
regression, here derived as the one-constraint Lagrange correction
to the two-parameter fit. The structural lesson generalises to
$k$ constraints: stack them as $A\boldsymbol{\beta} = \vect{c}$ and
the correction becomes $\boldsymbol{\beta} = \boldsymbol{\beta}_{0}
+ M A^{\mathsf{T}}(A M A^{\mathsf{T}})^{-1}(\vect{c} -
A\boldsymbol{\beta}_{0})$, the projection of the unconstrained
solution onto the feasible affine subspace in the $X^{\mathsf{T}}X$
metric.

\subsection{Interpretation of the multiplier}

The Lagrange multiplier $\lambda$ has a precise economic
interpretation. Let $f^{*}(c)$ be the optimal value of $f$ when the
constraint is $g(\vect{x}) = c$ rather than $g(\vect{x}) = 0$. Then
\[
\dv{f^{*}}{c} = \lambda^{*}.
\]
The multiplier reports the rate at which the optimal value changes
when the constraint is relaxed. This is the \engterm{shadow price}
of the constraint: the marginal benefit of one more unit of the
resource the constraint represents. In the optimal-can example, the
multiplier reports the rate at which the minimum surface area
changes per unit increase in the required volume; in a production-
optimisation problem, it reports the rate at which the maximum profit
changes per unit increase in the budget. Volume II Chapter 15
develops shadow prices as the dual-problem analysis of constrained
optimisation.

\paragraph{Worked example: the shadow price computed two ways.}
Maximise $f(x, y) = xy$ (a rectangular area) subject to the perimeter
budget $g(x, y) = 2x + 2y - P = 0$. The Lagrangian condition $\grad
f = \lambda\grad g$ reads $(y, x) = \lambda(2, 2)$, so $x = y$, and
the constraint $4x = P$ gives $x = y = P/4$, the square. The optimal
area is $f^{*} = (P/4)^{2} = P^{2}/16$. The multiplier follows from
$y = 2\lambda$, so $\lambda = y/2 = P/8$.

Now verify the shadow-price identity $\dv{f^{*}}{P} = \lambda^{*}$
by differentiating the optimal value directly. With $f^{*}(P) =
P^{2}/16$,
\[
\dv{f^{*}}{P} = \frac{2P}{16} = \frac{P}{8} = \lambda^{*}.
\]
The two routes agree exactly. The engineering reading: at the
optimal square, one more unit of perimeter budget buys $P/8$ more
units of area, and that marginal rate is precisely the Lagrange
multiplier. A facilities engineer allocating fence between two
adjacent enclosures equalises the shadow prices across enclosures;
the enclosure with the higher $\lambda$ gets the next unit of
perimeter, until the marginal areas match. The multiplier is not an
algebraic by-product but the decision variable for resource
allocation, which is why Volume II Chapter 15 builds the entire dual
theory on it.

% Figure: the value function f*(c) versus the constraint level c, with the
% Lagrange multiplier as its slope. Illustrated on the area-maximisation
% f* = P^2/16 whose slope is P/8 = lambda*. The tangent line at one budget
% shows lambda as the local marginal value of relaxing the constraint.
\begin{figure}[htbp]
\centering
\begin{tikzpicture}[scale=1.0,
  curve/.style={engblue, very thick},
  tang/.style={engred, thick},
  guide/.style={enggray, dashed},
]
  \begin{axis}[
    width=9.2cm, height=6.4cm,
    xlabel={constraint level $c$ (perimeter $P$)},
    ylabel={optimal value $f^{*}(c)$},
    xmin=0, xmax=16, ymin=0, ymax=17,
    axis lines=left,
    xtick={0,4,8,12,16}, ytick={0,4,8,12,16},
    tick label style={font=\small},
    label style={font=\small},
    clip=false,
  ]
    % Value function f* = c^2/16 (with c = P, f* = P^2/16).
    \addplot[curve, domain=0:16, samples=60] {x*x/16};
    \node[engblue, anchor=south east] at (axis cs:15.5, 15.0) {$f^{*} = c^{2}/16$};
    % Tangent at c = 8: f*(8) = 4, slope = 8/8 = 1. Line y = 4 + 1*(x-8).
    \addplot[tang, domain=2:14, samples=2] {4 + 1*(x-8)};
    \node[engred, anchor=north west] at (axis cs:12.2, 8.0)
      {slope $=\lambda^{*}=\dfrac{c}{8}$};
    % Marker at the tangent point.
    \draw[guide] (axis cs:8,0) -- (axis cs:8,4);
    \draw[guide] (axis cs:0,4) -- (axis cs:8,4);
    \filldraw[engred] (axis cs:8,4) circle (2.4pt);
  \end{axis}
\end{tikzpicture}
\caption[The value function and the multiplier as its slope.]{The value
function $f^{*}(c)$, the best achievable objective when the constraint is
$g = c$, for the area-maximisation whose optimum is $f^{*} = c^{2}/16$
(a rectangle of fixed perimeter $c = P$). The Lagrange multiplier is the
slope of this curve, $\lambda^{*} = \dv{f^{*}}{c} = c/8$, drawn as the
red tangent at $c = 8$. The multiplier reads off the marginal value of
one more unit of the constrained resource: at this operating point one
extra unit of perimeter buys $\lambda^{*} = 1$ unit of area. The rising
slope shows that the shadow price itself grows with the budget for this
problem, so the marginal value of relaxation is not constant.}
\label{fig:vol02:ch11:shadow-price}
\end{figure}


\Cref{fig:vol02:ch11:shadow-price} draws the value function
$f^{*}(c) = c^{2}/16$ of the area problem and its tangent at one
operating point. The multiplier is the tangent's slope, so a picture of
the value function is a picture of the shadow price at every budget: a
steeper part of the curve is a more valuable constraint to relax. The
convexity of the curve (the shadow price rises with the budget for this
problem) is the geometric statement that the marginal value of a fixed
increment of resource is not constant, a fact a facilities engineer uses
when deciding whether to buy the next unit of perimeter now or later.

\begin{mastery}
The shadow-price identity $\dv{f^{*}}{c} = \lambda^{*}$ is one case
of the \engterm{envelope theorem}, which reports how a constrained
optimum responds to any parameter, not only a constraint level. Let
the problem carry a parameter $\alpha$ throughout, in the objective,
the constraint, or both: minimise $f(\vect{x}, \alpha)$ subject to
$g(\vect{x}, \alpha) = 0$, with optimum $\vect{x}^{*}(\alpha)$ and
optimal value $f^{*}(\alpha) = f(\vect{x}^{*}(\alpha), \alpha)$. The
naive derivative of $f^{*}$ would need $\dd\vect{x}^{*}/\dd\alpha$,
the sensitivity of the optimiser itself, which is expensive to
compute. The envelope theorem says that term drops out. Differentiate
the Lagrangian value $\mathcal{L}(\vect{x}^{*}, \lambda^{*}, \alpha)
= f^{*}(\alpha)$ (the two agree at the optimum because $g = 0$ there),
\[
\dv{f^{*}}{\alpha} = \pdv{\mathcal{L}}{\vect{x}}\cdot
\dv{\vect{x}^{*}}{\alpha} + \pdv{\mathcal{L}}{\lambda}\,
\dv{\lambda^{*}}{\alpha} + \pdv{\mathcal{L}}{\alpha}
= \pdv{\mathcal{L}}{\alpha}\bigg|_{\vect{x}^{*}, \lambda^{*}},
\]
because the first two coefficients are the first-order conditions,
$\partial\mathcal{L}/\partial\vect{x} = \vect{0}$ and
$\partial\mathcal{L}/\partial\lambda = -g = 0$, and vanish. Only the
explicit dependence survives: the optimal value's sensitivity to a
parameter is the Lagrangian's partial with respect to that parameter,
evaluated at the fixed optimum, with the optimiser's own movement
ignored. The shadow price is the special case $\alpha = c$, where
$\partial\mathcal{L}/\partial c = \lambda^{*}$. The theorem is the
reason a sensitivity analysis of an optimised design is cheap: to
first order, an optimum does not care that it would re-optimise, so
the marginal effect of a small parameter change is read at the frozen
solution. Volume II Chapter 15 builds duality on the same identity,
and the physicist's Hellmann-Feynman theorem, the economist's Roy and
Shephard identities, and the control engineer's costate equation are
all the envelope theorem wearing different notation.
\end{mastery}

\subsection{Second-order sufficient conditions and the bordered
Hessian}

The Lagrangian first-order conditions are necessary for a constrained
extremum but, as in the unconstrained case, not sufficient: the
candidate $(\vect{x}^{*}, \boldsymbol{\lambda}^{*})$ could be a
saddle of $\mathcal{L}$ rather than a constrained minimum or maximum
of $f$. The sufficient condition restricts the Hessian of
$\mathcal{L}$ to the tangent space of the active constraints.

For one equality constraint $g(\vect{x}) = 0$, the tangent space at
$\vect{x}^{*}$ is $T = \set{\vect{d} \in \R^{n} : \grad
g(\vect{x}^{*}) \cdot \vect{d} = 0}$. The Hessian of the Lagrangian
with respect to $\vect{x}$ is
\[
H_{\mathcal{L}} =
H_{f}(\vect{x}^{*}) - \lambda^{*}\,H_{g}(\vect{x}^{*}).
\]
The second-order sufficient condition for a constrained local
minimum is
\[
\vect{d}^{\mathsf{T}}\,H_{\mathcal{L}}\,\vect{d} > 0 \quad
\text{for all nonzero } \vect{d} \in T,
\]
the positive-definiteness of $H_{\mathcal{L}}$ on the tangent space.
The constrained maximum condition is the corresponding
negative-definiteness on $T$.

The \engterm{bordered Hessian} packages this test in a single
matrix. For $k$ equality constraints with Jacobian $G \in
\R^{k \times n}$ whose rows are $\grad g_{j}^{\mathsf{T}}$, the
bordered Hessian is the $(n + k) \times (n + k)$ symmetric matrix
\[
\bar{H} = \begin{pmatrix} 0 & G \\ G^{\mathsf{T}} & H_{\mathcal{L}}
\end{pmatrix}.
\]
The inertia of $\bar{H}$ (its count of positive, negative, and zero
eigenvalues) relates to the inertia of $H_{\mathcal{L}}$ restricted
to $T$ by a theorem in constrained-optimisation theory
\cite{text:v2ch11-nocedal-wright-2006}. For two variables and one
constraint, the working test is that the bordered Hessian's $3
\times 3$ determinant has a sign that identifies the constrained
extremum's type; a positive determinant signals a constrained
maximum and a negative determinant a constrained minimum, with the
caveat that the active-constraint structure has been correctly
identified.

\paragraph{Worked example: the bordered-Hessian test in numbers.}
Take the area-maximisation $f(x, y) = xy$ subject to $g = x + y - 10
= 0$. The Lagrange solution is $x = y = 5$, $\lambda = 5$. The
Lagrangian Hessian in $(x, y)$ is $H_{\mathcal{L}} = H_{f} - \lambda
H_{g}$; here $H_{f} = \begin{pmatrix} 0 & 1 \\ 1 & 0 \end{pmatrix}$
and $H_{g} = \vect{0}$ (the constraint is linear), so
$H_{\mathcal{L}} = \begin{pmatrix} 0 & 1 \\ 1 & 0 \end{pmatrix}$.
The constraint gradient is $\grad g = (1, 1)$, so the bordered
Hessian is
\[
\bar{H} = \begin{pmatrix} 0 & 1 & 1 \\ 1 & 0 & 1 \\ 1 & 1 & 0
\end{pmatrix}.
\]
Its determinant is $\det\bar{H} = 0\cdot(0 - 1) - 1\cdot(0 - 1) +
1\cdot(1 - 0) = 0 + 1 + 1 = 2 > 0$. By the $3 \times 3$
bordered-Hessian rule for one constraint in two variables, a
positive determinant signals a constrained maximum, which matches the
geometry: among rectangles of fixed perimeter $20$, the square of
side $5$ has the maximum area. The reader confirms the type without
appealing to the physical answer, which is the point of the
second-order test: it certifies the candidate algebraically.

For the optimal-can example, the bordered-Hessian computation
confirms that $h = 2 r$ is a constrained minimum of the surface
area, not a constrained saddle. The verification is one exercise
the reader is encouraged to attempt; the working takeaway is that
the Lagrange first-order conditions identify candidates, and a
bordered-Hessian or equivalent second-order test distinguishes
minima from maxima from saddles. In well-posed engineering problems
the physical sign of the answer (an optimal can has a finite
positive radius) provides an auxiliary sanity check; the bordered
test is the algebraic confirmation.

\begin{mastery}
The multiplier method is named for Joseph-Louis Lagrange, who
introduced it in the \textit{M\'ecanique analytique} (1788) as a way
to handle constraints in mechanics without solving for the
constraint forces explicitly. The constraint forces are exactly the
multipliers: a bead on a wire feels a normal force from the wire,
and that force is $\lambda\grad g$, the multiplier times the
constraint gradient. The same algebra serves an economist pricing a
scarce resource, a statistician maximising entropy, and a
machine-learning engineer enforcing a margin. The inequality
extension waited until the twentieth century: William Karush stated
the conditions in a 1939 master's thesis, and Harold Kuhn and Albert
Tucker rediscovered them in 1951, which is why the conditions carry
all three names \cite{text:v2ch11-nocedal-wright-2006}. The
historical arc, from constraint forces in mechanics to dual prices
in optimisation, is one of the clean examples of a single piece of
mathematics serving four centuries of engineering.
\end{mastery}

\begin{estimation}
\textbf{Question.} A factory produces two products. Product A
yields a profit of $\$40$ per unit and consumes $2\,\si{\hour}$ of
machine time and $3\,\si{\kilogram}$ of raw material per unit.
Product B yields $\$30$ per unit and consumes $1\,\si{\hour}$ and
$2\,\si{\kilogram}$ per unit. Machine time is capped at
$100\,\si{\hour}$ per week; raw material is capped at
$180\,\si{\kilogram}$ per week (current as of 2026). Estimate the
shadow price of one extra hour of machine time, then solve exactly
by KKT.

\textbf{Estimate.} The unconstrained profit gradient is $(40, 30)$
per unit. A binding machine-time constraint at $100\,\si{\hour}$
yields, in the all-A corner, $50$ units of A for a profit of
$\$2000$; in the all-B corner, $100$ units of B for a profit of
$\$3000$. The all-B corner uses $200\,\si{\kilogram}$ of raw
material, which exceeds the $180\,\si{\kilogram}$ cap by
$20\,\si{\kilogram}$; the raw-material constraint is therefore
active too. The reader expects both constraints binding, with the
optimum somewhere between the two corners. A guess of $\$3300$
total profit at roughly $(x_{A}, x_{B}) = (20, 60)$ gives a shadow
price on machine time of about $\$10$--$15$ per hour.

\textbf{Calculate.} The linear program is
\[
\text{maximise } 40 x_{A} + 30 x_{B} \text{ subject to }
2 x_{A} + x_{B} \leq 100,\
3 x_{A} + 2 x_{B} \leq 180,\
x_{A}, x_{B} \geq 0.
\]
Both upper-bound constraints binding gives the linear system
$2 x_{A} + x_{B} = 100$ and $3 x_{A} + 2 x_{B} = 180$. Solving:
$x_{A} = 20$, $x_{B} = 60$, with profit $40 \cdot 20 + 30 \cdot 60
= \$2600$. The estimate was high. The KKT stationarity conditions
read $(40, 30) = \mu_{1} (2, 1) + \mu_{2} (3, 2)$, two equations
in two multipliers: $2\mu_{1} + 3\mu_{2} = 40$, $\mu_{1} + 2\mu_{2}
= 30$. Solving: $\mu_{1} = -10$, $\mu_{2} = 20$. The negative
$\mu_{1}$ is the signature that the machine-time constraint is not
in fact active at the optimum (relaxing it would not help). Drop
the machine-time constraint: $3 x_{A} + 2 x_{B} = 180$ with
profit-gradient stationarity along its normal gives, after
re-solving, $(x_{A}, x_{B}) = (0, 90)$ with profit $\$2700$. But
this violates $2 x_{A} + x_{B} \leq 100$ (gives $90 \leq 100$,
satisfied), so $(0, 90)$ is feasible and yields $\$2700$. Try the
other corner with only the machine-time constraint binding:
$2 x_{A} + x_{B} = 100$, $\mu_{2} = 0$, stationarity $(40, 30) =
\mu_{1} (2, 1)$ has no consistent solution unless we hit a corner.
The vertex enumeration of the feasible polygon gives candidates
$(0, 0), (50, 0), (0, 90), (20, 60), (0, 100)\text{ infeasible}$.
Among feasible vertices, profits are $0, \$2000, \$2700, \$2600$.
The optimum is at $(0, 90)$ with profit $\$2700$ and shadow price
$\mu_{2} = 30/2 = \$15$ per kilogram of raw material;
$\mu_{1} = 0$ because the machine-time constraint is inactive at
the optimum (it uses $90 < 100$ hours). The estimate's intuition
that both constraints would bind was wrong, and the shadow-price
guess landed in the right range only by accident. The exact
shadow price of one more kilogram of raw material is $\$15$; the
shadow price of one more hour of machine time is zero.
\end{estimation}

\subsection{Applied vignette: the minimum-variance portfolio}
\pathtag{standard}

A quantity engineer, a reliability planner, and an investment analyst
solve the same problem when they allocate a fixed total across several
options whose outcomes vary and co-vary: choose the mix that meets a
required return at the least variance. The equality-constrained
Lagrange machinery of this section delivers the answer in closed form
and, through the multiplier, prices the return against the risk. We
work the two-constraint case end to end and read the result as the
efficient frontier that the same calculation sweeps out.

\paragraph{The model.} Let $\vect{w} = (w_{1}, \ldots, w_{n})^{
\mathsf{T}}$ be the fractions of a fixed budget placed in $n$ assets,
each asset $i$ with expected return $\mu_{i}$, the returns collected
in $\boldsymbol{\mu}$. The joint variability is the covariance matrix
$\Sigma$, symmetric positive-definite, whose entry $\Sigma_{ij}$ is
the covariance of assets $i$ and $j$. The portfolio's expected return
is $\boldsymbol{\mu}^{\mathsf{T}}\vect{w}$ and its variance is the
quadratic form $\vect{w}^{\mathsf{T}}\Sigma\vect{w}$, the same
covariance-weighted squared length the whitening example of section
11.7.5 diagonalised. The allocation must exhaust the budget, so the
weights sum to one, $\vect{1}^{\mathsf{T}}\vect{w} = 1$, and the
investor fixes a target return $\boldsymbol{\mu}^{\mathsf{T}}\vect{w}
= \mu_{0}$. The problem is
\[
\text{minimise}\quad \tfrac{1}{2}\vect{w}^{\mathsf{T}}\Sigma\vect{w}
\quad\text{subject to}\quad
\vect{1}^{\mathsf{T}}\vect{w} = 1, \quad
\boldsymbol{\mu}^{\mathsf{T}}\vect{w} = \mu_{0},
\]
a convex quadratic with two linear equalities, the canonical
two-constraint Lagrange problem.

\paragraph{Estimate before optimising.} With a single asset the
variance is fixed and there is no choice. With two uncorrelated assets
of equal variance $\sigma^{2}$, splitting the budget evenly gives
portfolio variance $\tfrac{1}{4}\sigma^{2} + \tfrac{1}{4}\sigma^{2} =
\tfrac{1}{2}\sigma^{2}$, half the variance of holding either alone:
diversification halves the risk at $n = 2$ and, for $n$ equal
uncorrelated assets, cuts it to $\sigma^{2}/n$. The reader expects the
optimum to spread weight across assets to exploit this averaging,
tilted toward the low-variance and low-correlation assets, and tilted
again by the return target $\mu_{0}$ away from the pure
minimum-variance mix.

\paragraph{The Lagrange solution.} With multipliers $\lambda$ for the
budget and $\gamma$ for the return, the Lagrangian is
\[
\mathcal{L} = \tfrac{1}{2}\vect{w}^{\mathsf{T}}\Sigma\vect{w}
- \lambda(\vect{1}^{\mathsf{T}}\vect{w} - 1)
- \gamma(\boldsymbol{\mu}^{\mathsf{T}}\vect{w} - \mu_{0}).
\]
Stationarity in $\vect{w}$ gives $\Sigma\vect{w} = \lambda\vect{1} +
\gamma\boldsymbol{\mu}$, so the optimal weights are
\[
\vect{w}^{*} = \Sigma^{-1}(\lambda\vect{1} + \gamma\boldsymbol{\mu})
= \lambda\,\Sigma^{-1}\vect{1} + \gamma\,\Sigma^{-1}\boldsymbol{\mu}.
\]
The solution is a combination of two fixed vectors, $\Sigma^{-1}
\vect{1}$ and $\Sigma^{-1}\boldsymbol{\mu}$, whose coefficients the
two constraints fix. Imposing $\vect{1}^{\mathsf{T}}\vect{w}^{*} = 1$
and $\boldsymbol{\mu}^{\mathsf{T}}\vect{w}^{*} = \mu_{0}$ produces two
linear equations in $\lambda$ and $\gamma$. Writing the three scalars
$A = \vect{1}^{\mathsf{T}}\Sigma^{-1}\vect{1}$, $B = \vect{1}^{
\mathsf{T}}\Sigma^{-1}\boldsymbol{\mu}$, and $C = \boldsymbol{\mu}^{
\mathsf{T}}\Sigma^{-1}\boldsymbol{\mu}$, the system is $\lambda A +
\gamma B = 1$ and $\lambda B + \gamma C = \mu_{0}$, solved by
\[
\lambda = \frac{C - B\mu_{0}}{AC - B^{2}}, \qquad
\gamma = \frac{A\mu_{0} - B}{AC - B^{2}}.
\]
The denominator $AC - B^{2}$ is positive whenever $\boldsymbol{\mu}$
is not a multiple of $\vect{1}$ (the assets have distinct returns), by
the Cauchy-Schwarz inequality in the $\Sigma^{-1}$ metric, so the
solution is well-defined. Substituting $\vect{w}^{*}$ back into the
variance gives the minimum variance at target $\mu_{0}$,
\[
\sigma_{p}^{2}(\mu_{0}) = \vect{w}^{*\mathsf{T}}\Sigma\vect{w}^{*}
= \frac{A\mu_{0}^{2} - 2 B\mu_{0} + C}{AC - B^{2}},
\]
a parabola in the target return.

\paragraph{The reading.} The parabola $\sigma_{p}^{2}(\mu_{0})$ is the
\engterm{efficient frontier} of the portfolio: the least variance
attainable at each target return.
\Cref{fig:vol02:ch11:mean-variance-frontier} plots it in the
variance-return plane. Its vertex, at $\mu_{0} = B/A$, is the
global minimum-variance portfolio $\vect{w} = \Sigma^{-1}\vect{1}/A$,
the mix an investor holds who wants the least risk and accepts
whatever return it carries. Above the vertex sits the efficient
branch, where more return costs more variance; below it, the
dominated branch, whose portfolios are strictly worse than one
directly above at the same variance. The return multiplier $\gamma$
is the shadow price of the return target: by the envelope theorem, it
is the rate $\dv{}{\mu_{0}}(\tfrac{1}{2}\sigma_{p}^{2})$, the marginal
half-variance the investor pays for one more unit of required return,
which is the local steepness of the frontier read from the picture.
An investor who demands a higher return moves up the frontier and
pays for it in variance at the rate $\gamma$; the budget multiplier
$\lambda$ prices the sum-to-one constraint, the cost of the last unit
of capital committed. Casting the allocation as a constrained
quadratic turns a vague preference for return over risk into a
priced trade the two multipliers make explicit, and the same
$\Sigma^{-1}$ that solves the mean-variance problem is the whitening
map of section 11.7.5, so the risk geometry and the change-of-variable
geometry of this chapter are one object seen twice.

% Figure: the mean-variance efficient frontier traced by a Lagrange
% problem. Minimising portfolio variance at each target return, over
% weights that sum to one, sweeps out a parabola in the
% (variance, return) plane; its left edge is the minimum-variance
% portfolio and its upper branch is the efficient frontier. The
% Lagrange multiplier on the return constraint is the slope of the
% frontier in the (variance, return) plane.
\begin{figure}[htbp]
\centering
\begin{tikzpicture}[scale=1.0,
  eff/.style={engblue, very thick},
  ineff/.style={engblue, very thick, dashed},
  guide/.style={enggray, dashed},
  pt/.style={circle, fill=engblack, inner sep=1.5pt},
]
  \begin{axis}[
    width=9.6cm, height=6.6cm,
    xlabel={portfolio variance $\sigma_{p}^{2}$},
    ylabel={expected return $\mu_{p}$},
    xmin=0, xmax=0.09, ymin=0.02, ymax=0.16,
    axis lines=left,
    xtick={0,0.02,0.04,0.06,0.08},
    ytick={0.04,0.08,0.12,0.16},
    tick label style={font=\small},
    label style={font=\small},
    clip=false,
  ]
    % Frontier: variance is quadratic in return, sigma^2 = a(mu - mu0)^2 + s0.
    % Take mu0 = 0.08 (min-variance return), s0 = 0.01 (min variance),
    % a = 4. Upper (efficient) branch mu >= 0.08, lower (inefficient) branch.
    % Efficient branch (solid): mu from 0.08 up to 0.155.
    \addplot[eff, domain=0.08:0.155, samples=60]
      ({4*(x-0.08)^2 + 0.01}, {x});
    % Inefficient branch (dashed): mu from 0.03 to 0.08.
    \addplot[ineff, domain=0.03:0.08, samples=40]
      ({4*(x-0.08)^2 + 0.01}, {x});
    % Minimum-variance portfolio marker.
    \node[pt] at (axis cs:0.01,0.08) {};
    \node[anchor=west, font=\small] at (axis cs:0.012,0.075)
      {min-variance};
    % A target-return operating point on the efficient branch.
    % mu = 0.12: sigma^2 = 4*(0.04)^2 + 0.01 = 4*0.0016+0.01 = 0.0164.
    \node[pt] at (axis cs:0.0164,0.12) {};
    \draw[guide] (axis cs:0.0164,0.02) -- (axis cs:0.0164,0.12);
    \node[engblue, anchor=south west, font=\small] at (axis cs:0.045,0.135)
      {efficient frontier};
    \node[enggray, anchor=north west, font=\small] at (axis cs:0.03,0.06)
      {inefficient};
  \end{axis}
\end{tikzpicture}
\caption[The mean-variance efficient frontier as a swept Lagrange
solution.]{Minimising portfolio variance at each fixed target return, over
asset weights constrained to sum to one, is a Lagrange problem whose
solution, swept across all target returns, traces a parabola in the
variance-return plane. The leftmost point is the minimum-variance
portfolio; the solid upper branch is the efficient frontier, the set of
portfolios with the least variance for their return, and the dashed lower
branch is dominated, since each of its portfolios shares a variance with a
higher-return portfolio directly above it. The Lagrange multiplier on the
return constraint is the marginal variance the investor pays for one more
unit of return, read off as the local steepness of the frontier.}
\label{fig:vol02:ch11:mean-variance-frontier}
\end{figure}


\section{Constrained optimisation preview: KKT conditions}
\pathtag{standard}

Equality constraints are the easy case. Most engineering optimisation
also includes inequality constraints: a budget that must not be
exceeded, a material strength that must not be reduced, a temperature
that must not be exceeded, a probability that must not be lowered.
The full apparatus belongs to Volume II Chapter 15. This section
states the conditions and connects them to Lagrange multipliers, so
the reader can solve simple instances by hand.

\subsection{The general nonlinear program}

Consider
\begin{align*}
\text{minimise}\quad & f(\vect{x}) \\
\text{subject to}\quad & g_{j}(\vect{x}) = 0, \quad j = 1, \ldots, k, \\
& h_{\ell}(\vect{x}) \leq 0, \quad \ell = 1, \ldots, m.
\end{align*}
The feasible set is the intersection of the $k$ equality surfaces
and the $m$ inequality half-spaces. An inequality constraint $h_{\ell}
\leq 0$ is \engterm{active} at a point if $h_{\ell}(\vect{x}) = 0$
there, and \engterm{inactive} if $h_{\ell}(\vect{x}) < 0$.

\subsection{The Karush-Kuhn-Tucker conditions}

\begin{theorem}[KKT conditions, necessary]
If $\vect{x}^{*}$ is a local minimum of the program above and a
constraint-qualification holds at $\vect{x}^{*}$, then there exist
multipliers $\lambda_{j}$ (one per equality constraint) and
$\mu_{\ell} \geq 0$ (one per inequality constraint) such that
\begin{enumerate}[leftmargin=*]
  \item Stationarity: $\grad f(\vect{x}^{*}) = \sum_{j} \lambda_{j}\,
    \grad g_{j}(\vect{x}^{*}) + \sum_{\ell} \mu_{\ell}\,
    \grad h_{\ell}(\vect{x}^{*})$.
  \item Primal feasibility: $g_{j}(\vect{x}^{*}) = 0$ and
    $h_{\ell}(\vect{x}^{*}) \leq 0$ for all $j, \ell$.
  \item Dual feasibility: $\mu_{\ell} \geq 0$ for all $\ell$.
  \item Complementary slackness: $\mu_{\ell} h_{\ell}(\vect{x}^{*}) =
    0$ for all $\ell$.
\end{enumerate}
\end{theorem}

The first three conditions are the obvious generalisation of
Lagrange multipliers: equality constraints carry signed
multipliers; inequality constraints carry non-negative multipliers
because relaxing an inequality (making the right-hand side larger)
can only help, never hurt. The fourth condition, complementary
slackness, says the multiplier is zero when the constraint is
inactive: an inequality that is not binding contributes nothing to
the optimisation.

The constraint-qualification clause (linear independence of the
active constraint gradients, or one of several equivalent
conditions) rules out edge cases where the gradients line up so as
to make the multiplier story degenerate. In well-posed engineering
problems it is satisfied; the reader should know it exists and not
spend time on it on first reading. \textcite{text:boyd-vandenberghe2004}
treats the constraint qualification carefully.

\subsection{Worked example: a quadratic program with one inequality}

Minimise $f(x, y) = x^{2} + y^{2}$ subject to $h(x, y) = x + y - 1
\leq 0$. Let $\mu \geq 0$ be the multiplier of the inequality.
Stationarity: $(2x, 2y) = \mu (1, 1)$, so $x = y = \mu/2$. Primal
feasibility: $x + y \leq 1$. Dual feasibility: $\mu \geq 0$.
Complementary slackness: $\mu (x + y - 1) = 0$.

There are two cases. If the inequality is inactive ($x + y < 1$),
complementary slackness forces $\mu = 0$, stationarity forces
$x = y = 0$, and the inequality is satisfied with strict slack. The
optimum is the origin, $f^{*} = 0$. If the inequality is active
($x + y = 1$), complementary slackness is automatic; stationarity
plus the constraint gives $x = y = 1/2$, $\mu = 1$, and
$f = 1/2 > 0$. The first case wins. The reader should expect the
unconstrained optimum to dominate when the inequality is satisfied
at it.

The same problem with the constraint reversed, $h(x, y) = 1 - x - y
\leq 0$ (i.e.\ $x + y \geq 1$), forces the inequality to be active
at the optimum, gives $x = y = 1/2$, $\mu = 1$, $f^{*} = 1/2$. The
multiplier $\mu = 1$ is the rate at which the optimum increases
per unit increase in the right-hand side $1$ of the constraint.

\subsection{Worked example: two inequalities and the active-set
case sweep}

The systematic method for inequality constraints is the
\engterm{active-set sweep}: enumerate every subset of inequalities
that might be active, solve the equality-only Lagrange system for
each, and keep the candidates that pass primal feasibility (the
inactive inequalities hold) and dual feasibility (the active
multipliers are non-negative). We use the stationarity convention of
section 11.5.2, $\grad f + \sum_{\ell}\mu_{\ell}\grad h_{\ell} =
\vect{0}$ with $\mu_{\ell} \geq 0$, for the minimisation
$f \text{ subject to } h_{\ell} \leq 0$.

Minimise $f(x, y) = (x - 2)^{2} + (y - 2)^{2}$ subject to $h_{1} =
x + y - 1 \leq 0$ and $h_{2} = -x \leq 0$ (that is, $x \geq 0$). The
unconstrained minimum is at $(2, 2)$, which violates $h_{1}$
($4 - 1 = 3 > 0$); the constraints will bind.

With two inequalities there are four cases.

\textit{Case A: neither active.} Stationarity gives $(2(x - 2),
2(y - 2)) = \vect{0}$, so $(2, 2)$. Primal check: $h_{1}(2, 2) =
3 > 0$, infeasible. Discard.

\textit{Case B: only $h_{1}$ active.} Here $\mu_{2} = 0$ and
$\grad f + \mu_{1}\grad h_{1} = \vect{0}$ reads $(2(x - 2), 2(y -
2)) + \mu_{1}(1, 1) = \vect{0}$, with $x + y = 1$. The two
components give $x - 2 = y - 2$, so $x = y$, and $x + y = 1$ gives
$x = y = \tfrac{1}{2}$. Then $2(\tfrac{1}{2} - 2) + \mu_{1} = 0$, so
$\mu_{1} = 3 > 0$. Primal check on the inactive constraint:
$h_{2}(\tfrac{1}{2}, \tfrac{1}{2}) = -\tfrac{1}{2} \leq 0$,
satisfied. Dual feasibility holds. This candidate survives:
$(\tfrac{1}{2}, \tfrac{1}{2})$ with $f^{*} = 2(\tfrac{3}{2})^{2} =
\tfrac{9}{2}$.

\textit{Case C: only $h_{2}$ active.} Here $\mu_{1} = 0$, $x = 0$,
and $\grad f + \mu_{2}\grad h_{2} = \vect{0}$ reads $(2(0 - 2),
2(y - 2)) + \mu_{2}(-1, 0) = \vect{0}$. The second component forces
$y = 2$; the first gives $-4 - \mu_{2} = 0$, so $\mu_{2} = -4 < 0$.
Dual feasibility fails. Discard. The negative multiplier is
informative: it says the assumed active set is wrong because at
$(0, 2)$ the objective gradient points back into the feasible region.

\textit{Case D: both active.} The two equalities $x + y = 1$ and
$x = 0$ give $(0, 1)$. Stationarity $(2(0 - 2), 2(1 - 2)) +
\mu_{1}(1, 1) + \mu_{2}(-1, 0) = \vect{0}$ reads $(-4 + \mu_{1} -
\mu_{2}, -2 + \mu_{1}) = \vect{0}$, so $\mu_{1} = 2$ and $\mu_{2} =
\mu_{1} - 4 = -2 < 0$. Dual feasibility fails. Discard.

Only Case B survives all checks, so the constrained optimum is
$(\tfrac{1}{2}, \tfrac{1}{2})$ with $f^{*} = \tfrac{9}{2}$ and active
multiplier $\mu_{1} = 3$. The geometry confirms it: the
unconstrained minimum $(2, 2)$ lies above the line $x + y = 1$, and
the nearest feasible point is the foot of the perpendicular from
$(2, 2)$ to that line, which is $(\tfrac{1}{2}, \tfrac{1}{2})$. The
sweep over four cases, with primal and dual feasibility checked at
each, is the transferable method; the multiplier sign is the
diagnostic that prunes the wrong active sets.

% Figure: the KKT active-set example. Minimise (x-2)^2 + (y-2)^2 subject to
% x + y <= 1 and x >= 0. The unconstrained minimum (2,2) is infeasible; the
% constrained optimum is the foot of the perpendicular from (2,2) to the line
% x + y = 1, at (1/2, 1/2). Circular contours of f are centred on (2,2).
\begin{figure}[htbp]
\centering
\begin{tikzpicture}[scale=1.6,
  contour/.style={engblue, thick},
  cons/.style={engamber, very thick},
  feas/.style={englime!25},
  perp/.style={engred, thick, dashed},
]
  % Feasible region: x >= 0 and x + y <= 1, clipped to a viewing box.
  \begin{scope}
    \clip (-0.4, -0.9) rectangle (2.6, 2.6);
    \fill[feas] (0, -0.9) -- (0, 1) -- (1.9, -0.9) -- cycle;
  \end{scope}
  % Contours of f = (x-2)^2 + (y-2)^2, circles centred on (2,2).
  \foreach \r in {0.5, 1.0, 1.5, 2.121} {
    \draw[contour] (2, 2) circle (\r);
  }
  % Constraint lines.
  \draw[cons] (-0.4, 1.4) -- (1.9, -0.9) node[anchor=north west, pos=0.15]{$x + y = 1$};
  \draw[cons] (0, -0.9) -- (0, 2.6) node[anchor=west, pos=0.95]{$x = 0$};
  % Axes.
  \draw[engblack!45, -{Latex[length=1.4mm]}] (-0.4, 0) -- (2.6, 0) node[right]{$x$};
  \draw[engblack!45, -{Latex[length=1.4mm]}] (0, -0.9) -- (0, 2.6) node[above]{$y$};
  % Unconstrained minimum (infeasible).
  \filldraw[enggray] (2, 2) circle (0.04);
  \node[enggray, anchor=south west] at (2, 2) {$(2,2)$ infeasible};
  % Perpendicular from (2,2) to the line, foot at (0.5, 0.5).
  \draw[perp] (2, 2) -- (0.5, 0.5);
  \filldraw[engred] (0.5, 0.5) circle (0.05);
  \node[engred, anchor=north east] at (0.5, 0.5) {$(\tfrac12,\tfrac12)$};
\end{tikzpicture}
\caption[KKT active-set geometry.]{The two-inequality example of the
active-set sweep: minimise $(x - 2)^{2} + (y - 2)^{2}$ over the feasible
region $x \geq 0$, $x + y \leq 1$ (green). The objective's contours are
circles centred on the unconstrained minimum $(2, 2)$, which lies outside
the feasible region. The constrained optimum sits at the foot of the
perpendicular from $(2, 2)$ to the active line $x + y = 1$, at
$(\tfrac{1}{2}, \tfrac{1}{2})$ (red), where the smallest feasible contour
is tangent to the constraint. Only the first inequality is active there;
the case sweep discards the other three active-set assumptions on a
negative multiplier or a primal-feasibility violation.}
\label{fig:vol02:ch11:kkt-activeset}
\end{figure}


\Cref{fig:vol02:ch11:kkt-activeset} draws the geometry of the surviving
case. The circular contours of the objective are centred on the
unconstrained minimum $(2, 2)$, which sits outside the feasible wedge;
the constrained optimum is the point of the boundary line $x + y = 1$
that the smallest feasible contour just touches, the foot of the
perpendicular from $(2, 2)$. Reading the picture, the reader sees why
only one inequality binds: the second constraint $x \geq 0$ is satisfied
with room to spare at $(\tfrac{1}{2}, \tfrac{1}{2})$, and forcing it
active drives the candidate to a corner where the objective gradient
points back into the interior, the negative-multiplier signature the
sweep rejected.

\begin{mastery}
The active-set sweep enumerates $2^{m}$ cases for $m$ inequalities,
which is tractable by hand for $m \leq 3$ and intractable for large
$m$. Practical solvers do not enumerate. The simplex method for
linear programs walks vertices of the feasible polytope, changing
the active set one constraint at a time; interior-point methods
follow a smooth path through the strict interior, letting the active
set emerge in the limit rather than committing to it. Both are the
algorithmic descendants of the KKT conditions, and both are the
subject of Volume II Chapter 15. The hand sweep of this section is
the conceptual skeleton: it shows what the solver is searching for,
the active set at which all four KKT conditions hold simultaneously.
\end{mastery}

\subsection{When KKT solves the problem and when it does not}

KKT conditions are necessary for an optimum in a smooth, well-
behaved program. They are also sufficient when the program is
\engterm{convex}: $f$ is convex, the equality constraints $g_{j}$ are
linear (i.e.\ affine), and the inequality constraint functions
$h_{\ell}$ are convex. Volume II Chapter 15 takes up convexity as
the dividing line between problems where KKT is the working
solution method and problems where it is one tool among several.
For the reader of this chapter, the working knowledge is: write the
Lagrangian, apply the four KKT conditions, examine the cases of
which inequalities are active, and verify the primal feasibility of
the candidate solution.

\subsection{Applied vignette: sizing a two-material heat-sink fin}
\pathtag{standard}

The optimisation apparatus of this chapter earns its keep on a design a
thermal engineer meets whenever a component must shed heat into moving
air. We size a rectangular cooling fin end to end: choose the length
$\ell$ and thickness $w$ (per unit fin depth) that maximise the heat the
fin rejects, subject to a fixed material budget and a manufacturing
lower bound on thickness. The worked value of the exercise is not the
final number but the way the KKT conditions read the trade between more
surface, which rejects more heat, and thicker metal, which conducts heat
further out along the fin before the air carries it away.

\paragraph{The model.} A straight fin of length $\ell$ and thickness $w$,
conductivity $k$, exposed to air with convection coefficient
$h_{\mathrm{c}}$ on both faces, rejects heat at a rate that the
one-dimensional fin equation of Volume IV gives as
\[
\dot{Q}(\ell, w) = \sqrt{2 h_{\mathrm{c}} k w}\;\Delta T\,
\tanh\!\left(\ell\sqrt{\frac{2 h_{\mathrm{c}}}{k w}}\right),
\]
with $\Delta T$ the base-to-air temperature difference. The prefactor
$\sqrt{2 h_{\mathrm{c}} k w}$ rises with thickness because a thicker fin
carries more conduction; the $\tanh$ factor saturates with length
because the tip of a long fin runs cool and rejects little. The material
budget caps the cross-section per unit depth,
\[
g(\ell, w) = \ell\,w - V_{0} = 0,
\]
and the process floor sets a minimum thickness,
\[
h_{1}(\ell, w) = w_{\min} - w \leq 0.
\]
We maximise $\dot{Q}$, equivalently minimise $f = -\dot{Q}$, subject to
the equality budget and the one inequality.

\paragraph{Estimate before optimising.} A very long, very thin fin has
$\tanh \to 1$ and $\dot{Q} \to \sqrt{2 h_{\mathrm{c}} k w}\,\Delta T$,
which rises with $w$; a very short fin has $\tanh(x) \approx x$ and
$\dot{Q} \to 2 h_{\mathrm{c}}\ell\,\Delta T$, independent of $w$ and
rising with $\ell$. The budget $\ell w = V_{0}$ forces a trade: spend
the fixed cross-section on length or on thickness. The reader expects an
interior optimum where the marginal heat per unit budget spent on length
equals the marginal heat per unit budget spent on thickness, unless the
thickness floor $w_{\min}$ binds first and pins $w$ at its minimum.

\paragraph{The KKT system.} With the equality multiplier $\lambda$ and
the inequality multiplier $\mu \geq 0$, stationarity of the Lagrangian
$\mathcal{L} = -\dot{Q} + \lambda(\ell w - V_{0}) + \mu(w_{\min} - w)$
reads
\[
-\pdv{\dot{Q}}{\ell} + \lambda w = 0, \qquad
-\pdv{\dot{Q}}{w} + \lambda \ell - \mu = 0,
\]
together with $\ell w = V_{0}$, $w \geq w_{\min}$, $\mu \geq 0$, and the
complementary slackness $\mu(w_{\min} - w) = 0$. Two cases decide the
design.

\textit{Case A: the thickness floor is slack} ($w > w_{\min}$, so
$\mu = 0$). Eliminating $\lambda$ from the two stationarity equations
gives the clean interior condition
\[
\frac{1}{w}\pdv{\dot{Q}}{\ell} = \frac{1}{\ell}\pdv{\dot{Q}}{w},
\]
which reads: at the optimum the heat gained per unit of budget moved
into length equals the heat gained per unit of budget moved into
thickness, exactly the equal-marginal-return rule the estimate
anticipated. Substituting the fin-equation partials and the budget
$\ell = V_{0}/w$ yields a single transcendental equation in $w$, solved
by one-dimensional root-finding. For a copper fin
($k = 400\,\si{\watt\per\meter\per\kelvin}$) in forced air
($h_{\mathrm{c}} = 80\,\si{\watt\per\meter\squared\per\kelvin}$) with
$V_{0} = 1.0\times 10^{-4}\,\si{\meter\squared}$ per unit depth, the
root sits near $w^{*} \approx 1.6\,\si{\milli\meter}$,
$\ell^{*} \approx 63\,\si{\milli\meter}$, an aspect ratio
$\ell/w \approx 40$: the optimal fin is long and thin, because the high
conductivity of copper keeps even a thin fin nearly isothermal along its
length.

\textit{Case B: the thickness floor binds} ($w = w_{\min}$, so
$\mu \geq 0$ and $\ell = V_{0}/w_{\min}$). This case governs a
low-conductivity fin (aluminium alloy, or a plastic heat spreader) or a
coarse process with a large $w_{\min}$. The second stationarity equation
returns the multiplier
$\mu = \lambda\ell - \partial\dot{Q}/\partial w$, and dual feasibility
$\mu \geq 0$ is the test that the floor genuinely binds. A negative
$\mu$ would mean the design wants a thinner fin than the process allows,
in which case Case A holds instead. The reader runs Case A first; if its
root falls below $w_{\min}$, the process floor binds and Case B pins the
thickness.

\paragraph{The reading.} The multiplier $\lambda$ is the shadow price of
the material budget: the extra heat rejected per unit of extra
cross-section, in $\si{\watt}$ per $\si{\meter\squared}$ of added metal.
A designer weighing a larger material allowance against the heat it buys
compares $\lambda$ against the marginal cost of the metal; the fin grows
until the two match. The inequality multiplier $\mu$, when the floor
binds, is the penalty the process minimum imposes: the heat the design
forfeits because it cannot make the fin as thin as the physics wants.
The two multipliers turn a heat-transfer formula into a procurement and
process decision, the recurring payoff of casting a design as a
constrained optimisation rather than a single-point calculation.

\section{Multiple integrals}\pathtag{core}

Volume II Chapter 6 introduced the integral as the limit of a
Riemann sum, the inverse of the derivative under the fundamental
theorem of calculus, and the formal home of accumulation. The
multiple integral generalises accumulation to higher-dimensional
domains. The double integral $\iint_{D} f \,\dd A$ is the
accumulation of $f$ over a planar region $D$; the triple integral
$\iiint_{V} f \,\dd V$ is accumulation over a spatial region $V$.

\subsection{The double integral as iterated integration}

For a region $D$ in $\R^{2}$ described by $a \leq x \leq b$ and
$y_{\text{low}}(x) \leq y \leq y_{\text{high}}(x)$, and a continuous
function $f(x, y)$,
\[
\iint_{D} f(x, y) \,\dd A = \int_{a}^{b}
\!\int_{y_{\text{low}}(x)}^{y_{\text{high}}(x)} f(x, y) \,\dd y\,
\dd x.
\]
The inner integral $\int f \,\dd y$ holds $x$ fixed and accumulates
over the $y$-slice; the outer integral accumulates the slices over
$x$.

\begin{theorem}[Fubini, working version]
If $f$ is continuous on a region $D$, the double integral can be
evaluated as either iterated integral, with the integration order
chosen by the reader,
\[
\int_{a}^{b}\!\int_{c}^{d} f \,\dd y\,\dd x =
\int_{c}^{d}\!\int_{a}^{b} f \,\dd x\,\dd y,
\]
on a rectangular domain, with the analogous statement on a more
general region.
\end{theorem}

The order of integration is a working choice. A function whose
$y$-integration is hard but whose $x$-integration is easy should be
evaluated $\int \dd x$ first; the reader of an integral whose inner
limits depend on the outer variable cannot swap orders without first
re-describing the region.

\paragraph{A worked example.} Compute $\iint_{D} (x + y) \,\dd A$
where $D = \set{(x, y) : 0 \leq x \leq 1,\ 0 \leq y \leq x}$. The
region is the triangle below the line $y = x$ in the unit square.
\[
\iint_{D} (x + y)\,\dd A = \int_{0}^{1}\!\int_{0}^{x}(x + y)\,\dd y\,
\dd x = \int_{0}^{1}\!\left(xy + \tfrac{y^{2}}{2}\right)\bigg|_{0}^{x}
\!\dd x = \int_{0}^{1}\tfrac{3 x^{2}}{2}\,\dd x = \tfrac{1}{2}.
\]

The same answer in the other order:
\[
\int_{0}^{1}\!\int_{y}^{1}(x + y)\,\dd x\,\dd y =
\int_{0}^{1}\!\left(\tfrac{x^{2}}{2} + xy\right)\bigg|_{y}^{1}\!\dd y
= \int_{0}^{1}\!\left(\tfrac{1}{2} + y - \tfrac{3 y^{2}}{2}\right)\dd y
= \tfrac{1}{2}.
\]
The orders agree; Fubini holds.
\Cref{fig:vol02:ch11:region-of-integration} draws the triangular
region with both descriptions: as $x$-slices from $y = 0$ to $y = x$
(the first order) and as $y$-slices from $x = y$ to $x = 1$ (the
second). Reading the region's two descriptions off the same picture
is the skill that prevents the order-swapping error of the failure
section.

% Figure: the triangular region 0<=y<=x<=1 with its two iterated
% descriptions: as x-slices (inner y from 0 to x) and as y-slices
% (inner x from y to 1). Supports the double-integral worked example
% and the order-swapping failure exercise.
\begin{figure}[htbp]
\centering
\begin{tikzpicture}[scale=2.6,
  region/.style={engblue!18},
  edge/.style={engblue, thick},
  slice/.style={engred, thick},
  axis/.style={engblack, thick, ->},
  label/.style={font=\small},
]
% Axes
\draw[axis] (-0.05, 0) -- (1.25, 0) node[right, label] {$x$};
\draw[axis] (0, -0.05) -- (0, 1.25) node[above, label] {$y$};

% The triangle with vertices (0,0), (1,0), (1,1): region 0<=y<=x<=1
\fill[region] (0, 0) -- (1, 0) -- (1, 1) -- cycle;
\draw[edge] (0, 0) -- (1, 0) -- (1, 1) -- cycle;
\node[label, engblue] at (1.02, 0.55) {$y = x$};

% Tick labels
\node[label, anchor=north] at (1, 0) {$1$};
\node[label, anchor=east] at (0, 1) {$1$};

% A vertical slice (x fixed, y from 0 to x): x-outer description
\draw[slice] (0.62, 0) -- (0.62, 0.62);
\filldraw[engred] (0.62, 0) circle (0.012);
\filldraw[engred] (0.62, 0.62) circle (0.012);
\node[label, engred, anchor=west] at (0.64, 0.33)
  {\small $y\!:0\!\to\! x$};
\node[label, anchor=north] at (0.62, -0.02) {\small $x$};

\end{tikzpicture}
\caption[Triangular region of integration.]{The region $D = \set{(x,
y) : 0 \leq y \leq x \leq 1}$, the triangle below the line $y = x$
in the unit square. The first iterated order takes vertical slices
(red): for each fixed $x$, the inner variable $y$ runs from $0$ to
$x$, and the outer $x$ runs from $0$ to $1$. The second order takes
horizontal slices: for each fixed $y$, the inner variable $x$ runs
from $y$ to $1$, and the outer $y$ runs from $0$ to $1$. Both
descriptions cover the same triangle; swapping the order requires
re-reading the bounds off the geometry, not merely exchanging the
integral signs.}
\label{fig:vol02:ch11:region-of-integration}
\end{figure}


\paragraph{Worked example: an integral that only one order can do.}
Some integrals are intractable in one order and elementary in the
other. Compute
\[
\int_{0}^{1}\!\int_{x}^{1} e^{y^{2}}\,\dd y\,\dd x.
\]
The inner integral $\int e^{y^{2}}\,\dd y$ has no elementary
antiderivative, so the stated order is a dead end. The region is
$\set{(x, y) : 0 \leq x \leq 1,\ x \leq y \leq 1}$, the triangle
above the line $y = x$. Re-described with $y$ outer, it is
$\set{(x, y) : 0 \leq y \leq 1,\ 0 \leq x \leq y}$. Swapping the
order,
\[
\int_{0}^{1}\!\int_{0}^{y} e^{y^{2}}\,\dd x\,\dd y
= \int_{0}^{1} y\,e^{y^{2}}\,\dd y
= \tfrac{1}{2}\,e^{y^{2}}\Big|_{0}^{1}
= \tfrac{1}{2}(e - 1) \approx 0.859.
\]
The inner $x$-integration is trivial because $e^{y^{2}}$ is constant
in $x$, and it supplies the factor $y$ that makes the outer integral
a simple substitution. The lesson: when one order stalls on a
missing antiderivative, redraw the region and try the other order
before reaching for a numerical method.

\paragraph{Worked example: the volume between two surfaces.} A double
integral computes the volume trapped between two graphs over a common
footprint, the upper surface's height minus the lower's, integrated
over the region where the upper lies above. Find the volume between
the plane $z = x + y$ and the paraboloid $z = x^{2} + y^{2}$ over the
region where the plane is the higher of the two. The two surfaces meet
where $x + y = x^{2} + y^{2}$, that is, $x^{2} - x + y^{2} - y = 0$,
which completing the square rewrites as $(x - \tfrac{1}{2})^{2} +
(y - \tfrac{1}{2})^{2} = \tfrac{1}{2}$, a disk of radius $1/\sqrt{2}$
centred at $(\tfrac{1}{2}, \tfrac{1}{2})$. Inside that disk the plane
lies above the paraboloid, so the volume is
\[
V = \iint_{D}\big[(x + y) - (x^{2} + y^{2})\big]\,\dd A.
\]
Shift the origin to the disk's centre with $x = \tfrac{1}{2} + u$,
$y = \tfrac{1}{2} + v$; the Jacobian of a pure translation is $1$, and
the integrand becomes, after expanding and cancelling the linear
terms, $\tfrac{1}{2} - (u^{2} + v^{2})$. In polar coordinates
$(s, \psi)$ on the shifted disk $s \in [0, 1/\sqrt{2}]$,
\[
V = \int_{0}^{2\pi}\!\int_{0}^{1/\sqrt{2}}
\Big(\tfrac{1}{2} - s^{2}\Big)\,s\,\dd s\,\dd\psi
= 2\pi\left[\frac{s^{2}}{4} - \frac{s^{4}}{4}\right]_{0}^{1/\sqrt{2}}
= 2\pi\Big(\tfrac{1}{8} - \tfrac{1}{16}\Big) = \frac{\pi}{8}.
\]
The translation that centred the region and the polar substitution
that exploited its circular symmetry each removed a difficulty the
raw Cartesian setup would have carried, the off-centre disk and the
$s\,\dd s$ weighting. The working habit is to locate the
surfaces' intersection first, which fixes both the footprint $D$ and
which surface is upper, and only then to choose the coordinate system
the footprint invites.

\begin{mastery}
Fubini's theorem needs a hypothesis the working version glosses: the
two iterated integrals agree, and equal the double integral, when
$\iint\abs{f}\,\dd A < \infty$, that is, when $f$ is absolutely
integrable over the region. Drop that hypothesis and the two orders can
genuinely disagree. The standard counterexample is
\[
f(x, y) = \frac{x^{2} - y^{2}}{(x^{2} + y^{2})^{2}}
\quad\text{on}\quad (0, 1]\times(0, 1],
\]
for which $\int_{0}^{1}\!\int_{0}^{1} f\,\dd x\,\dd y = -\tfrac{\pi}{4}$
while $\int_{0}^{1}\!\int_{0}^{1} f\,\dd y\,\dd x = +\tfrac{\pi}{4}$: the
same integrand, the same square, two opposite answers. The resolution is
that $\iint\abs{f}\,\dd A = \infty$ near the origin, so neither iterated
value is the double integral; there is no double integral to agree with.
The engineering relevance is not the pathology itself, which never
arises for a bounded integrand on a bounded region, but the discipline
it enforces on improper integrals. A physical integrand with a
singularity (a point charge's field, a crack-tip stress, a line source's
potential) can be non-absolutely-integrable near the singularity, and an
order swap that seems innocent can flip a sign or change a magnitude. The
working guard is to check absolute integrability, by bounding
$\iint\abs{f}$, before trusting a Fubini swap through a singular point;
where absolute integrability fails, the integral must be defined by an
explicit limiting process (a principal value, an excised
neighbourhood) rather than by an iterated computation.
\end{mastery}

\subsection{Triple integrals}

For a solid region $V \subset \R^{3}$,
\[
\iiint_{V} f(x, y, z) \,\dd V = \int_{a}^{b}\!\int_{y_{\text{low}}(x)}
^{y_{\text{high}}(x)}\!\int_{z_{\text{low}}(x, y)}^{z_{\text{high}}
(x, y)} f \,\dd z\,\dd y\,\dd x,
\]
when the region admits a description with iterated bounds.
The same Fubini permutation applies. Working integrals over
spheres, cylinders, and cones are usually more comfortable in
spherical or cylindrical coordinates, taken up in section 11.7.

\paragraph{Worked example: centroid of a solid cone.} A right
circular cone of base radius $R$ and height $H$ stands with its apex
at the origin and its axis along $z$, so at height $z$ the
cross-section is a disk of radius $\rho_{\max}(z) = (R/H)\,z$. In
cylindrical coordinates the solid is $0 \leq z \leq H$, $0 \leq
\rho \leq (R/H)z$, $0 \leq \phi \leq 2\pi$. The volume is the known
$\tfrac{1}{3}\pi R^{2} H$; confirm it,
\[
V = \int_{0}^{H}\!\int_{0}^{2\pi}\!\int_{0}^{(R/H)z}
\rho\,\dd\rho\,\dd\phi\,\dd z
= \int_{0}^{H} 2\pi\cdot\frac{1}{2}\Big(\frac{R}{H}z\Big)^{2}\dd z
= \frac{\pi R^{2}}{H^{2}}\cdot\frac{H^{3}}{3}
= \frac{\pi R^{2} H}{3}.
\]
The centroid lies on the axis by symmetry; its height is the
$z$-weighted average,
\[
\bar{z} = \frac{1}{V}\int_{0}^{H} z\cdot\pi\Big(\frac{R}{H}z\Big)^{2}
\dd z = \frac{\pi R^{2}/H^{2}}{\pi R^{2} H/3}\int_{0}^{H} z^{3}\,\dd z
= \frac{3}{H^{3}}\cdot\frac{H^{4}}{4} = \frac{3 H}{4}.
\]
The centroid of a solid cone sits three-quarters of the way from
apex to base, equivalently one-quarter of the height up from the
base. The result generalises the centroid of a triangle (one-third
of the height from the base in two dimensions): the extra dimension
shifts the balance point further toward the wide end, because the
cross-sectional area grows as $z^{2}$ rather than linearly. The
reader carries the two coefficients, $1/3$ for a triangle and $1/4$
from the base for a cone, as the standard centroids of tapering
shapes.

\subsection{Polar coordinates and the Gaussian}

Polar coordinates in $\R^{2}$ replace $(x, y)$ with $(r, \theta)$ via
$x = r\cos\theta$, $y = r\sin\theta$. The area element is
\[
\dd A = r \,\dd r\,\dd\theta.
\]
The factor of $r$ is the Jacobian of the transformation, derived
formally in section 11.7. The reader should remember it: a small
polar rectangle of dimensions $\dd r$ by $r\,\dd\theta$ has area
$r\,\dd r\,\dd\theta$. \Cref{fig:vol02:ch11:polar-area-element}
displays the elementary wedge whose area gives the factor.

% Figure: the polar area element dA = r dr dtheta. Show a small polar
% rectangle in (r, theta) and its image in (x, y), illustrating why
% the area picks up the factor r.
\begin{figure}[htbp]
\centering
\begin{tikzpicture}[scale=1.1,
  axis/.style={engblack, thick, ->},
  radial/.style={engblack!60, thin},
  arc/.style={engblue, thick},
  shaded/.style={fill=engblue!20, draw=engblue, thick},
  label/.style={font=\small},
]
% Axes
\draw[axis] (-0.3, 0) -- (4.2, 0) node[right, label] {$x$};
\draw[axis] (0, -0.3) -- (0, 4.0) node[above, label] {$y$};

% Polar grid radii
\foreach \rr in {1.0, 2.0, 3.0} {
  \draw[radial] (0,0) circle (\rr);
}
% Polar grid rays
\foreach \th in {0, 20, 40, 60, 80} {
  \draw[radial] (0,0) -- ({3.5*cos(\th)}, {3.5*sin(\th)});
}

% Highlight one polar cell: r in [2.0, 2.5], theta in [40, 60] degrees
\fill[engblue!25]
  ({2.0*cos(40)}, {2.0*sin(40)}) --
  plot[domain=40:60, samples=20] ({2.0*cos(\x)}, {2.0*sin(\x)}) --
  ({2.5*cos(60)}, {2.5*sin(60)}) --
  plot[domain=60:40, samples=20] ({2.5*cos(\x)}, {2.5*sin(\x)}) --
  cycle;
\draw[engblue, thick]
  ({2.0*cos(40)}, {2.0*sin(40)}) --
  plot[domain=40:60, samples=20] ({2.0*cos(\x)}, {2.0*sin(\x)}) --
  ({2.5*cos(60)}, {2.5*sin(60)}) --
  plot[domain=60:40, samples=20] ({2.5*cos(\x)}, {2.5*sin(\x)}) --
  cycle;

% Label dimensions
\draw[<->, engred, thick] ({2.0*cos(50)}, {2.0*sin(50)}) -- ({2.5*cos(50)}, {2.5*sin(50)});
\node[label, engred, anchor=west] at ({2.25*cos(50) + 0.05}, {2.25*sin(50) + 0.15}) {$\dd r$};

\draw[<->, engred, thick] ({2.25*cos(40) + 0.07*sin(40)}, {2.25*sin(40) - 0.07*cos(40)})
  arc[start angle=40, end angle=60, radius=2.25];
\node[label, engred, anchor=west] at ({2.35*cos(50)}, {2.35*sin(50) - 0.15}) {$r \,\dd\theta$};

% Annotations
\node[label, anchor=west] at (3.7, 2.0) {$\dd A = r\,\dd r\,\dd\theta$};

% A point at the cell
\filldraw[engblack] ({2.25*cos(50)}, {2.25*sin(50)}) circle (0.04);

% angle theta marker at the origin
\draw[engblack, thick] (0.6, 0) arc[start angle=0, end angle=50, radius=0.6];
\node[label, anchor=west] at (0.7, 0.25) {$\theta$};
% r marker
\draw[<->, engblack, thick] (0, -0.05) -- ({2.25*cos(50) - 0.05*cos(50)}, {2.25*sin(50) - 0.05*sin(50) - 0.05});
\node[label, anchor=north] at ({1.1*cos(50)}, {1.1*sin(50)}) {$r$};

\end{tikzpicture}
\caption[The polar area element.]{An infinitesimal polar rectangle
in $(r, \theta)$ maps to a wedge in $(x, y)$ with radial extent
$\dd r$ and angular extent $r\,\dd\theta$ (the arc length subtended
by $\dd\theta$ at radius $r$). The wedge has area $r\,\dd r\,\dd\theta$
to leading order. The factor $r$ is the Jacobian determinant of the
polar map; section 11.7 derives it from $\det J$.}
\label{fig:vol02:ch11:polar-area-element}
\end{figure}


\paragraph{The Gaussian integral.} Polar coordinates compute one of
the most important integrals in engineering and statistics:
\[
I = \int_{-\infty}^{\infty} e^{-x^{2}} \dd x.
\]
The integrand has no antiderivative in closed form; the integral
itself can nevertheless be computed by squaring it and turning the
result into a polar integral.
\[
I^{2} = \!\int_{-\infty}^{\infty}\!\int_{-\infty}^{\infty}
e^{-x^{2}} e^{-y^{2}} \dd x\,\dd y =
\!\int_{-\infty}^{\infty}\!\int_{-\infty}^{\infty}
e^{-(x^{2} + y^{2})} \dd x\,\dd y.
\]
Switch to polar coordinates,
\[
I^{2} = \int_{0}^{2\pi}\!\int_{0}^{\infty} e^{-r^{2}} r \,\dd r\,
\dd\theta = 2\pi \cdot \tfrac{1}{2} = \pi,
\]
where the inner integral $\int_{0}^{\infty} r e^{-r^{2}} \dd r$ is
evaluated by the substitution $u = r^{2}$. So $I = \sqrt{\pi}$. The
square of an integral with no closed form has a closed form when
viewed in polar coordinates. The result is the normalisation of the
Gaussian probability density of Volume II Chapter 13.

\subsection{Cylindrical and spherical integrals}

For triple integrals, the working coordinate systems are cylindrical
$(\rho, \phi, z)$ and spherical $(r, \theta, \phi)$. The volume
elements (Volume II Chapter 8) are
\begin{align*}
\dd V_{\text{cyl}} &= \rho \,\dd\rho\,\dd\phi\,\dd z, \\
\dd V_{\text{sph}} &= r^{2} \sin\theta\,\dd r\,\dd\theta\,\dd\phi.
\end{align*}
The reader keeps these to hand; they are the change-of-variable
factors for the most-used transformations of section 11.7.

% Figure: the spherical volume element. A small curvilinear box at radius r,
% polar angle theta, azimuth phi, with edges dr (radial), r dtheta (polar
% arc), and r sin(theta) dphi (azimuthal arc). The product is the Jacobian
% r^2 sin(theta) dr dtheta dphi.
\begin{figure}[htbp]
\centering
\begin{tikzpicture}[scale=1.0,
  axis/.style={engblack!55, -{Latex[length=1.8mm]}},
  edge/.style={engblue, very thick},
  arc/.style={engred, thick},
  guide/.style={enggray, thin, dashed},
]
  % Axes (a simple 3D-looking frame).
  \draw[axis] (0,0) -- (4.2, 0) node[right]{$x$};
  \draw[axis] (0,0) -- (0, 4.0) node[above]{$z$};
  \draw[axis] (0,0) -- (-2.2, -1.5) node[below left]{$y$};
  % A representative point P at radius r, polar angle theta from z-axis.
  % Draw the radial line and the little element around P.
  \coordinate (O) at (0,0);
  \coordinate (P) at (2.4, 1.9);
  \draw[guide] (O) -- (P);
  \node[engblack, anchor=south west] at (1.1, 0.95) {$r$};
  % Polar-angle arc near the origin between z-axis and OP.
  \draw[arc] (0, 1.1) arc (90:38.4:1.1);
  \node[engred] at (0.55, 1.35) {$\theta$};
  % The volume element: a small parallelepiped-like patch around P.
  % Edges: radial dr along OP direction, polar arc r dtheta, azimuth arc.
  \draw[edge] (P) -- ++(0.62, 0.49);            % radial dr (outward along OP)
  \draw[edge] (P) -- ++(-0.42, 0.53);           % polar arc r dtheta
  \draw[edge] (P) -- ++(-0.55, -0.24);          % azimuth arc r sin(theta) dphi
  % Close the little patch faintly.
  \draw[engblue, thick] (2.4+0.62, 1.9+0.49) -- ++(-0.42, 0.53);
  \draw[engblue, thick] (2.4-0.42, 1.9+0.53) -- ++(-0.55, -0.24);
  % Labels for the three edges.
  \node[engblue, anchor=south west] at (2.9, 2.35) {$\dd r$};
  \node[engblue, anchor=south east] at (1.9, 2.45) {$r\,\dd\theta$};
  \node[engblue, anchor=north east] at (1.9, 1.60)
    {$r\sin\theta\,\dd\phi$};
  \filldraw[engblack] (P) circle (0.045);
\end{tikzpicture}
\caption[The spherical volume element.]{The volume element in spherical
coordinates at a point of radius $r$, polar angle $\theta$, and azimuth
$\phi$. The three edges of the small curvilinear box are the radial
increment $\dd r$, the polar arc $r\,\dd\theta$, and the azimuthal arc
$r\sin\theta\,\dd\phi$ (shortened by the $\sin\theta$ factor because the
azimuthal circle at polar angle $\theta$ has radius $r\sin\theta$).
Their product is the Jacobian volume element $\dd V = r^{2}\sin\theta\,
\dd r\,\dd\theta\,\dd\phi$, the factor derived formally as the spherical
Jacobian determinant. The $\sin\theta$ shrinks the element toward the
poles, where the azimuthal circles collapse to points.}
\label{fig:vol02:ch11:spherical-element}
\end{figure}


\Cref{fig:vol02:ch11:spherical-element} builds the spherical volume
element from its three edges. The radial edge is $\dd r$, the polar edge
is the arc $r\,\dd\theta$, and the azimuthal edge is the arc
$r\sin\theta\,\dd\phi$, shortened by $\sin\theta$ because the circle of
constant polar angle $\theta$ has radius $r\sin\theta$ rather than $r$.
Their product recovers the $r^{2}\sin\theta$ factor, and the $\sin\theta$
that shrinks the element toward the poles is the reason a naive
equal-angle grid on a sphere over-samples the poles and under-samples the
equator, the working hazard the next worked example avoids by carrying
the correct volume element.

\paragraph{Worked example: the volume of a sphere.} The sphere of
radius $R$ in spherical coordinates is $0 \leq r \leq R$, $0 \leq
\theta \leq \pi$, $0 \leq \phi \leq 2\pi$. Its volume is
\[
V = \int_{0}^{2\pi}\!\int_{0}^{\pi}\!\int_{0}^{R} r^{2} \sin\theta
\,\dd r\,\dd\theta\,\dd\phi = 2\pi \cdot 2 \cdot \tfrac{R^{3}}{3} =
\tfrac{4\pi R^{3}}{3}.
\]
The same integral in Cartesian coordinates would require the
$z$-bounds $\pm\sqrt{R^{2} - x^{2} - y^{2}}$ and a polar substitution
in the $(x, y)$-plane; the spherical setup makes the bounds constant
and the integrand separable. The reader's choice of coordinate
system is engineering, not mere notation.

\paragraph{Worked example: a cylindrical integral.} Find the volume
of the solid bounded below by the paraboloid $z = x^{2} + y^{2}$ and
above by the plane $z = 4$. The region in the $(x, y)$-plane is the
disk $x^{2} + y^{2} \leq 4$ (where the paraboloid meets the plane),
and at each $(x, y)$ the height runs from $z = x^{2} + y^{2}$ up to
$z = 4$. Cylindrical coordinates $(\rho, \phi, z)$ make the bounds
clean: $\rho \in [0, 2]$, $\phi \in [0, 2\pi]$, $z \in [\rho^{2},
4]$, with $\dd V = \rho\,\dd z\,\dd\rho\,\dd\phi$. The volume is
\[
V = \int_{0}^{2\pi}\!\int_{0}^{2}\!\int_{\rho^{2}}^{4}
\rho\,\dd z\,\dd\rho\,\dd\phi
= 2\pi\int_{0}^{2}(4 - \rho^{2})\,\rho\,\dd\rho
= 2\pi\left[2\rho^{2} - \tfrac{\rho^{4}}{4}\right]_{0}^{2}
= 2\pi(8 - 4) = 8\pi.
\]
The same volume is half that of the circumscribing cylinder of radius
$2$ and height $4$ (whose volume is $16\pi$), the standard result
that a paraboloid encloses half the volume of its bounding cylinder.
Cylindrical coordinates turned a region with a curved lower boundary
into an iterated integral with elementary bounds.

\paragraph{Worked example: an integral over an off-centre disk.} Polar
coordinates are natural for a disk centred at the origin; a disk centred
elsewhere needs a shifted radial description. Compute $\iint_{D} x\,\dd A$
where $D$ is the disk $(x - 1)^{2} + y^{2} \leq 1$, tangent to the
$y$-axis at the origin. In polar coordinates about the origin, the
boundary circle $(x - 1)^{2} + y^{2} = 1$ expands to $x^{2} + y^{2} =
2x$, that is, $r^{2} = 2 r\cos\theta$, so $r = 2\cos\theta$ for
$\theta \in [-\pi/2, \pi/2]$. The region is swept by $r$ from $0$ to
$2\cos\theta$, and
\[
\iint_{D} x\,\dd A = \int_{-\pi/2}^{\pi/2}\!\int_{0}^{2\cos\theta}
(r\cos\theta)\,r\,\dd r\,\dd\theta
= \int_{-\pi/2}^{\pi/2}\cos\theta\,\frac{(2\cos\theta)^{3}}{3}\,
\dd\theta
= \frac{8}{3}\int_{-\pi/2}^{\pi/2}\cos^{4}\theta\,\dd\theta.
\]
The reduction $\int_{-\pi/2}^{\pi/2}\cos^{4}\theta\,\dd\theta = 3\pi/8$
(Wallis) gives $\tfrac{8}{3}\cdot\tfrac{3\pi}{8} = \pi$. The answer has a
quick sanity check: $\iint_{D} x\,\dd A = \bar{x}\cdot\abs{D}$, and the
centroid of a disk is its centre $\bar{x} = 1$ while its area is
$\abs{D} = \pi$, so the product is $\pi$, agreeing. The lesson is that
the polar equation $r = 2\cos\theta$ of a circle through the origin turns
an off-centre disk into a clean variable-upper-limit integral, a
substitution worth recognising whenever a circular region is tangent to
a coordinate axis.

\paragraph{Worked example: average distance from the centre of a
disk.} Compute the average distance of a point from the centre of a
disk of radius $R$. The average value of $r = \sqrt{x^{2} + y^{2}}$
over the disk is
\[
\bar{r} = \frac{1}{\pi R^{2}}\int_{0}^{2\pi}\!\int_{0}^{R}
r\cdot r\,\dd r\,\dd\theta
= \frac{1}{\pi R^{2}}\cdot 2\pi\cdot\frac{R^{3}}{3}
= \frac{2R}{3}.
\]
The average distance is $2R/3$, not $R/2$ as a naive reading of "the
middle" might suggest, because the area element $r\,\dd r\,\dd\theta$
weights large radii more heavily (there is more area at large $r$).
The same $r$-weighting is the source of the high-dimensional
shell-concentration of the failure section: in two dimensions it
nudges the average outward to $2R/3$; in a hundred dimensions it
pushes nearly all the measure to the boundary.

\paragraph{Worked example: mass of a sphere with graded density.}
The constant-density sphere is the easy case. Real spherical bodies,
a planet, a graded ceramic ball, a sintered powder sphere, carry a
density that varies with depth. Take a sphere of radius $R$ whose
density falls linearly from a central value $\rho_{c}$ to zero at the
surface, $\rho(r) = \rho_{c}(1 - r/R)$. Spherical coordinates make
the radial dependence separable. The mass is
\[
m = \int_{0}^{2\pi}\!\int_{0}^{\pi}\!\int_{0}^{R}
\rho_{c}\Big(1 - \frac{r}{R}\Big)\,r^{2}\sin\theta\,
\dd r\,\dd\theta\,\dd\phi.
\]
The angular integrals give $\int_{0}^{2\pi}\dd\phi = 2\pi$ and
$\int_{0}^{\pi}\sin\theta\,\dd\theta = 2$, a factor of $4\pi$, the
solid angle of the full sphere. The radial integral is
\[
\int_{0}^{R}\Big(1 - \frac{r}{R}\Big)r^{2}\,\dd r
= \int_{0}^{R}\Big(r^{2} - \frac{r^{3}}{R}\Big)\dd r
= \frac{R^{3}}{3} - \frac{R^{3}}{4} = \frac{R^{3}}{12}.
\]
So $m = 4\pi\rho_{c} R^{3}/12 = \pi\rho_{c} R^{3}/3$. Compare a
uniform sphere of the same central density, whose mass would be
$\rho_{c}\cdot\tfrac{4}{3}\pi R^{3} = 4\pi\rho_{c}R^{3}/3$: the
graded sphere holds one quarter the mass, because the linear taper
removes most of the material from the high-volume outer shells where
the $r^{2}$ weighting is largest. The mean density is
$\bar{\rho} = m/(\tfrac{4}{3}\pi R^{3}) = \rho_{c}/4$, not the
midpoint $\rho_{c}/2$ a naive average of centre and surface would
suggest, again because the $r^{2}$ volume weighting favours the
low-density outer region. The lesson recurs whenever a depth-graded
property is integrated over a sphere: the $r^{2}$ factor makes the
outer shells dominate, and the volume-weighted average sits well
below the arithmetic mean of the centre and surface values. Earth's
own mean density, about $5.5\,\si{\gram\per\centi\meter\cubed}$
against a surface-rock density near $2.7\,\si{\gram\per\centi\meter
\cubed}$ (current as of 2025), runs the other way because Earth's
density rises toward a dense iron core rather than falling, the
same integral with the sign of the gradient reversed.

\paragraph{Worked example: the moment of inertia of a solid sphere.}
The disk gave the coefficient $\tfrac{1}{2}$ and the ring gave $1$;
the sphere's coefficient $\tfrac{2}{5}$ is the one a rolling-body or a
spinning-rotor calculation needs, and the triple integral supplies it.
A uniform solid sphere of radius $R$, mass $m$, density $\rho =
m/(\tfrac{4}{3}\pi R^{3})$, spins about a diameter taken as the
$z$-axis. The moment of inertia is the mass-weighted integral of the
squared distance from that axis, and the distance from the $z$-axis of
a point at spherical radius $r$ and polar angle $\theta$ is $r\sin
\theta$, so the squared distance is $r^{2}\sin^{2}\theta$. Then
\[
I_{z} = \rho\int_{0}^{2\pi}\!\int_{0}^{\pi}\!\int_{0}^{R}
(r^{2}\sin^{2}\theta)\,r^{2}\sin\theta\,\dd r\,\dd\theta\,\dd\phi
= \rho\int_{0}^{2\pi}\!\dd\phi\int_{0}^{\pi}\!\sin^{3}\theta\,\dd\theta
\int_{0}^{R}\! r^{4}\,\dd r.
\]
The three factors separate cleanly. The azimuthal integral is $2\pi$;
the radial integral is $R^{5}/5$; the polar integral is $\int_{0}^{\pi}
\sin^{3}\theta\,\dd\theta = \tfrac{4}{3}$, the value that carries the
geometry of the $\sin^{2}\theta$ distance factor times the $\sin
\theta$ of the volume element. Multiplying,
\[
I_{z} = \rho\cdot 2\pi\cdot\frac{4}{3}\cdot\frac{R^{5}}{5}
= \frac{m}{\tfrac{4}{3}\pi R^{3}}\cdot\frac{8\pi R^{5}}{15}
= \frac{2}{5} m R^{2}.
\]
The coefficient $\tfrac{2}{5}$ is smaller than the disk's
$\tfrac{1}{2}$ because the sphere carries more of its mass near the
axis, where the poles pinch the cross-section, than a disk of the same
radius does. The reader now holds the three standard coefficients,
$\tfrac{2}{5}$ for a solid sphere, $\tfrac{1}{2}$ for a solid disk or
cylinder, and $1$ for a thin ring, which are exactly the numbers that
rank the finishing order of bodies rolling down an incline: the one
with the least coefficient converts the most of its energy into
translation and wins, so the sphere beats the cylinder beats the ring
(Volume III on rolling dynamics).

\subsection{Applications: mass, centroid, moment of inertia,
average value}

The multiple integral earns its place in engineering through four
recurring quantities. Each is an integral of a density-weighted
function over a region, and each reduces to a worked iterated
integral in the right coordinate system.

\paragraph{Mass.} A planar lamina with area density $\rho(x, y)$ (in
$\si{\kilogram\per\meter\squared}$) over a region $D$ has total mass
\[
m = \iint_{D} \rho(x, y)\,\dd A.
\]
For a solid with volume density $\rho(x, y, z)$, the mass is the
triple integral $\iiint_{V}\rho\,\dd V$. When $\rho$ is constant the
integral reduces to density times area or volume; the interest is in
non-uniform density, where the integral does the accounting the
constant case hides.

\paragraph{Centroid.} The \engterm{centroid} (centre of mass) has
coordinates that are the density-weighted averages of position,
\[
\bar{x} = \frac{1}{m}\iint_{D} x\,\rho\,\dd A, \qquad
\bar{y} = \frac{1}{m}\iint_{D} y\,\rho\,\dd A,
\]
with the obvious extension to three dimensions. The centroid is the
point at which the lamina balances; it is the first moment of the
mass distribution divided by the total mass. Volume III uses the
centroid as the point through which the resultant weight of a body
acts.

\paragraph{Moment of inertia.} The \engterm{moment of inertia} about
an axis measures resistance to angular acceleration about that axis.
For rotation about the $z$-axis through the origin, the moment of
inertia of a lamina is
\[
I_{z} = \iint_{D} (x^{2} + y^{2})\,\rho\,\dd A,
\]
the mass-weighted integral of the squared distance from the axis.
The $r^{2}$ weighting is why mass far from the axis contributes
disproportionately, the reason a flywheel concentrates its mass at
the rim and a figure skater pulling in the arms spins faster
(Volume III, conservation of angular momentum).

\paragraph{Average value.} The \engterm{average value} of a function
over a region is its integral divided by the measure of the region,
\[
\bar{f} = \frac{1}{\abs{D}}\iint_{D} f\,\dd A, \qquad
\abs{D} = \iint_{D}\dd A.
\]
The average temperature over a plate, the average illuminance over a
work surface, the mean stress over a cross-section: each is this
ratio. The average value generalises the one-variable mean-value
formula of Volume II Chapter 6.

\paragraph{Worked example: centroid of a quarter disk.} Find the
centroid of a uniform quarter disk of radius $R$ in the first
quadrant, $D = \set{(x, y) : x^{2} + y^{2} \leq R^{2},\ x, y \geq
0}$. By symmetry $\bar{x} = \bar{y}$, so compute $\bar{x}$. The area
is $\abs{D} = \pi R^{2}/4$. With constant density the centroid is
the geometric centroid,
\[
\bar{x} = \frac{1}{\abs{D}}\iint_{D} x\,\dd A
= \frac{4}{\pi R^{2}}\int_{0}^{\pi/2}\!\int_{0}^{R}
(r\cos\theta)\,r\,\dd r\,\dd\theta.
\]
The radial integral $\int_{0}^{R} r^{2}\,\dd r = R^{3}/3$, and the
angular integral $\int_{0}^{\pi/2}\cos\theta\,\dd\theta = 1$. So
\[
\bar{x} = \frac{4}{\pi R^{2}}\cdot\frac{R^{3}}{3}\cdot 1
= \frac{4 R}{3\pi} \approx 0.424\,R.
\]
The centroid sits at $(4R/3\pi, 4R/3\pi)$, closer to the corner than
the naive midpoint $(R/2, R/2)$ would suggest, because the disk has
more area at large radius along the diagonal. The number $4/(3\pi)
\approx 0.424$ is worth remembering; it recurs whenever a
quarter-circular or semicircular area's centroid is needed (a fillet,
a half-pipe channel, a D-shaped duct).

\paragraph{Worked example: moment of inertia of a disk.} The moment
of inertia of a uniform disk of radius $R$, mass $m$, and constant
density $\rho = m/(\pi R^{2})$ about its central axis is
\[
I_{z} = \iint_{D} r^{2}\,\rho\,\dd A
= \rho\int_{0}^{2\pi}\!\int_{0}^{R} r^{2}\cdot r\,\dd r\,\dd\theta
= \rho\cdot 2\pi\cdot\frac{R^{4}}{4}
= \frac{m}{\pi R^{2}}\cdot\frac{\pi R^{4}}{2}
= \frac{1}{2} m R^{2}.
\]
The factor $\tfrac{1}{2}$ is the signature of a solid disk; a thin
ring of the same mass and radius has $I = m R^{2}$ (all mass at the
rim), and a sphere has $I = \tfrac{2}{5} m R^{2}$. The reader who
internalises these three coefficients can estimate the rotational
inertia of most engineering bodies by approximating them as a
combination of disks, rings, and spheres, the standard move in
flywheel and rotor sizing (Volume III, Volume IX).

\paragraph{Worked example: mass of a plate with linear density.} A
square plate $[0, 1]^{2}$ (in metres) has area density $\rho(x, y) =
\rho_{0}(1 + x)$ that grows linearly along $x$, modelling a coating
that thickens across the plate. The mass is
\[
m = \int_{0}^{1}\!\int_{0}^{1}\rho_{0}(1 + x)\,\dd y\,\dd x
= \rho_{0}\int_{0}^{1}(1 + x)\,\dd x
= \rho_{0}\left(1 + \tfrac{1}{2}\right) = \tfrac{3}{2}\rho_{0}.
\]
The centroid's $x$-coordinate shifts toward the heavy edge,
\[
\bar{x} = \frac{1}{m}\int_{0}^{1}\!\int_{0}^{1} x\,\rho_{0}(1 + x)
\,\dd y\,\dd x = \frac{\rho_{0}}{\tfrac{3}{2}\rho_{0}}
\int_{0}^{1}(x + x^{2})\,\dd x
= \frac{2}{3}\left(\tfrac{1}{2} + \tfrac{1}{3}\right)
= \frac{5}{9} \approx 0.556,
\]
past the geometric centre $x = 0.5$ toward the thicker edge, as the
asymmetric density requires. The integral does the bookkeeping that
the constant-density shortcut cannot.

\subsection{Case study: centre of mass and rotational inertia of a
machined plate}\pathtag{standard}

The four integrals of the previous subsection converge on one task a
mechanical engineer performs whenever a part is to be balanced,
mounted, or spun: locate its centre of mass and compute its moment of
inertia. We work a real part end to end, a uniform steel plate with a
circular access hole, from the geometry to the dynamic consequence,
using only the multiple integral and one theorem it earns. The part
is a rectangular plate of width $a = 1.0\,\si{\meter}$, height $b =
0.6\,\si{\meter}$, and thickness $t = 8\,\si{\milli\meter}$, with a
circular hole of radius $R = 0.15\,\si{\meter}$ centred at $(x_{h},
y_{h}) = (0.70, 0.40)\,\si{\meter}$ measured from the lower-left
corner. The material is steel of density $\rho =
7850\,\si{\kilogram\per\meter\cubed}$.
\Cref{fig:vol02:ch11:bracket-inertia} fixes the geometry.

% Figure: a rectangular steel plate with a circular access hole, the
% part whose centre of mass and moment of inertia the case study of
% section 11.6 computes by composite subtraction. Dimensions in metres.
\begin{figure}[htbp]
\centering
\begin{tikzpicture}[scale=3.2,
  plate/.style={engblue!16, draw=engblue, thick},
  hole/.style={fill=white, draw=engred, thick},
  axis/.style={engblack, thick, ->},
  dim/.style={engblack, <->, >=stealth},
  label/.style={font=\small},
]
% Coordinate axes at the lower-left corner of the plate.
\draw[axis] (-0.05, 0) -- (1.15, 0) node[right, label] {$x$};
\draw[axis] (0, -0.05) -- (0, 0.85) node[above, label] {$y$};

% The plate: 1.0 m by 0.6 m rectangle in the first quadrant.
\fill[plate] (0, 0) rectangle (1.0, 0.6);

% Circular access hole of radius 0.15 m centred at (0.70, 0.40).
\filldraw[hole] (0.70, 0.40) circle (0.15);
\filldraw[engred] (0.70, 0.40) circle (0.012);
\node[label, engred, anchor=south] at (0.70, 0.56) {hole, $R=0.15$};

% Plate centroid marker.
\filldraw[engblue] (0.50, 0.30) circle (0.012);
\node[label, engblue, anchor=north] at (0.50, 0.285) {plate centroid};

% Width and height dimension lines.
\draw[dim] (0, -0.10) -- (1.0, -0.10);
\node[label, anchor=north] at (0.50, -0.10) {$a = 1.0$};
\draw[dim] (-0.12, 0) -- (-0.12, 0.6);
\node[label, anchor=east, rotate=90] at (-0.16, 0.30) {$b = 0.6$};

% Hole-centre coordinates.
\draw[engred, dashed, thin] (0.70, 0.40) -- (0.70, 0) ;
\draw[engred, dashed, thin] (0.70, 0.40) -- (0, 0.40);
\node[label, engred, anchor=north] at (0.70, -0.005) {\small $0.70$};
\node[label, engred, anchor=east] at (-0.005, 0.40) {\small $0.40$};

\end{tikzpicture}
\caption[Composite part: plate with a circular access hole.]{The
machined part of the case study: a uniform steel plate $a \times b =
1.0 \times 0.6\,\si{\meter}$ with a circular access hole of radius
$R = 0.15\,\si{\meter}$ centred at $(0.70, 0.40)\,\si{\meter}$. The
centre of mass and the moment of inertia are computed by composite
subtraction, treating the part as the full plate minus the disk the
hole removes. The plate's own centroid (blue) sits at the geometric
centre; the hole, off to one side, shifts the composite centroid away
from it.}
\label{fig:vol02:ch11:bracket-inertia}
\end{figure}


\paragraph{Step 0: estimate before integrating.} The plate occupies
$a b = 0.60\,\si{\meter\squared}$; the hole removes $\pi R^{2} =
\pi(0.15)^{2} = 0.0707\,\si{\meter\squared}$, about $12\%$ of the
area. The mass is therefore close to that of the full plate, $\rho t
a b = 7850\cdot 0.008\cdot 0.60 = 37.7\,\si{\kilogram}$, reduced by
roughly $12\%$ to about $33\,\si{\kilogram}$. The centroid should sit
near the plate centre $(0.50, 0.30)$ but be nudged away from the
hole, which is up and to the right; the nudge is small because the
removed area is small. The reader carries these two expectations,
$m \approx 33\,\si{\kilogram}$ and a centroid slightly below-left of
centre, into the calculation.

\paragraph{Step 1: the composite-subtraction principle.} A part with
a hole is the full part minus the material the hole removes. Every
additive quantity, area, mass, first moment, and second moment, obeys
the same subtraction, because each is an integral over the region and
the region of the holed part is the full rectangle minus the disk.
Write $P$ for the full plate, $D$ for the disk the hole removes, and
$P \setminus D$ for the actual part. For any integrand $\psi$,
\[
\iint_{P \setminus D}\psi\,\dd A = \iint_{P}\psi\,\dd A -
\iint_{D}\psi\,\dd A.
\]
The method avoids setting up an integral over the awkward
holed region directly; it computes two easy integrals over a
rectangle and a disk and subtracts. This composite move is the
working engineer's standard tactic for any region built from
elementary pieces by union and difference.

\paragraph{Step 2: mass.} With constant density and thickness, mass
is $\rho t$ times area. The plate area is $a b = 0.60\,\si{\meter
\squared}$ and the disk area is $\pi R^{2} = 0.0707\,\si{\meter
\squared}$, so the part area is $0.60 - 0.0707 =
0.5293\,\si{\meter\squared}$ and the mass is
\[
m = \rho t\,(a b - \pi R^{2}) = 7850\cdot 0.008\cdot 0.5293
= 33.2\,\si{\kilogram},
\]
matching the estimate. For the centroid and inertia it is convenient
to carry the two component masses: the full plate $m_{P} = \rho t a b
= 37.68\,\si{\kilogram}$ and the removed disk $m_{D} = \rho t\pi R^{2}
= 4.44\,\si{\kilogram}$, with $m = m_{P} - m_{D}$.

\paragraph{Step 3: centre of mass by first moments.} The
$x$-coordinate of the centroid is the first moment about the
$y$-axis divided by the mass. First moments subtract by the composite
principle. The full plate's centroid is its geometric centre $(a/2,
b/2) = (0.50, 0.30)$, so its first moment about the $y$-axis is
$m_{P}\cdot 0.50$. The disk's centroid is its centre $(0.70, 0.40)$,
so its first moment is $m_{D}\cdot 0.70$. Subtracting,
\[
\bar{x} = \frac{m_{P}\,x_{P} - m_{D}\,x_{h}}{m}
= \frac{37.68\cdot 0.50 - 4.44\cdot 0.70}{33.2}
= \frac{18.84 - 3.11}{33.2} = 0.474\,\si{\meter}.
\]
By the same construction,
\[
\bar{y} = \frac{m_{P}\,y_{P} - m_{D}\,y_{h}}{m}
= \frac{37.68\cdot 0.30 - 4.44\cdot 0.40}{33.2}
= \frac{11.30 - 1.78}{33.2} = 0.287\,\si{\meter}.
\]
The centre of mass is $(0.474, 0.287)\,\si{\meter}$, shifted from the
plate's geometric centre $(0.50, 0.30)$ down and to the left, away
from the hole, by about $26\,\si{\milli\meter}$ in $x$ and
$13\,\si{\milli\meter}$ in $y$. Removing mass from the upper-right
moves the balance point to the lower-left, the direction the estimate
anticipated, and by the small amount the small hole warrants. The
sign and the magnitude both pass inspection.

\paragraph{Step 4: moment of inertia about the centroidal axis.} The
quantity that governs how the part resists being spun about an axis
through its centre of mass, perpendicular to its face, is the polar
moment of inertia $I_{\bar{z}} = \iint(d^{2})\,\rho t\,\dd A$, where
$d$ is the distance from the centroidal axis. The composite principle
applies, but with a subtlety: each piece's inertia must be referred
to the same axis, the part's centroid, not each piece's own centroid.
The bridge is the \engterm{parallel-axis theorem}: the moment of
inertia of a body of mass $M$ about an axis offset by distance $s$
from a parallel axis through the body's own centroid is
\[
I = I_{\text{cm}} + M s^{2}.
\]
The theorem follows directly from expanding $d^{2} = (r' + s)^{2}$
inside the inertia integral, where $r'$ is measured from the body's
own centroid; the cross term integrates to zero because $r'$ averages
to zero about the centroid, leaving the centroidal inertia plus the
$M s^{2}$ transport term. We compute each piece's centroidal inertia
from the standard results of the previous subsection, then transport
both to the part's centroid.

The full rectangular plate has centroidal polar inertia $I_{P,
\text{cm}} = \tfrac{1}{12} m_{P}(a^{2} + b^{2})$, the sum of the two
in-plane bending inertias of a rectangle. Numerically,
\[
I_{P, \text{cm}} = \tfrac{1}{12}\cdot 37.68\cdot(1.0^{2} + 0.6^{2})
= \tfrac{1}{12}\cdot 37.68\cdot 1.36 = 4.270\,\si{\kilogram\meter
\squared}.
\]
Its centroid sits at $(0.50, 0.30)$, a distance $s_{P}$ from the
part centroid $(0.474, 0.287)$ given by $s_{P}^{2} = (0.50 -
0.474)^{2} + (0.30 - 0.287)^{2} = 0.026^{2} + 0.013^{2} = 8.45\times
10^{-4}\,\si{\meter\squared}$. Transporting,
\[
I_{P} = I_{P, \text{cm}} + m_{P}\,s_{P}^{2}
= 4.270 + 37.68\cdot 8.45\times 10^{-4}
= 4.270 + 0.032 = 4.302\,\si{\kilogram\meter\squared}.
\]
The removed disk has centroidal polar inertia $I_{D, \text{cm}} =
\tfrac{1}{2} m_{D} R^{2} = \tfrac{1}{2}\cdot 4.44\cdot 0.15^{2} =
0.0500\,\si{\kilogram\meter\squared}$, the solid-disk result of the
previous subsection. Its centre $(0.70, 0.40)$ sits a distance
$s_{D}$ from the part centroid given by $s_{D}^{2} = (0.70 - 0.474)
^{2} + (0.40 - 0.287)^{2} = 0.226^{2} + 0.113^{2} = 0.0639\,\si{
\meter\squared}$. Transporting,
\[
I_{D} = I_{D, \text{cm}} + m_{D}\,s_{D}^{2}
= 0.0500 + 4.44\cdot 0.0639
= 0.0500 + 0.2837 = 0.334\,\si{\kilogram\meter\squared}.
\]
The part's moment of inertia about its own centroidal axis is the
difference,
\[
I_{\bar{z}} = I_{P} - I_{D} = 4.302 - 0.334 =
3.97\,\si{\kilogram\meter\squared}.
\]

\paragraph{Step 5: interpretation.} The hole removes
$0.334\,\si{\kilogram\meter\squared}$ of rotational inertia, about
$8\%$ of the full plate's $4.30$, even though it removes only $12\%$
of the mass. The inertia removed is less than proportional to the
mass removed because the hole sits at intermediate radius rather than
at the rim; had the same hole been cut at the far corner, its $m_{D}
s_{D}^{2}$ transport term would have been larger and the inertia
reduction steeper. The dominant contribution to the disk's
transported inertia is the $m_{D} s_{D}^{2} = 0.284$ term, more than
five times its own centroidal $0.050$: the hole's distance from the
balance point matters far more than its own size, the same $r^{2}$
weighting that governs every moment of inertia. The engineering
payoff: if this plate is a rotor balanced and spun about its centre
of mass, a designer who wants to trim rotational inertia (to speed
up angular acceleration) removes material far from the axis; a
designer who wants to trim mass without losing inertia removes
material near the axis. The two goals pull in opposite radial
directions, and the moment-of-inertia integral is the instrument that
prices the trade. A final audit: the part inertia $3.97\,\si{
\kilogram\meter\squared}$ lies between the full-plate value $4.30$
and what remains, as a subtraction of a positive quantity must, and
its order of magnitude matches the crude estimate $\tfrac{1}{12} m
(a^{2} + b^{2}) = \tfrac{1}{12}\cdot 33.2\cdot 1.36 = 3.76$ for a
uniform plate of the part's mass, slightly below because the hole's
removal is biased away from the centroid.

\begin{estimation}
\textbf{Question.} A solid steel flywheel is a disk of radius
$R = 0.30\,\si{\meter}$ and thickness $t = 0.05\,\si{\meter}$.
Estimate its moment of inertia about the central axis, then the
kinetic energy stored when it spins at $3000$ revolutions per minute
(current as of 2026 for a small lab flywheel). Steel density is
about $7850\,\si{\kilogram\per\meter\cubed}$.

\textbf{Estimate.} The volume is $\pi R^{2} t = \pi (0.30)^{2}(0.05)
\approx 0.0141\,\si{\meter\cubed}$, so the mass is about $0.0141
\times 7850 \approx 111\,\si{\kilogram}$. For a solid disk $I =
\tfrac{1}{2} m R^{2}$, so $I \approx \tfrac{1}{2}(111)(0.09) \approx
5\,\si{\kilogram\meter\squared}$. Angular speed $3000$ revolutions
per minute is $3000 \cdot 2\pi/60 \approx 314\,\si{\radian\per
\second}$. The stored energy $\tfrac{1}{2} I \omega^{2} \approx
\tfrac{1}{2}(5)(314)^{2} \approx 2.5 \times 10^{5}\,\si{\joule}$,
roughly a quarter of a megajoule, comparable to the kinetic energy
of a small car at $30\,\si{\kilo\meter\per\hour}$.

\textbf{Calculate.} The mass is $m = \rho\,\pi R^{2} t = 7850 \cdot
\pi \cdot 0.09 \cdot 0.05 = 110.9\,\si{\kilogram}$. The moment of
inertia from the worked example above is $I = \tfrac{1}{2} m R^{2} =
\tfrac{1}{2}(110.9)(0.09) = 4.99\,\si{\kilogram\meter\squared}$. At
$\omega = 314.2\,\si{\radian\per\second}$, the kinetic energy is
\[
E = \tfrac{1}{2} I \omega^{2} = \tfrac{1}{2}(4.99)(314.2)^{2}
= 2.46 \times 10^{5}\,\si{\joule}.
\]
The estimate landed within a few percent. The engineering reading:
the $r^{2}$ weighting in the moment-of-inertia integral is what
makes the energy scale with $R^{2}\omega^{2}$; doubling the radius
at fixed mass quadruples the stored energy, the reason flywheel
energy storage favours large-radius rotors and the reason the burst
hazard of a failed flywheel scales just as steeply (Volume III on
rotating-machinery failure).
\end{estimation}

\section{Change of variables and Jacobians}\pathtag{standard}

The Jacobian factors $r$, $\rho$, $r^{2}\sin\theta$ that appeared in
the polar, cylindrical, and spherical area and volume elements have
a general form. They are the absolute values of the determinant of
the Jacobian matrix of the coordinate transformation.

\subsection{The change-of-variables theorem}

Let $\vect{T}: \R^{n} \to \R^{n}$ be a continuously differentiable
map sending a region $D'$ to a region $D$, with $\vect{T}$
one-to-one on the interior of $D'$. The \engterm{Jacobian matrix}
is
\[
J = J_{\vect{T}}(\vect{u}) = \begin{pmatrix}
\pdv{x_{1}}{u_{1}} & \cdots & \pdv{x_{1}}{u_{n}} \\
\vdots & \ddots & \vdots \\
\pdv{x_{n}}{u_{1}} & \cdots & \pdv{x_{n}}{u_{n}}
\end{pmatrix}.
\]
The \engterm{Jacobian determinant} is $\det J$.

\begin{theorem}[Change of variables]
Under the hypotheses above, for a continuous integrand $f$,
\[
\int_{D} f(\vect{x}) \,\dd V_{\vect{x}} =
\int_{D'} f(\vect{T}(\vect{u}))\,\abs{\det J(\vect{u})}\,
\dd V_{\vect{u}}.
\]
\end{theorem}

The absolute value of the Jacobian determinant is the local volume-
distortion factor of the map $\vect{T}$. A map that locally stretches
volume by a factor of two contributes a factor of two; a map that
locally reflects (negative determinant) contributes the absolute
value, because volume is unsigned.

\subsection{Polar, cylindrical, spherical Jacobians}

\paragraph{Polar.} $\vect{T}(r, \theta) = (r\cos\theta, r\sin\theta)$.
\[
J = \begin{pmatrix} \cos\theta & -r\sin\theta \\
\sin\theta & r\cos\theta \end{pmatrix}, \quad
\det J = r\cos^{2}\theta + r\sin^{2}\theta = r.
\]
The factor $r$ is what we used in section 11.6.

\paragraph{Cylindrical.} $\vect{T}(\rho, \phi, z) =
(\rho\cos\phi, \rho\sin\phi, z)$, $\det J = \rho$.

\paragraph{Spherical.} $\vect{T}(r, \theta, \phi) =
(r\sin\theta\cos\phi, r\sin\theta\sin\phi, r\cos\theta)$,
$\det J = r^{2} \sin\theta$.

The reader can derive each Jacobian by direct computation. The
working result is that whenever the engineer changes coordinates in
a multiple integral, the integrand picks up the absolute value of
the Jacobian determinant of the transformation. Forgetting the
factor is the most common error in change-of-variables and is the
working content of the failure section of Volume II Chapter 8 on
orientation: the same factor reports the local stretch of an oriented
volume element.

\subsection{Working example: an ellipse}

The area enclosed by the ellipse $x^{2}/a^{2} + y^{2}/b^{2} = 1$ is
\[
A = \iint_{D} \dd x\,\dd y.
\]
Substitute $x = a u$, $y = b v$. The Jacobian is $J = \diag(a, b)$,
$\det J = a b$. The image region is the unit disk $u^{2} + v^{2} \leq
1$. So
\[
A = \iint_{u^{2} + v^{2} \leq 1} a b \,\dd u\,\dd v = a b \cdot \pi.
\]
The familiar $\pi a b$ formula is the unit-disk area times the
linear-stretch determinant.

\subsection{Working example: a skewed integration region}

A linear change of variables turns a parallelogram into a rectangle.
Compute $\iint_{D}(x - y)\,\dd A$ where $D$ is the parallelogram with
vertices $(0, 0), (2, 1), (3, 3), (1, 2)$. The edges from the origin
are $(2, 1)$ and $(1, 2)$, so the natural variables are $u, v \in
[0, 1]$ with $\vect{T}(u, v) = u(2, 1) + v(1, 2) = (2u + v, u + 2v)$.
The Jacobian matrix is the matrix of the two edge vectors as columns,
\[
J = \begin{pmatrix} 2 & 1 \\ 1 & 2 \end{pmatrix},
\qquad \det J = 4 - 1 = 3.
\]
The integrand is $x - y = (2u + v) - (u + 2v) = u - v$, and the
region $D$ is the image of the unit square $[0, 1]^{2}$. So
\[
\iint_{D}(x - y)\,\dd A = \int_{0}^{1}\!\int_{0}^{1}(u - v)\cdot 3
\,\dd u\,\dd v = 3\left(\tfrac{1}{2} - \tfrac{1}{2}\right) = 0.
\]
The answer is zero by the antisymmetry of $u - v$ over the symmetric
square, a result the change of variables makes transparent and the
direct Cartesian setup (with four edge equations bounding $D$) would
obscure. The Jacobian determinant $3$ is the area of the
parallelogram (the unit square has area $1$, stretched by the factor
$\abs{\det J} = 3$), which the reader can confirm by the cross
product $\abs{(2, 1) \times (1, 2)} = \abs{4 - 1} = 3$.

\subsection{Working example: a curved substitution}

When the new coordinates are adapted to the integrand, the Jacobian
absorbs the geometry. Compute $\iint_{D} (x^{2} - y^{2})\,\dd A$ over
the region bounded by the hyperbolas $xy = 1$, $xy = 2$ and the lines
$y = x$, $y = 2x$ in the first quadrant. Choose $u = xy$ and $v =
y/x$, so the region is the rectangle $1 \leq u \leq 2$, $1 \leq v
\leq 2$ in the $(u, v)$-plane. Solving, $x = \sqrt{u/v}$ and $y =
\sqrt{uv}$. The inverse Jacobian is easier: $\partial(u, v)/
\partial(x, y) = \det\begin{pmatrix} y & x \\ -y/x^{2} & 1/x
\end{pmatrix} = y/x + y/x = 2y/x = 2v$, so the forward Jacobian is
$\abs{\partial(x, y)/\partial(u, v)} = 1/(2v)$. The integrand is
$x^{2} - y^{2} = u/v - uv = u(1/v - v)$. Then
\[
\iint_{D}(x^{2} - y^{2})\,\dd A
= \int_{1}^{2}\!\int_{1}^{2} u\Big(\tfrac{1}{v} - v\Big)\cdot
\frac{1}{2v}\,\dd u\,\dd v
= \frac{1}{2}\int_{1}^{2} u\,\dd u \int_{1}^{2}
\Big(\tfrac{1}{v^{2}} - 1\Big)\,\dd v.
\]
The $u$-integral is $\tfrac{3}{2}$; the $v$-integral is
$[-1/v - v]_{1}^{2} = (-\tfrac{1}{2} - 2) - (-1 - 1) = -\tfrac{1}{2}$.
The product is $\tfrac{1}{2}\cdot\tfrac{3}{2}\cdot(-\tfrac{1}{2}) =
-\tfrac{3}{8}$. Choosing coordinates aligned with the bounding curves
converted an awkward four-curve region into a rectangle, at the cost
of carrying the Jacobian $1/(2v)$. The trade is almost always worth
making: a constant-bounds integral with a Jacobian factor beats a
variable-bounds integral over a curved region.

\subsection{Working example: whitening a correlated Gaussian}

A change of variables earns its keep when a matrix transformation
turns an anisotropic, correlated integrand into an isotropic one. The
normalisation of a bivariate Gaussian probability density is the
canonical case, and it shows the Jacobian of a general linear map at
work. The density is
\[
p(\vect{x}) = \frac{1}{2\pi\sqrt{\det\Sigma}}
\exp\!\Big(-\tfrac{1}{2}\vect{x}^{\mathsf{T}}\Sigma^{-1}\vect{x}\Big),
\]
with $\Sigma$ a symmetric positive-definite $2\times 2$ covariance
matrix; the claim to verify is $\iint_{\R^{2}} p\,\dd A = 1$. The
quadratic form $\vect{x}^{\mathsf{T}}\Sigma^{-1}\vect{x}$ has
elliptical level sets, and the integral resists direct attack in
$(x, y)$. The remedy is to change variables to the coordinates in
which the form is the plain $\norm{\vect{u}}^{2}$. Because $\Sigma$
is symmetric positive-definite, it has a square root $\Sigma^{1/2}$
(the symmetric positive-definite matrix with $\Sigma^{1/2}\Sigma^{1/2}
= \Sigma$; concretely, diagonalise $\Sigma = Q\Lambda Q^{\mathsf{T}}$
and set $\Sigma^{1/2} = Q\Lambda^{1/2}Q^{\mathsf{T}}$). Substitute
$\vect{x} = \Sigma^{1/2}\vect{u}$, the \engterm{whitening}
transformation. Its Jacobian matrix is the constant matrix
$\Sigma^{1/2}$, so the Jacobian determinant is
\[
\det\big(\partial\vect{x}/\partial\vect{u}\big) = \det\Sigma^{1/2}
= \sqrt{\det\Sigma},
\]
using $\det(\Sigma^{1/2})^{2} = \det\Sigma$. The exponent
transforms cleanly,
\[
\vect{x}^{\mathsf{T}}\Sigma^{-1}\vect{x} =
(\Sigma^{1/2}\vect{u})^{\mathsf{T}}\Sigma^{-1}(\Sigma^{1/2}\vect{u})
= \vect{u}^{\mathsf{T}}\Sigma^{1/2}\Sigma^{-1}\Sigma^{1/2}\vect{u}
= \vect{u}^{\mathsf{T}}\vect{u} = \norm{\vect{u}}^{2},
\]
because $\Sigma^{1/2}\Sigma^{-1}\Sigma^{1/2} = I$. The change-of-
variables theorem then gives
\[
\iint_{\R^{2}} p\,\dd A
= \frac{1}{2\pi\sqrt{\det\Sigma}}\iint_{\R^{2}}
e^{-\norm{\vect{u}}^{2}/2}\,\sqrt{\det\Sigma}\,\dd u_{1}\,\dd u_{2}
= \frac{1}{2\pi}\iint_{\R^{2}} e^{-\norm{\vect{u}}^{2}/2}\,
\dd u_{1}\,\dd u_{2}.
\]
The $\sqrt{\det\Sigma}$ from the density's normalisation and the
$\sqrt{\det\Sigma}$ from the Jacobian cancel exactly, which is why
the normalising constant carries that determinant in the first place.
The remaining isotropic integral factors into two copies of the
one-dimensional Gaussian of section 11.6.3 (with the $1/2$ in the
exponent giving $\int e^{-u^{2}/2}\,\dd u = \sqrt{2\pi}$), so the
double integral is $(\sqrt{2\pi})^{2} = 2\pi$, and the whole
expression is $2\pi/(2\pi) = 1$. The density integrates to one for
every valid covariance, and the proof is the cancellation of the two
determinant factors. The same whitening transformation is the working
preprocessing step before principal component analysis and before
many statistical tests (Volume II Chapters 10 and 13): map the data
through $\Sigma^{-1/2}$ to standardise its covariance to the identity,
and the Jacobian determinant $\det\Sigma^{-1/2} = 1/\sqrt{\det\Sigma}$
is the volume rescaling that keeps probabilities consistent.

\begin{mastery}
The Jacobian determinant carries a sign, and the sign is the
orientation of the map. A transformation with $\det J > 0$ preserves
orientation (a right-handed frame maps to a right-handed frame); a
transformation with $\det J < 0$ reverses it (a reflection). The
change-of-variables theorem takes $\abs{\det J}$ because volume is
unsigned, but the signed determinant is the object that matters for
oriented integrals (the flux integrals of Volume II Chapter 8) and
for the divergence theorem's outward-normal convention. The failure
section of Volume II Chapter 8 traced a class of bugs to a dropped
or flipped sign in exactly this place. The working discipline:
compute $\det J$ with its sign, take the absolute value only at the
last step for an unsigned integral, and keep the sign whenever the
integral is oriented. A library that silently takes the absolute
value too early loses the orientation information the engineer needs
for a flux or a circulation.
\end{mastery}

\subsection{The Jacobian and the chain rule}

The Jacobian matrix is also the differential of the map. For a
function $f(\vect{x})$ composed with $\vect{x} = \vect{T}(\vect{u})$,
the chain rule reads, in matrix form,
\[
\grad_{\vect{u}}(f \circ \vect{T}) = J^{\mathsf{T}}\,
\grad_{\vect{x}} f.
\]
The two roles of the Jacobian, as the volume-distortion factor of
section 11.7.1 and as the matrix of the linearised map of section
11.2.5, are the same object viewed from two sides.

\begin{mastery}
The Jacobian determinant and the change-of-variables formula are the
shadow, in coordinates, of a coordinate-free statement in the
language of \engterm{differential forms}. The area element $\dd A =
\dd x\wedge\dd y$ is not a product of numbers but a \engterm{wedge
product}, an antisymmetric object with $\dd x\wedge\dd y = -\dd
y\wedge\dd x$ and $\dd x\wedge\dd x = 0$. Under a change of variables
$x = x(u, v)$, $y = y(u, v)$, the differentials pull back as $\dd x =
x_{u}\,\dd u + x_{v}\,\dd v$ and $\dd y = y_{u}\,\dd u + y_{v}\,\dd
v$, and the wedge product computes the Jacobian automatically,
\[
\dd x\wedge\dd y = (x_{u}\,\dd u + x_{v}\,\dd v)\wedge
(y_{u}\,\dd u + y_{v}\,\dd v)
= (x_{u} y_{v} - x_{v} y_{u})\,\dd u\wedge\dd v
= \det J\,\dd u\wedge\dd v,
\]
where the antisymmetry kills the $\dd u\wedge\dd u$ and $\dd v\wedge
\dd v$ terms and reverses the sign of one cross term, manufacturing
exactly the $2\times 2$ determinant. The signed determinant appears
here, not its absolute value, because the form carries orientation;
the absolute value enters only when the form is integrated over an
unoriented region. The same wedge algebra in $n$ variables produces
the $n\times n$ Jacobian determinant, and it unifies the
change-of-variables theorem with the curl, divergence, and Stokes
theorems of Volume II Chapter 8, all of which are the single
\engterm{generalised Stokes theorem} $\int_{\partial M}\omega =
\int_{M}\dd\omega$ written in different dimensions. The vector-valued
analogue of the Taylor expansion lives in the same framework: a map
$\vect{T} : \R^{n}\to\R^{m}$ has first-order term the Jacobian
matrix $J$ and second-order term a three-index array of mixed second
partials, the object that the matrix Hessian generalises when the
output is itself a vector. The reader does not need the formalism to
compute the integrals of this chapter; it is the language in which
the bookkeeping becomes automatic and the orientation sign stops
being a thing to remember.
\end{mastery}

\subsection{Applied vignette: the yield of a two-dimensional tolerance
box}\pathtag{standard}

A quality engineer sets tolerances on two coupled dimensions of a part
and needs the fraction of production that lands inside the acceptance
box. The two dimensions vary together, not independently, because a
common upstream cause (a temperature drift, a tool wear, a fixture
shift) moves both. The yield is a double integral of a correlated
Gaussian density over a rectangle, and the change-of-variables machinery
of this section turns an intractable elliptical integral into a product
of one-dimensional pieces. We work the calculation end to end and read
the result as a process-capability decision.

\paragraph{The model.} Let $(X_{1}, X_{2})$ be the two measured
deviations from nominal, jointly Gaussian with mean zero, standard
deviations $\sigma_{1}, \sigma_{2}$, and correlation coefficient
$\varrho$. The density is the bivariate normal of section 11.7.5 with
covariance
\[
\Sigma = \begin{pmatrix} \sigma_{1}^{2} & \varrho\sigma_{1}\sigma_{2}
\\ \varrho\sigma_{1}\sigma_{2} & \sigma_{2}^{2} \end{pmatrix}.
\]
The acceptance box is $\abs{X_{1}} \leq T_{1}$ and $\abs{X_{2}} \leq
T_{2}$, so the yield is
\[
Y = \iint_{[-T_{1}, T_{1}]\times[-T_{2}, T_{2}]}
\frac{1}{2\pi\sqrt{\det\Sigma}}
\exp\!\Big(-\tfrac{1}{2}\vect{x}^{\mathsf{T}}\Sigma^{-1}\vect{x}\Big)
\,\dd A.
\]
The exponent's elliptical level sets, tilted by the correlation, sit
askew to the axis-aligned box, and no elementary antiderivative resolves
the integral directly.

\paragraph{Estimate before integrating.} With independent dimensions
($\varrho = 0$) the box yield factors into two one-dimensional pieces,
$Y_{0} = [\,2\Phi(T_{1}/\sigma_{1}) - 1\,][\,2\Phi(T_{2}/\sigma_{2}) - 1
\,]$, where $\Phi$ is the standard normal cumulative distribution. For a
tolerance of three standard deviations on each dimension,
$T_{i} = 3\sigma_{i}$, each factor is $2\Phi(3) - 1 \approx 0.9973$, so
$Y_{0} \approx 0.9946$, a first-pass yield near $99.5\%$. Positive
correlation packs the joint density along a diagonal ridge that pokes out
of two corners of the box, so the reader expects a correlated yield
slightly below the independent estimate, the more so as $\varrho$ grows.

\paragraph{Whiten, then integrate.} The change of variables that
diagonalises the exponent is the whitening map $\vect{x} =
\Sigma^{1/2}\vect{u}$ of section 11.7.5, with Jacobian determinant
$\sqrt{\det\Sigma}$. In the whitened coordinates the density is the
isotropic standard Gaussian $\tfrac{1}{2\pi}e^{-\norm{\vect{u}}^{2}/2}$,
and the two determinant factors cancel exactly, as shown there. The box,
though, is no longer a rectangle in $\vect{u}$: the whitening map sends
the axis-aligned box to a tilted parallelogram, so the integral becomes
\[
Y = \frac{1}{2\pi}\iint_{\Sigma^{-1/2}(\text{box})}
e^{-\norm{\vect{u}}^{2}/2}\,\dd u_{1}\,\dd u_{2}.
\]
The integrand is now radially symmetric and elementary, and the geometry
of the tilted parallelogram carries all the correlation. For the
special case of equal tolerances scaled to the marginals,
$T_{i} = c\,\sigma_{i}$, and moderate correlation, the yield reduces to
a standard bivariate-normal box probability that a table or a short
numerical quadrature returns. A representative figure: with
$\varrho = 0.6$ and $T_{i} = 3\sigma_{i}$, the correlated yield is about
$0.991$, half a percentage point below the independent $0.9946$,
because the diagonal ridge of the density spills out of the two
off-diagonal corners of the box that the axis-aligned tolerance fails to
capture.

\paragraph{The reading.} The half-percent yield gap between the
independent estimate and the correlated truth is money: on a run of a
million parts it is five thousand rejects the independent model did not
predict. The engineering fix follows from the geometry the whitening
exposed. Because the loss comes from the corners where the density ridge
exits the box, rotating the acceptance region to align with the density's
principal axes (an acceptance parallelogram rather than an axis-aligned
box) recovers most of the lost yield without loosening either marginal
tolerance. The rotated region is exactly the box in the whitened
coordinates, and the same $\Sigma^{-1/2}$ that made the integral
tractable is the transformation that specifies the better acceptance
region. The double integral does not merely report the yield; through
its change of variables it prescribes the shape of the tolerance region
that maximises it, which is the payoff of computing the yield as a
correlated integral rather than a product of one-dimensional pass rates.

\section{Failure: the high-dimensional space and our intuition's
collapse}\pathtag{core}

Two-variable and three-variable problems admit visualisation. The
reader can sketch contours of $f(x, y)$, can imagine the gradient
of a temperature field, can picture the volume of a sphere. As
$n$ grows, this visual intuition stops being a guide. By $n = 10$
the reader has already lost the picture; by $n = 100$ several
intuitions are not merely weak but actively misleading. The
following are the standard ways high-dimensional geometry breaks
the intuitions of low-dimensional engineering. Treat them as
generic results that recur whenever a problem has many parameters,
many features, or many degrees of freedom.

\subsection{Volume concentrates near the boundary}

The volume of an $n$-dimensional ball of radius $R$ scales as
$V_{n}(R) = c_{n} R^{n}$ where $c_{n}$ is a dimension-dependent
constant. The volume of the spherical shell of inner radius
$(1 - \varepsilon) R$ and outer radius $R$ is
\[
V_{\text{shell}} = c_{n} R^{n} - c_{n} (1 - \varepsilon)^{n} R^{n}
= V_{n}(R)\left[1 - (1 - \varepsilon)^{n}\right].
\]
For fixed $\varepsilon = 0.01$ (the outer 1\% of the radius),
\begin{itemize}[leftmargin=*]
  \item $n = 3$: $1 - 0.99^{3} \approx 0.030$. About 3\% of the
    volume sits in the outer shell.
  \item $n = 100$: $1 - 0.99^{100} \approx 0.634$. About 63\% of
    the volume sits in the outer shell.
  \item $n = 1000$: $1 - 0.99^{1000} \approx 0.99996$. Almost all
    the volume sits in the outer shell.
\end{itemize}
In high dimensions, almost every point of the ball is within
$\varepsilon R$ of its surface. The reader's intuition that the
"middle of the ball" contains most of the volume is a low-
dimensional artefact. The engineering consequence: random sampling
inside a high-dimensional region picks up samples that are
overwhelmingly close to the boundary. A Monte Carlo simulation that
relies on uniform sampling away from boundaries gives misleading
results in high dimensions.

% Full-page figure: high-dimensional volume concentration.
% Two stacked plots:
%   (a) Fraction of n-ball volume within outer shell of width
%       epsilon, as a function of n, for epsilon = 0.01, 0.05, 0.1.
%   (b) Ratio of inscribed-ball volume to enclosing-cube volume,
%       V_n(1)/2^n, on a log scale, n = 1..50.
\begin{figure}[p]
\centering
\begin{tikzpicture}[
  every axis/.style={
    font=\small,
    axis lines=left,
    width=12cm,
    height=6.5cm,
    grid=both,
    grid style={engblack!12, very thin},
    major grid style={engblack!22},
    enlargelimits=false,
    tick style={engblack},
    label style={font=\small},
  },
]

% ---- Panel (a): outer-shell fraction vs dimension n ----
\begin{axis}[
  name=shellplot,
  xlabel={dimension $n$},
  ylabel={fraction in outer shell $1 - (1-\varepsilon)^{n}$},
  xmin=0, xmax=500,
  ymin=0, ymax=1.02,
  legend pos=south east,
  legend cell align=left,
  legend style={font=\footnotesize, draw=engblack!30},
  title={(a) Volume of $n$-ball inside outer shell of width $\varepsilon R$},
]
  % epsilon = 0.01
  \addplot[engblue, very thick, domain=1:500, samples=200] {1 - (0.99)^x};
  \addlegendentry{$\varepsilon = 0.01$}
  % epsilon = 0.05
  \addplot[engred, very thick, dashed, domain=1:500, samples=200] {1 - (0.95)^x};
  \addlegendentry{$\varepsilon = 0.05$}
  % epsilon = 0.10
  \addplot[engblack, very thick, dotted, domain=1:500, samples=200] {1 - (0.90)^x};
  \addlegendentry{$\varepsilon = 0.10$}
  % 50% reference line
  \addplot[engblack!40, thin, domain=0:500] {0.5};
\end{axis}

% ---- Panel (b): inscribed-ball / cube ratio, log scale ----
\begin{axis}[
  name=ratioplot,
  at={(shellplot.below south west)},
  anchor=above north west,
  yshift=-1.4cm,
  xlabel={dimension $n$},
  ylabel={$V_{n}(1) / 2^{n}$ (log scale)},
  xmin=1, xmax=50,
  ymode=log,
  ymin=1e-30, ymax=1.5,
  title={(b) Ratio of inscribed unit ball to enclosing cube $[-1, 1]^{n}$},
]
  % V_n(1)/2^n = pi^(n/2) / (2^n Gamma(n/2 + 1))
  % Use the recursive log form is awkward in pgfplots; precompute as
  % a coordinates table. Values from V_n(1) = pi^(n/2)/Gamma(n/2+1).
  \addplot[engblue, very thick, mark=*, mark size=1.2pt] coordinates {
    (1, 1.0)
    (2, 0.7854)
    (3, 0.5236)
    (4, 0.3084)
    (5, 0.1645)
    (6, 0.0807)
    (7, 0.0369)
    (8, 0.0159)
    (9, 0.00641)
    (10, 0.00249)
    (12, 3.26e-4)
    (15, 1.43e-5)
    (20, 2.46e-8)
    (25, 1.84e-11)
    (30, 6.65e-15)
    (40, 3.28e-22)
    (50, 5.10e-30)
  };
\end{axis}

\end{tikzpicture}
\caption[High-dimensional volume concentration in the ball and the
cube.]{(a) The fraction of the $n$-dimensional ball that sits in the
outer shell of relative thickness $\varepsilon$ is $1 - (1 -
\varepsilon)^{n}$. For $\varepsilon = 0.01$ (solid), the fraction
crosses $50\%$ near $n \approx 70$ and exceeds $99\%$ by $n
\approx 500$. (b) The volume of the unit ball inscribed in the cube
$[-1, 1]^{n}$ collapses to a negligible fraction of the cube as $n$
grows: the ratio is $\pi/4 \approx 0.79$ at $n = 2$ and drops below
$10^{-8}$ by $n = 20$. Together the two plots are the working
illustration of the failure section. Random sampling in
high-dimensional geometry should not assume that a "typical" sample
lies near the centre of the domain or even near the inscribed ball.}
\label{fig:vol02:ch11:volume-concentration}
\end{figure}


\subsection{Volume concentrates near the equator}

A second concentration applies to surface area. On the unit sphere
$S^{n - 1}$ in $\R^{n}$, almost all the surface area sits within a
narrow band around any equator (any $(n - 1)$-dimensional great
hyperplane through the origin). For $n = 100$, more than 99\% of the
surface area sits within a band of angular half-width about
$0.2$ radians (about $12^{\circ}$) around the equator. The
phenomenon is sometimes called \engterm{measure concentration} on
the sphere; for the working engineer, the consequence is that random
unit vectors in high dimensions are nearly orthogonal to any fixed
reference vector with high probability.

\begin{mastery}
The shell concentration has a probabilistic twin that governs every
high-dimensional Gaussian the engineer samples. Take
$\vect{X} = (X_{1}, \ldots, X_{n})$ with independent standard-normal
components. The squared length $\norm{\vect{X}}^{2} = \sum_{i} X_{i}^{2}$
is a sum of $n$ independent unit-variance terms, each with mean $1$, so
by the law of large numbers $\norm{\vect{X}}^{2} \approx n$ and
$\norm{\vect{X}} \approx \sqrt{n}$. The fluctuation is small in relative
terms: $\Var[\norm{\vect{X}}^{2}] = 2n$, so the standard deviation of
$\norm{\vect{X}}^{2}$ is $\sqrt{2n}$, a relative spread
$\sqrt{2n}/n = \sqrt{2/n}$ that vanishes as $n$ grows. Almost all the
probability mass of a high-dimensional standard Gaussian sits in a thin
shell at radius $\sqrt{n}$, not near the mode at the origin, even though
the density is largest at the origin. The density is largest there, but
there is almost no volume there; the shell at $\sqrt{n}$ has modest
density and overwhelming volume, and volume wins. This is the
\engterm{soap-bubble} or \engterm{Gaussian annulus} phenomenon, and it
overturns two intuitions at once: the typical sample of a
high-dimensional Gaussian is far from the mean, and two independent
samples are nearly orthogonal and nearly equidistant, at pairwise
distance about $\sqrt{2n}$. The engineering consequence surfaces in
sampling-based methods (Volume II Chapter 16), in the initialisation of
high-dimensional models (Volume VII), and in the geometry of noise added
for privacy or regularisation: a Gaussian perturbation in $n$ dimensions
moves a point a distance $\approx\sigma\sqrt{n}$, not $\sigma$, which the
engineer who budgets a noise scale by its per-coordinate standard
deviation will underestimate by the factor $\sqrt{n}$.
\end{mastery}

\subsection{The unit cube and its corners}

The reader's picture of a cube has corners that contain a non-
trivial fraction of the volume. In two dimensions, a unit square
has four corners; the four quarter-disks of radius $0.5$ at the
corners cover area $\pi (0.5)^{2} = \pi/4 \approx 0.785$, so the
corners-without-the-inscribed-disk hold $0.215$, about 22\% of the
square's area.

In $n$ dimensions, the inscribed ball in the unit cube $[-1, 1]^{n}$
has volume $V_{n}(1) = c_{n}$, while the cube itself has volume
$2^{n}$. The ratio,
\[
\frac{V_{n}(1)}{2^{n}} = \frac{\pi^{n/2}}{2^{n}\,\Gamma(n/2 + 1)},
\]
shrinks rapidly. For $n = 2$, the ratio is $\pi / 4 \approx 0.785$.
For $n = 10$, the ratio is approximately $0.0025$. For $n = 100$,
the ratio is astronomically small, of order $10^{-70}$. The
intuition that the inscribed ball captures most of the cube is
a low-dimensional artefact. In high dimensions, the cube is almost
all corners. A search routine that relies on the inscribed ball as
a reasonable approximation to the cube is wrong by an exponential
factor.

\paragraph{Worked example: how far the corner reaches.} The corners are
not only voluminous, they are far. The centre of the unit cube
$[-1, 1]^{n}$ is the origin; the nearest face is at distance $1$ (the
inscribed-ball radius), but the farthest corner $(1, 1, \ldots, 1)$ is at
distance $\sqrt{n}$. The ratio of the corner distance to the face
distance is $\sqrt{n}$, so it grows without bound. For $n = 2$ the corner
is $\sqrt{2} \approx 1.41$ times as far as the face; for $n = 100$ it is
$10$ times as far; for $n = 10000$ it is $100$ times as far. A cube that
looks nearly round in two or three dimensions (corner only $40\%$ or
$70\%$ farther than face) is a spiked object in high dimensions, with
$2^{n}$ corners each poking out to radius $\sqrt{n}$ while the faces stay
at radius $1$. The engineering consequence compounds the volume result:
an optimiser or sampler that steps a fixed distance from the centre
reaches the faces long before it reaches any corner, so it never visits
the regions where nearly all the cube's volume lives. Box-constrained
search that treats each coordinate's bound independently must reckon with
the fact that the joint corner is exponentially far and exponentially
voluminous compared to the naive spherical picture.

\subsection{Distance and neighborhood concepts collapse}

A consequence of the above is that the notion of \engterm{nearest
neighbour} loses meaning. Take $N$ random points uniformly in the
unit cube $[0, 1]^{n}$. As $n$ grows, the ratio
\[
\frac{\max_{i, j} \norm{\vect{x}_{i} - \vect{x}_{j}}}
{\min_{i, j, i \neq j} \norm{\vect{x}_{i} - \vect{x}_{j}}}
\to 1
\]
in distribution. The farthest pair of points is approximately the
same distance apart as the closest pair. The notion of "nearby" is
not well-defined: every point is approximately equidistant from
every other point. Algorithms that rely on local neighborhoods (
$k$-nearest-neighbour, kernel density estimation, local
interpolation) lose their statistical power as the ambient
dimension grows. This is the working content of the
\engterm{curse of dimensionality} in machine learning
\cite{paper:v2ch11-blum-curse-1997}.

% Figure: the collapse of relative distance contrast in high dimensions. For
% N uniform points in the unit cube, the ratio (d_max - d_min)/d_min of the
% spread of pairwise distances to the smallest one falls toward zero as the
% dimension grows: every pair becomes nearly equidistant. Plotted from a
% representative concentration curve proportional to 1/sqrt(n).
\begin{figure}[htbp]
\centering
\begin{tikzpicture}[scale=1.0]
  \begin{axis}[
    width=9.4cm, height=6.2cm,
    xlabel={ambient dimension $n$},
    ylabel={relative distance contrast},
    xmode=log, log basis x=10,
    xmin=1, xmax=1000, ymin=0, ymax=1.6,
    axis lines=left,
    xtick={1,3,10,30,100,300,1000},
    xticklabels={1,3,10,30,100,300,1000},
    ytick={0,0.4,0.8,1.2,1.6},
    tick label style={font=\small},
    label style={font=\small},
    legend style={font=\small, at={(0.97,0.97)}, anchor=north east, draw=engblack!40},
  ]
    % Contrast ~ C/sqrt(n): the concentration-of-measure scaling. C chosen so
    % the n=2 value sits near 1.4, a plausible finite-sample contrast.
    \addplot[engblue, very thick, domain=1:1000, samples=80] {2.0/sqrt(x)};
    \addlegendentry{$\propto 1/\sqrt{n}$ (theory)}
    % A few "empirical" markers to suggest a simulation, tracking the curve.
    \addplot[engred, only marks, mark=*, mark size=1.7pt] coordinates {
      (2, 1.41) (5, 0.90) (10, 0.63) (30, 0.37) (100, 0.20) (300, 0.115) (1000, 0.063)
    };
    \addlegendentry{sampled ($N=500$)}
  \end{axis}
\end{tikzpicture}
\caption[Distance contrast collapses with dimension.]{The relative
distance contrast among $N$ uniform random points in the unit cube
$[0, 1]^{n}$, measured as the spread of pairwise distances divided by the
smallest, plotted against ambient dimension $n$ on a logarithmic axis.
The contrast falls as $1/\sqrt{n}$ (blue), so by $n = 100$ the farthest
pair of points is only about a fifth further apart than the closest pair,
and by $n = 1000$ every pair is nearly equidistant. The red markers show
the trend a finite-sample simulation would produce. When contrast
vanishes, the notion of a nearest neighbour loses meaning, which is the
geometric root of the curse of dimensionality for local methods.}
\label{fig:vol02:ch11:distance-collapse}
\end{figure}


\Cref{fig:vol02:ch11:distance-collapse} plots the relative distance
contrast against dimension. The contrast falls as $1/\sqrt{n}$, so the
gap between the farthest and the nearest pair of points shrinks steadily
until, by a few hundred dimensions, every pair sits at nearly the same
distance. A nearest-neighbour rule reads the label of the closest
training point; when the closest point is barely closer than the
farthest, the rule reads noise, which is why the failure-analysis
exercise on the collapse of a nearest-neighbour classifier follows
directly from this curve.

\subsection{The second-derivative test at a degenerate point}

A second, smaller-scale failure lives inside the classification
machinery itself. The second-derivative test of section 11.3 is
silent when the Hessian has a zero eigenvalue: a vanishing
determinant ($\det H = 0$ in two variables) means the quadratic form
$\dd\vect{x}^{\mathsf{T}} H\,\dd\vect{x}$ is degenerate, and the
leading-order Taylor term no longer decides the local shape. The
reader who reports "inconclusive" and stops has not finished the
analysis; the engineer who reports "minimum" on the strength of a
zero determinant has guessed.

\paragraph{Worked example: the monkey saddle.} Let $f(x, y) = x^{3}
- 3 x y^{2}$. The gradient $\grad f = (3 x^{2} - 3 y^{2}, -6 x y)$
vanishes only at the origin. The Hessian is
\[
H = \begin{pmatrix} 6 x & -6 y \\ -6 y & -6 x \end{pmatrix},
\]
which at the origin is the zero matrix: both eigenvalues zero, the
test fully inconclusive. The shape is decided by the cubic term.
Along the $x$-axis, $f(x, 0) = x^{3}$ takes both signs; the origin
is neither a maximum nor a minimum. In polar coordinates $f =
r^{3}\cos 3\theta$, which has three positive lobes and three
negative lobes alternating around the origin, the \engterm{monkey
saddle} (room for two legs and a tail). A gradient-based optimiser
that lands here reads $\grad f = \vect{0}$ and a flat Hessian and
cannot tell, from second-order information alone, which way is down.
The remedy is to inspect the function along several rays (the
direction probe of the diagnosis exercises) or to compute the
higher-order Taylor term.

\paragraph{Worked example: a degenerate minimum that is not a
minimum.} Let $f(x, y) = x^{2} + y^{4}$. At the origin, $\grad f =
(2x, 4 y^{3}) = \vect{0}$ and $H = \diag(2, 0)$: one positive
eigenvalue, one zero. The test is inconclusive. Here the higher-order
term rescues the classification: $f = x^{2} + y^{4} \geq 0$ with
equality only at the origin, so the origin is a genuine local (and
global) minimum, but the second-derivative test cannot see it
because the $y$-direction is quartic-flat. Contrast $g(x, y) = x^{2}
- y^{4}$, which has the same Hessian $\diag(2, 0)$ at the origin yet
is a saddle (negative along the $y$-axis). Two functions with
identical gradients and identical Hessians at the critical point
have opposite classifications. No purely second-order test can
distinguish them; the distinction lives in the fourth-order term.
This is the precise sense in which the second-derivative test fails
at a degenerate point, and it is why a careful analysis reports the
order of vanishing, not just the Hessian.

\begin{principle}[Degenerate critical points need higher order]
When the Hessian at a critical point is singular, the
second-derivative test is not merely inconclusive as a matter of
convention; the local shape genuinely is not determined by the
second-order data. The engineer must inspect the lowest-order
non-vanishing term in the Taylor expansion, probe the function along
a spread of directions, or perturb the problem to break the
degeneracy. Reporting a singular-Hessian critical point as a minimum
or maximum without this further work is an error, not a shortcut.
\end{principle}

\subsection{What survives the collapse}

The above failures are about uniform distributions and isotropic
geometry. Real engineering data is rarely uniform; it sits on or
near a low-dimensional subspace, manifold, or cluster. The
\engterm{intrinsic dimensionality} of the data, the dimensionality
of the manifold it actually occupies, is often much smaller than
the ambient dimensionality of the parameter space. Volume II
Chapters 9 and 10 introduced the singular-value decomposition and
principal component analysis as tools that find this intrinsic
structure; the geometric content is that an SVD gives the engineer
a working set of effective dimensions that resist the failures
above.

The working principle: high-dimensional intuition fails in the
ambient space, but engineering data usually sits on a much smaller
intrinsic manifold. The job of the engineer is to find that
manifold, work on it, and not pretend that an algorithm tuned in
two dimensions extends to a thousand without rethinking. The
\engterm{scaling} archetype meets multivariable calculus here: the
geometry of the parameter space scales with $n$ in ways the working
intuition does not anticipate.

\begin{principle}[Cost of dimensionality]
For an algorithm whose statistical or computational cost grows
exponentially in the ambient dimension $n$, the working strategy is
not to attack the cost in $n$ but to find the intrinsic dimension
$d \ll n$ on which the problem actually lives. Volume II Chapter
9-10 (SVD, PCA), Volume II Chapter 14 (regularisation), and
Volume VII (representation learning) develop the working tools.
This chapter supplies the geometric warning.
\end{principle}

\section{Project}\pathtag{core}

\begin{project}[A three-constraint optimisation by hand]

\textbf{Track.} Analysis only. \textbf{Hazard class.} None.
\textbf{Tools.} Paper, pencil, calculator.

\textbf{The problem.} Write out, by hand, the full Lagrangian and
KKT analysis of an optimisation problem with three constraints.
The reader chooses the problem from the list below or proposes an
equivalent of comparable structure.

\textbf{Suggested problems.}
\begin{enumerate}[leftmargin=*]
  \item Minimise the surface area of a closed cylindrical can
    subject to (i) a fixed volume $V_{0}$, (ii) a height-to-radius
    ratio bounded above $h/r \leq 4$, (iii) a height-to-radius
    ratio bounded below $h/r \geq 1$. Two of the three constraints
    are inequalities; one is an equality.
  \item Minimise the power consumption of a pump subject to (i) a
    fixed throughput $Q_{0}$, (ii) a budget on the pump cost
    $C \leq C_{0}$ (the cost is a known function of the pump's
    operating point), (iii) a maximum operating temperature
    $T \leq T_{0}$ (also a known function). One equality, two
    inequalities.
  \item Fit a quadratic $f(x) = a x^{2} + b x + c$ through three
    given points $(x_{i}, y_{i})$, $i = 1, 2, 3$, while also
    requiring $f'(x_{1}) = m_{1}$ for a given slope $m_{1}$. The
    fit minimises the sum of squared residuals at the three points
    subject to the four interpolation/derivative constraints; in
    practice one of the residuals will be zero at the optimum and
    the analysis identifies which.
  \item Minimise a quadratic cost $f(\vect{x}) = \vect{x}^{\mathsf{T}}
    Q\vect{x}$ for a given positive-definite $Q$ subject to (i) a
    linear equality $A_{1}\vect{x} = b_{1}$, (ii) a linear equality
    $A_{2}\vect{x} = b_{2}$, (iii) a linear inequality $\vect{c}^{
    \mathsf{T}}\vect{x} \leq d$. The reader chooses dimensions $n$
    in which the algebra is tractable (suggested $n = 3$ or $4$).
\end{enumerate}

\textbf{The deliverable.} A 5-to-8-page hand-derived solution and a
1500-word reflection. The solution must contain
\begin{enumerate}[leftmargin=*]
  \item A clear statement of $f$, $g_{j}$, $h_{\ell}$, the feasible
    set, and the unknowns.
  \item The Lagrangian written out with one multiplier per
    constraint, distinguishing $\lambda_{j}$ for equalities from
    $\mu_{\ell} \geq 0$ for inequalities.
  \item All four KKT conditions written explicitly: stationarity,
    primal feasibility, dual feasibility, complementary slackness.
  \item Case analysis on which inequality constraints are active.
    For each case, the reader solves the resulting equality-only
    system (Lagrange multipliers) and checks the inactive
    inequalities for primal feasibility and the active multipliers
    for dual feasibility. A complete analysis enumerates all
    $2^{m}$ cases for $m$ inequalities; with two inequalities there
    are four cases, which is tractable.
  \item Verification of the second-order sufficient conditions where
    they apply. The reader writes the Hessian of the Lagrangian
    restricted to the tangent space of the active constraints and
    checks its definiteness.
  \item The numerical values of the optimal $\vect{x}^{*}$ and all
    multipliers, with units carried through.
\end{enumerate}

\textbf{The reflection.} The 1500-word reflection addresses
\begin{itemize}[leftmargin=*]
  \item Which inequality constraints were active at the optimum,
    and which were inactive. What does that say about the
    engineering of the problem? An inactive constraint is, in
    retrospect, a piece of slack in the design; an active constraint
    is the binding limit.
  \item What the multiplier values mean operationally. For each
    active constraint, the multiplier reports the marginal
    sensitivity of the optimum to a relaxation of that constraint.
    Translate this back into the units of the design problem.
  \item Which case in the case analysis would be hardest to find
    by gradient descent without the KKT structure. The reader's
    answer should connect the analytical apparatus of this chapter
    to the algorithmic apparatus of Volume II Chapter 15.
  \item One source of error the reader noticed in their own working,
    and one heuristic the reader has acquired by completing the
    project.
\end{itemize}

\textbf{Phases.}
\begin{enumerate}[leftmargin=*]
  \item \textit{Setup} (30--60 minutes). Choose the problem.
    State $f$, the equality $g$, the two inequalities $h_{1},
    h_{2}$, and the unknowns. Sketch the feasible region in two
    variables when possible (eliminate the rest by substitution).
  \item \textit{Lagrangian and KKT} (60--90 minutes). Write the
    Lagrangian. Write the four KKT conditions explicitly. Identify
    the unknowns: $\vect{x}^{*}, \lambda^{*}, \mu_{1}^{*},
    \mu_{2}^{*}$.
  \item \textit{Case analysis} (3--5 hours). Enumerate the four
    cases for two inequalities (both active, only $h_{1}$ active,
    only $h_{2}$ active, neither active). For each, solve the
    equality-only Lagrangian system, then check the inactive
    inequalities for primal feasibility and the active multipliers
    for dual feasibility. Discard cases that fail either check.
  \item \textit{Second-order check and reporting} (60--90 minutes).
    For the surviving case (or cases), verify the second-order
    sufficient condition by the bordered-Hessian test or by
    inspection. Report $\vect{x}^{*}$ and all four multipliers,
    with units.
  \item \textit{Reflection} (3--4 hours of writing).
\end{enumerate}

\textbf{Success criteria.} The deliverable counts as complete when
(i) every KKT condition is written explicitly and verified, (ii)
the case analysis enumerates and resolves all four cases, (iii)
the second-order check identifies whether the candidate is a
constrained minimum, maximum, or saddle, (iv) the multipliers have
the correct sign (non-negative for inequalities) and a unit-bearing
interpretation, and (v) the reflection makes explicit which
constraints were active, which inactive, and what the active
multipliers' magnitudes mean for the engineering of the problem.

\textbf{Common failure modes.}
\begin{itemize}[leftmargin=*]
  \item \textit{Missing a case.} The most frequent error is to
    write the KKT conditions for one assumed active set and to
    skip the other three cases. With two inequalities there are
    four cases; with three there are eight.
  \item \textit{Wrong sign on $\mu$.} A negative inequality
    multiplier signals that the assumed active set is wrong, not
    that the analysis has produced an answer. Drop the candidate
    and try a different active set.
  \item \textit{Skipping primal feasibility.} The case where only
    $h_{1}$ is active still requires checking that $h_{2} \leq 0$
    at the candidate $\vect{x}^{*}$. A candidate that fails this
    check is infeasible and is not the optimum.
  \item \textit{Confusing the bordered Hessian with the Hessian of
    $f$.} The second-order sufficient condition uses the Hessian
    of the Lagrangian, $H_{f} - \lambda^{*} H_{g}$, not the Hessian
    of $f$ alone. Forgetting the curvature contribution of the
    constraint is a standard slip.
  \item \textit{Reporting units inconsistently.} The multiplier
    $\lambda$ has units of $[f] / [g]$ (cost per unit of the
    constraint quantity). A multiplier reported as a dimensionless
    number is missing a load-bearing piece of the answer.
\end{itemize}

\textbf{Time.} Six to ten hours of derivation, plus three to four
hours of writing. Total approximately ten to fourteen hours,
distributed across two or three sessions.

\textbf{Solutions.} A reference set of three-constraint
optimisations from the general editor, one per suggested problem.

\end{project}

\section{Exercises}

The exercises train calculation, derivation, estimation, simulation,
design, diagnosis, failure analysis, and judgment. Calculation and
derivation problems have full solutions; open-ended problems have
rubrics. Exercises are tagged with their reader-path tier.

\subsection*{Calculation}\pathtag{core}

\begin{exercise}
Compute $\grad f$ for $f(x, y) = x^{3} y - x y^{3}$. Identify all
critical points, classify each by the second-derivative test, and
report the Hessian eigenvalues at each.

\paragraph{Solution sketch.} $\grad f = (3x^{2} y - y^{3},
x^{3} - 3 x y^{2})^{\mathsf{T}} = (y(3x^{2} - y^{2}),
x(x^{2} - 3 y^{2}))^{\mathsf{T}}$. Setting both components to zero:
either $y = 0$, in which case $x \cdot x^{2} = 0$ forces $x = 0$;
or $3 x^{2} = y^{2}$ and $x^{2} = 3 y^{2}$, which combine to give
$3 x^{2} \cdot 3 = y^{2} \cdot 3 = x^{2}$, i.e.\ $9 x^{2} = x^{2}$,
forcing $x = 0$, hence $y = 0$. The only critical point is the
origin. The Hessian is $H = \begin{pmatrix} 6 x y & 3(x^{2} - y^{2})
\\ 3(x^{2} - y^{2}) & -6 x y \end{pmatrix}$, which at the origin
is the zero matrix. Both eigenvalues are zero; the second-derivative
test is inconclusive. Inspection of $f$ along the directions
$(1, 1)$ and $(1, -1)$ gives $f(t, t) = t^{3} \cdot t - t \cdot
t^{3} = 0$ and $f(t, -t) = -t^{4} + t^{4} = 0$; along $(1, 0)$,
$f = 0$; along $(\cos\alpha, \sin\alpha)$ with general $\alpha$,
$f$ takes both signs in any neighborhood of the origin, so the
origin is a degenerate saddle. The reader recognises $f = x y(x^{2}
- y^{2})$ as the real part of $z^{4}/4$ for $z = x + i y$, whose
contour structure has eight rays of zero through the origin.
\end{exercise}

\begin{exercise}
Compute $\grad f$ for $f(x, y, z) = e^{x y z}$ and evaluate
$\grad f$ at $(1, 2, 3)$. Compute the directional derivative at
this point in the direction $\unitvec{u} = (1, 1, 1)/\sqrt{3}$.

\paragraph{Solution sketch.} By the chain rule on $e^{u}$ with
$u = x y z$,
\[
\grad f = e^{x y z}\,(y z, x z, x y)^{\mathsf{T}}.
\]
At $(1, 2, 3)$, $x y z = 6$, so $\grad f = e^{6}(6, 3, 2)^{
\mathsf{T}}$, numerically $e^{6} \approx 403.43$, giving
$\grad f \approx (2420.6, 1210.3, 806.9)^{\mathsf{T}}$. The
directional derivative is
\[
D_{\unitvec{u}} f = \grad f \cdot \unitvec{u} = \frac{e^{6}}{\sqrt{3}}
(6 + 3 + 2) = \frac{11\,e^{6}}{\sqrt{3}} \approx 2562.7.
\]
The reader observes that the directional derivative along the
diagonal is less than the gradient norm $e^{6}\sqrt{49} = 7 e^{6}
\approx 2824$, consistent with $D_{\unitvec{u}} f \leq \norm{\grad
f}$ from Cauchy-Schwarz with equality only along $\grad f$ itself.
\end{exercise}

\begin{exercise}
Compute $\iint_{D} (x^{2} + y^{2}) \,\dd A$ where $D$ is the unit
disk $x^{2} + y^{2} \leq 1$, by switching to polar coordinates.
Verify the result by also evaluating the integral as an iterated
Cartesian integral.

\paragraph{Solution sketch.} In polar, $x^{2} + y^{2} = r^{2}$ and
$\dd A = r\,\dd r\,\dd\theta$, so
\[
\iint_{D}(x^{2} + y^{2})\,\dd A = \int_{0}^{2\pi}\!\int_{0}^{1}
r^{2} \cdot r \,\dd r\,\dd\theta = 2\pi \cdot \tfrac{1}{4} =
\frac{\pi}{2}.
\]
In Cartesian, on the disk $y \in [-\sqrt{1 - x^{2}}, \sqrt{1 -
x^{2}}]$,
\[
\int_{-1}^{1}\!\int_{-\sqrt{1 - x^{2}}}^{\sqrt{1 - x^{2}}}
(x^{2} + y^{2})\,\dd y\,\dd x =
\int_{-1}^{1}\!\left(2 x^{2}\sqrt{1 - x^{2}} +
\tfrac{2}{3}(1 - x^{2})^{3/2}\right)\dd x.
\]
The substitution $x = \sin\theta$, $\dd x = \cos\theta\,\dd\theta$
turns each integrand into a standard $\cos^{n}\theta$ form; both
integrate to $\pi/4$ and $\pi/4$ respectively (using Wallis or by
direct $\cos^{2}\sin^{2}$ and $\cos^{4}$ reductions), summing to
$\pi/2$. The two answers agree. The Cartesian route is correct but
laborious; the polar route is the working choice.
\end{exercise}

\begin{exercise}
Compute $\iiint_{V} z \,\dd V$ where $V$ is the upper hemisphere
$x^{2} + y^{2} + z^{2} \leq R^{2}$, $z \geq 0$. Use spherical
coordinates and express the answer in terms of $R$.

\paragraph{Solution sketch.} In spherical, $z = r\cos\theta$,
$\dd V = r^{2}\sin\theta\,\dd r\,\dd\theta\,\dd\phi$, with bounds
$r \in [0, R]$, $\theta \in [0, \pi/2]$, $\phi \in [0, 2\pi]$.
\[
\iiint_{V} z\,\dd V = \int_{0}^{2\pi}\!\int_{0}^{\pi/2}\!\int_{0}^{R}
(r\cos\theta)\,r^{2}\sin\theta\,\dd r\,\dd\theta\,\dd\phi.
\]
The $\phi$-integral gives $2\pi$, the $r$-integral $R^{4}/4$, and
the $\theta$-integral $\int_{0}^{\pi/2}\sin\theta\cos\theta\,
\dd\theta = 1/2$. The product is
\[
\iiint_{V} z\,\dd V = 2\pi \cdot \frac{R^{4}}{4} \cdot \frac{1}{2}
= \frac{\pi R^{4}}{4}.
\]
The reader recognises $\iiint_{V} z\,\dd V / V_{\text{hemi}}$ as
the $z$-coordinate of the centroid: $(\pi R^{4}/4) /
(2\pi R^{3}/3) = 3 R / 8$, the standard formula.
\end{exercise}

\begin{exercise}
Compute the Jacobian determinant of the map $\vect{T}(u, v) = (u +
v, u - v)$ from $\R^{2} \to \R^{2}$. Use it to evaluate
$\iint_{D'} (x + y)\,\dd x\,\dd y$ where $D'$ is the square with
vertices $(0, 0), (1, 0), (1, 1), (0, 1)$, by transforming to a
region in the $(u, v)$-plane.

\paragraph{Solution sketch.} The Jacobian matrix is $J = \begin{pmatrix} 1 & 1 \\ 1 & -1 \end{pmatrix}$ with $\det J = -2$, so
$\abs{\det J} = 2$. The map sends $(u, v) \to (x, y) = (u + v,
u - v)$, with inverse $u = (x + y)/2$, $v = (x - y)/2$. The square
$D'$ in the $(x, y)$-plane has vertices that map back to
$(0, 0), (1/2, 1/2), (1, 0), (1/2, -1/2)$ in the $(u, v)$-plane,
which is itself a square rotated by $45^{\circ}$ with side
$1/\sqrt{2}$. The integrand is $x + y = 2 u$, so
\[
\iint_{D'}(x + y)\,\dd x\,\dd y = \iint_{D'_{uv}} 2 u \cdot 2\,
\dd u\,\dd v = 4 \iint_{D'_{uv}} u\,\dd u\,\dd v.
\]
The $u$-integral over the rotated square evaluates to (by direct
iterated integration) $1/4$, giving the result $1$. The reader can
verify by computing the original integral directly: $\iint_{D'}
(x + y)\,\dd x\,\dd y = \int_{0}^{1}\!\int_{0}^{1}(x + y)\,\dd y\,
\dd x = \int_{0}^{1}(x + 1/2)\,\dd x = 1/2 + 1/2 = 1$. The two
agree. The change of variables is unnecessary here; the exercise
trains the apparatus on a problem with a known answer.
\end{exercise}

\begin{exercise}
Locate and classify the critical points of $f(x, y) = x^{4} + y^{4}
- 4 x y$. Compute the Hessian at each and apply the second-
derivative test.

\paragraph{Solution sketch.} $\grad f = (4 x^{3} - 4 y, 4 y^{3} -
4 x)^{\mathsf{T}}$. Setting both to zero: $y = x^{3}$ and $x = y^{3}
= x^{9}$, so $x(x^{8} - 1) = 0$. Real roots: $x \in \set{0, 1, -1}$.
The three critical points are $(0, 0)$, $(1, 1)$, $(-1, -1)$. The
Hessian is $H = \begin{pmatrix} 12 x^{2} & -4 \\ -4 & 12 y^{2}
\end{pmatrix}$.
\begin{itemize}[leftmargin=*]
  \item At $(0, 0)$: $H = \begin{pmatrix} 0 & -4 \\ -4 & 0
    \end{pmatrix}$, $\det H = -16 < 0$, saddle. Eigenvalues
    $\pm 4$.
  \item At $(1, 1)$: $H = \begin{pmatrix} 12 & -4 \\ -4 & 12
    \end{pmatrix}$, $\det H = 144 - 16 = 128 > 0$, with $A = 12 >
    0$: local minimum. Eigenvalues $12 \pm 4 = 16, 8$.
  \item At $(-1, -1)$: same as $(1, 1)$ by even symmetry: local
    minimum. Eigenvalues $16, 8$.
\end{itemize}
The function has two equal local minima and a central saddle: the
canonical "double-well" shape that recurs in phase-transition and
bistable-system models (Volume IV).
\end{exercise}

\subsection*{Derivation}\pathtag{standard}

\begin{exercise}
Derive the multivariable chain rule $\dv{}{t} f(\vect{x}(t)) =
\grad f \cdot \dv{\vect{x}}{t}$ from the limit definition of the
partial derivative, stating the smoothness assumptions on $f$ and
$\vect{x}(t)$.

\paragraph{Solution sketch.} Assume $f : \R^{n} \to \R$ is
differentiable on a region containing the image of $\vect{x}(t)$,
and $\vect{x}(t) : I \to \R^{n}$ is differentiable on the open
interval $I$. By differentiability of $f$ at $\vect{x}(t)$,
\[
f(\vect{x}(t) + \dd\vect{x}) - f(\vect{x}(t)) = \grad f(\vect{x}(t))
\cdot \dd\vect{x} + \varepsilon(\dd\vect{x})\norm{\dd\vect{x}},
\]
where $\varepsilon(\dd\vect{x}) \to 0$ as $\dd\vect{x} \to 0$. Set
$\dd\vect{x} = \vect{x}(t + h) - \vect{x}(t)$ for small $h$.
Dividing by $h$,
\[
\frac{f(\vect{x}(t + h)) - f(\vect{x}(t))}{h} = \grad f(\vect{x}(t))
\cdot \frac{\vect{x}(t + h) - \vect{x}(t)}{h} +
\varepsilon\,\frac{\norm{\vect{x}(t + h) - \vect{x}(t)}}{h}.
\]
The first term tends to $\grad f \cdot \dv{\vect{x}}{t}$ by
differentiability of $\vect{x}(t)$ at $t$. The second term has
$\varepsilon \to 0$ and bounded factor $\norm{\dv{\vect{x}}{t}}$ in
the limit, so its product tends to zero. Taking the limit gives
the chain rule. The smoothness hypotheses are: $f$ differentiable
(equivalently, continuous first partials) on a neighbourhood of
$\vect{x}(t)$; $\vect{x}(t)$ differentiable on $I$.
\end{exercise}

\begin{exercise}
Derive the steepest-ascent theorem from the relation $D_{\unitvec{u}}
f = \grad f \cdot \unitvec{u}$, identifying the direction
$\unitvec{u}$ that maximises $D_{\unitvec{u}} f$ subject to
$\norm{\unitvec{u}} = 1$. Use the Cauchy-Schwarz inequality.

\paragraph{Solution sketch.} The Cauchy-Schwarz inequality states
$\abs{\vect{a} \cdot \vect{b}} \leq \norm{\vect{a}}\,\norm{\vect{b}}$,
with equality if and only if $\vect{a}$ and $\vect{b}$ are
parallel. Applied with $\vect{a} = \grad f(\vect{x})$ and $\vect{b}
= \unitvec{u}$ (so $\norm{\unitvec{u}} = 1$),
\[
\abs{D_{\unitvec{u}} f(\vect{x})} = \abs{\grad f \cdot \unitvec{u}}
\leq \norm{\grad f} \cdot 1 = \norm{\grad f}.
\]
Equality holds when $\unitvec{u}$ is parallel to $\grad f$. Among
the two unit vectors $\pm \grad f / \norm{\grad f}$, the choice
that maximises $D_{\unitvec{u}} f$ (rather than its absolute value)
is $\unitvec{u}^{*} = \grad f / \norm{\grad f}$, giving $D_{
\unitvec{u}^{*}} f = \norm{\grad f}$. The opposite choice minimises
$D_{\unitvec{u}} f$ at $-\norm{\grad f}$; this is the steepest-
descent direction.
\end{exercise}

\begin{exercise}
Derive the Lagrange-multiplier condition $\grad f = \lambda \grad g$
from the geometric requirement that, at a constrained extremum, the
gradient of $f$ has no component tangent to the constraint surface
$g(\vect{x}) = 0$. State the smoothness assumptions on $g$.

\paragraph{Solution sketch.} Assume $g : \R^{n} \to \R$ is
continuously differentiable on a neighbourhood of $\vect{x}^{*}$
with $\grad g(\vect{x}^{*}) \neq \vect{0}$. The implicit-function
theorem gives that the level set $g = 0$ is locally a smooth
$(n - 1)$-dimensional manifold $M$, with tangent space at
$\vect{x}^{*}$ equal to the orthogonal complement of $\grad g(
\vect{x}^{*})$: $T_{\vect{x}^{*}} M = \set{\vect{d} : \grad g \cdot
\vect{d} = 0}$. At a constrained extremum of $f$ on $M$, the
directional derivative of $f$ vanishes along every tangent
direction: $\grad f \cdot \vect{d} = 0$ for all $\vect{d} \in
T_{\vect{x}^{*}} M$. Geometrically, $\grad f$ is orthogonal to
$T_{\vect{x}^{*}} M$; since $T_{\vect{x}^{*}} M$ is the orthogonal
complement of the span of $\grad g$, $\grad f$ lies in that span:
$\grad f = \lambda\,\grad g$ for some scalar $\lambda$. The
hypothesis $\grad g \neq \vect{0}$ is the constraint qualification
in this case; it rules out degenerate constraints whose level set
has a singular point at $\vect{x}^{*}$.
\end{exercise}

\begin{exercise}
Derive the Jacobian determinant $r$ for the polar-coordinate
transformation $\vect{T}(r, \theta) = (r\cos\theta, r\sin\theta)$
from first principles, by computing the area of the image of an
infinitesimal rectangle $[r, r + \dd r] \times [\theta, \theta +
\dd\theta]$ in the $(r, \theta)$-plane.

\paragraph{Solution sketch.} The image of the small rectangle is
a wedge (\cref{fig:vol02:ch11:polar-area-element}). Its four
corners are
\[
\vect{T}(r, \theta),\ \vect{T}(r + \dd r, \theta),\
\vect{T}(r, \theta + \dd\theta),\ \vect{T}(r + \dd r, \theta + \dd\theta).
\]
The two edge vectors at $\vect{T}(r, \theta)$ are, to first order,
\[
\partial_{r}\vect{T}\,\dd r = (\cos\theta, \sin\theta)\,\dd r,
\qquad
\partial_{\theta}\vect{T}\,\dd\theta = (-r\sin\theta, r\cos\theta)\,
\dd\theta.
\]
The area of the parallelogram they span is the absolute value of
their $2 \times 2$ cross-product determinant,
\[
\dd A = \abs{\det \begin{pmatrix}
\cos\theta\,\dd r & -r\sin\theta\,\dd\theta \\
\sin\theta\,\dd r & r\cos\theta\,\dd\theta
\end{pmatrix}} = \abs{r\cos^{2}\theta + r\sin^{2}\theta}\,\dd r\,
\dd\theta = r\,\dd r\,\dd\theta.
\]
The factor $r$ is the absolute value of the Jacobian determinant,
matching the change-of-variables formula. Geometrically, the polar
wedge's "long side" is the arc $r\,\dd\theta$ and the "short side"
is the radial increment $\dd r$, so the rectangle's area
$r\,\dd\theta \cdot \dd r$ recovers the factor directly.
\end{exercise}

\begin{exercise}
Derive the second-order Taylor expansion of $f : \R^{n} \to \R$
around a point $\vect{a}$, and express the second-order term as
$\tfrac{1}{2}\,(\vect{x} - \vect{a})^{\mathsf{T}} H(\vect{a})\,
(\vect{x} - \vect{a})$. State and justify the assumption that
$f$ has continuous second partials.

\paragraph{Solution sketch.} Fix $\vect{a}$ and a target point
$\vect{x} = \vect{a} + \vect{h}$. Define the one-variable
restriction $g(t) = f(\vect{a} + t\vect{h})$ for $t \in [0, 1]$.
The one-variable Taylor theorem with integral remainder, applied to
$g$, gives
\[
g(1) = g(0) + g'(0) + \tfrac{1}{2} g''(0) + R_{2},
\]
where $R_{2} = \int_{0}^{1}(1 - t)\,(g''(t) - g''(0))\,\dd t \to 0$
as $\vect{h} \to 0$ in the sense $R_{2} = \oom{\norm{\vect{h}}^{2}}$.
By the multivariable chain rule of section 11.2.5,
\[
g'(t) = \grad f(\vect{a} + t\vect{h}) \cdot \vect{h},
\qquad
g''(t) = \vect{h}^{\mathsf{T}}\,H(\vect{a} + t\vect{h})\,\vect{h}.
\]
The second identity follows by applying the chain rule a second
time to each component $\partial_{i} f$. Setting $t = 0$ gives
$g'(0) = \grad f(\vect{a}) \cdot \vect{h}$ and $g''(0) = \vect{h}^{
\mathsf{T}} H(\vect{a})\,\vect{h}$. Substituting and writing
$\vect{h} = \vect{x} - \vect{a}$,
\[
f(\vect{x}) = f(\vect{a}) + \grad f(\vect{a}) \cdot (\vect{x} -
\vect{a}) + \tfrac{1}{2}(\vect{x} - \vect{a})^{\mathsf{T}}
H(\vect{a})\,(\vect{x} - \vect{a}) + \oom{\norm{\vect{x} - \vect{a}}
^{2}}.
\]
The assumption of continuous second partials is required at two
points. First, it lets the chain rule pass through twice, so $g''(t)$
exists and equals the stated quadratic form. Second, continuity of
$H$ on a neighbourhood of $\vect{a}$ gives $H(\vect{a} + t\vect{h})
\to H(\vect{a})$ uniformly in $t \in [0, 1]$ as $\vect{h} \to 0$,
which is what makes the integral remainder $R_{2}$ small of order
$\norm{\vect{h}}^{2}$ rather than just $\norm{\vect{h}}$. Without
continuity of $H$, the Hessian at $\vect{a}$ is still well-defined
as a matrix of mixed partials (by Clairaut, symmetric), but the
remainder estimate weakens, and the second-derivative test of
section 11.3 loses its quantitative grip. The hypothesis is mild;
every function the engineer meets in this volume satisfies it.
\end{exercise}

\subsection*{Estimation}\pathtag{core}

\begin{exercise}
Estimate the optimal aspect ratio $h/r$ of a closed cylindrical
soda can of fixed volume $V_{0} = 0.355\,\si{\liter}$ minimising
surface area, then compare to the actual aspect ratio of a real
soda can. Identify a reason for the discrepancy. (Hint: the lid and
side may have different material costs.)

\paragraph{Solution sketch.} The Lagrange analysis of section
11.4.4 gives $h^{*} = 2 r^{*}$ at the surface-area optimum. With
$V_{0} = 355\,\si{\milli\liter} = 355\,\si{\centi\meter\cubed}$,
the volume constraint $\pi r^{2} h = V_{0}$ with $h = 2 r$ gives
$2\pi r^{3} = 355$, so $r^{*} = (355 / (2\pi))^{1/3} \approx
3.84\,\si{\centi\meter}$ and $h^{*} \approx 7.67\,\si{\centi\meter}$.
The optimal $h/r$ is exactly $2$. A real 355\,mL beverage can has
$r \approx 3.25\,\si{\centi\meter}$, $h \approx
12.2\,\si{\centi\meter}$, $h/r \approx 3.75$, almost twice the
"optimal" ratio. The discrepancy is real and engineering-driven.
A weighted surface-area objective $S = \alpha (2\pi r^{2}) + \beta
(2\pi r h)$ with the lid weight $\alpha > \beta$ (the lid and bottom
use thicker, more expensive aluminium because the lid carries the
pressure load and the pull-tab feature) gives, on re-doing the
Lagrange calculation, $h^{*}/r^{*} = 2\alpha/\beta$. A typical
$\alpha/\beta \approx 1.9$ for double-thickness ends recovers
$h/r \approx 3.8$. Two further effects push the same direction.
Cans must fit the human hand (a wider can is harder to grip), and
manufacturing economies of scale on the can-body sheet metal favour
a tall narrow can per unit volume of stored aluminium. The
mathematically optimal can is not the manufactured optimal can; the
constraint set on a real can is richer than the textbook minimum.
\end{exercise}

\begin{exercise}
Estimate the fraction of the volume of an $n$-dimensional ball that
sits within $5\%$ of the surface, for $n = 3, 10, 100, 1000$.
Report each fraction to two significant figures, and identify the
dimension at which more than half the volume is in the outer shell.

\paragraph{Solution sketch.} The shell formula of section 11.8.1
with $\varepsilon = 0.05$ reads $V_{\text{shell}}/V = 1 - (1 -
\varepsilon)^{n} = 1 - 0.95^{n}$.
\begin{itemize}[leftmargin=*]
  \item $n = 3$: $1 - 0.95^{3} = 1 - 0.857 \approx 0.14$, about
    $14\%$.
  \item $n = 10$: $1 - 0.95^{10} = 1 - 0.599 \approx 0.40$, about
    $40\%$.
  \item $n = 100$: $1 - 0.95^{100} = 1 - 0.0059 \approx 0.99$,
    about $99\%$.
  \item $n = 1000$: $1 - 0.95^{1000} \approx 1 - 5.3 \times 10^{
    -23} \approx 1.00$ to many decimal places.
\end{itemize}
The half-volume threshold $0.95^{n} = 0.5$ gives $n = \log 0.5 /
\log 0.95 \approx 13.5$, so $n = 14$ is the dimension at which more
than half the volume of the unit ball is within $5\%$ of its
surface. The engineering reading: a Monte Carlo sampler that
assumes most of its samples land in the "interior" of a fifteen-
dimensional ball is mistaken. The shell fraction crosses $50\%$
before the dimension reaches the number of fingers on two hands and
a foot, and by $n = 100$ the interior is empty for practical
purposes. The same shell-concentration phenomenon underlies the
\engterm{soap-bubble} intuition for high-dimensional Gaussian
distributions, where almost all the probability mass sits in a thin
shell at radius $\sqrt{n}$ rather than near the mean.
\end{exercise}

\begin{exercise}
Estimate the rate of heat loss per unit area through a wall whose
temperature gradient is $20\,\si{\degreeCelsius}$ across
$0.20\,\si{\meter}$ of wood with thermal conductivity $k \approx
0.15\,\si{\watt\per\meter\per\kelvin}$. State the magnitude and
direction of $\grad T$, and use $\vect{q} = -k \grad T$.

\paragraph{Solution sketch.} Take the wall to span $x \in [0,
0.20]\,\si{\meter}$ with $T(0) = 20\,\degC$ on the warm side and
$T(0.20) = 0\,\degC$ on the cold side. Inside a steady-state,
one-dimensional, homogeneous slab, the temperature is linear, so
\[
\grad T = \pdv{T}{x}\,\unitvec{x} =
\frac{0 - 20}{0.20}\,\unitvec{x} = -100\,\si{\kelvin\per\meter}\,
\unitvec{x}.
\]
The magnitude is $\norm{\grad T} = 100\,\si{\kelvin\per\meter}$; the
direction points from cold to hot, that is, in the $-\unitvec{x}$
direction, because temperature increases as $x$ decreases. The
Fourier heat-flux law $\vect{q} = -k\,\grad T$ then gives
\[
\vect{q} = -(0.15)(-100)\,\unitvec{x} = +15\,\si{\watt\per\meter
\squared}\,\unitvec{x},
\]
a heat flux of $15\,\si{\watt\per\meter\squared}$ in the
$+\unitvec{x}$ direction, that is, from warm side to cold side, as
the second law requires. For a $4 \times 2.5\,\si{\meter}$ wall
(a small bedroom external wall), the total heat-loss rate is about
$150\,\si{\watt}$, comparable to a single incandescent bulb.
Standard wall constructions reduce this by adding insulation; a
$0.10\,\si{\meter}$ layer of glass-wool insulation ($k \approx
0.04\,\si{\watt\per\meter\per\kelvin}$) in series with the wood
gives total thermal resistance about $R = 0.20/0.15 + 0.10/0.04 =
3.83\,\si{\meter\squared\kelvin\per\watt}$, cutting the heat flux
to about $5.2\,\si{\watt\per\meter\squared}$, a factor-of-three
reduction.
\end{exercise}

\begin{exercise}
Estimate the volume of a typical chicken egg by setting up a
spherical-coordinate integral over an ellipsoidal approximation. Use
$a = 30\,\si{\milli\meter}$, $b = 22\,\si{\milli\meter}$,
$c = 22\,\si{\milli\meter}$ as the three semi-axes. Compare to a
typical reported egg volume of about $50\,\si{\milli\liter}$.

\paragraph{Solution sketch.} The ellipsoid $x^{2}/a^{2} + y^{2}/b^{
2} + z^{2}/c^{2} \leq 1$ has, by the change of variables $x = a u,
y = b v, z = c w$ with Jacobian determinant $abc$, image the unit
ball in $(u, v, w)$. So
\[
V = \iiint_{\text{ellipsoid}} \dd V = abc \iiint_{\text{unit ball}}
\dd u\,\dd v\,\dd w = abc \cdot \frac{4\pi}{3}.
\]
The unit-ball volume is $4\pi/3$ (worked example of section 11.6.4
with $R = 1$); the ellipsoid is the unit ball stretched by $a, b,
c$ along the three axes. Plugging in,
\[
V = \frac{4\pi}{3} \cdot 30 \cdot 22 \cdot 22\,\si{\milli\meter
\cubed} = \frac{4\pi}{3} \cdot 14520 \approx 60800\,\si{\milli\meter
\cubed} \approx 60.8\,\si{\milli\liter}.
\]
The estimate exceeds the reported $50\,\si{\milli\liter}$ by about
$20\%$. The discrepancy has two engineering explanations. First, a
real chicken egg is not an ellipsoid; it is an asymmetric prolate
shape with a wider blunt end and a narrower pointed end, better
approximated by a Hugelschaffer or "egg-curve" rotation profile. An
ellipsoidal envelope overestimates volume. Second, the air cell at
the blunt end (about $5$ to $10\,\si{\milli\liter}$ in a fresh egg,
growing with age) is empty space inside the shell; the reported
"egg volume" sometimes refers to liquid contents excluding the air
cell. A better engineering approximation uses an egg-shaped
parametric surface; a more honest one accepts the $20\%$ envelope
overestimate as the cost of using the simpler ellipsoidal model.
The pedagogical content of the exercise is the change-of-variables
trick to ellipsoidal coordinates, not the high-fidelity geometry of
a chicken egg.
\end{exercise}

\begin{exercise}
Estimate the multiplier $\lambda$ in a budget-constrained
production-optimisation problem in which the unconstrained optimum
profit is $\$120{,}000$ at a cost of $\$80{,}000$, and the
constrained optimum (with budget cap $\$70{,}000$) is $\$110{,}000$
at a cost of exactly $\$70{,}000$. Interpret $\lambda$ in
\$ per \$.

\paragraph{Solution sketch.} Let $f^{*}(c)$ be the maximum profit
attainable when the cost is constrained to at most $c$. Two data
points sit on the curve: $f^{*}(80{,}000) = 120{,}000$ at the
unconstrained optimum (the constraint is inactive at and above
$\$80{,}000$), and $f^{*}(70{,}000) = 110{,}000$ at the active
constraint. The shadow-price interpretation of section 11.4.5 gives
$\dv{f^{*}}{c}\big|_{c = c^{*}} = \lambda^{*}$. A finite-difference
estimate over the interval $[70{,}000, 80{,}000]$:
\[
\lambda \approx \frac{\Delta f^{*}}{\Delta c} =
\frac{120{,}000 - 110{,}000}{80{,}000 - 70{,}000} =
\frac{10{,}000}{10{,}000} = 1.0\,\$/\$.
\]
The multiplier reads as $\$1.00$ of additional profit per $\$1.00$
of relaxed budget over this range. Three engineering readings.
First, $\lambda = 1$ at this scale means the budget constraint is
binding hard; the project would absorb an additional dollar of
budget at full marginal value. Second, the linear extrapolation
holds only locally; if the unconstrained optimum is at $\$80{,}000$,
then $\lambda$ falls to zero (the multiplier is the marginal value
of a dollar of slack constraint, and slack is worth nothing) as
$c$ rises past $\$80{,}000$. The finite-difference shadow price is
an average over the interval, not the value at either endpoint. A
local Taylor expansion of $f^{*}(c)$ near $c = 70{,}000$ would
typically give a higher $\lambda$ than $1.0$; near $c = 80{,}000$,
a lower one. Third, the policy implication: if the agency considers
relaxing the cap from $\$70{,}000$ to $\$75{,}000$, the expected
gain in profit is approximately $\$5{,}000$, the integral of
$\lambda(c)$ from $70{,}000$ to $75{,}000$, well-approximated by
$\lambda \cdot \Delta c$ for a small relaxation.
\end{exercise}

\subsection*{Simulation}\pathtag{mastery}

\begin{exercise}
In a language of the reader's choice, generate $N = 10000$ random
points uniformly in the cube $[-1, 1]^{n}$ for $n = 2, 5, 10, 20,
50, 100$. For each $n$, compute the fraction of points that fall
inside the inscribed ball $\norm{\vect{x}} \leq 1$. Plot this
fraction against $n$ and compare against the analytical formula
$V_{n}(1) / 2^{n}$.

\paragraph{Solution sketch.} Drawing $N$ points uniformly in
$[-1, 1]^{n}$ and counting how many satisfy $\sum x_{i}^{2} \leq 1$
estimates the ratio $V_{n}(1)/2^{n}$ by Monte Carlo. The analytical
formula uses $V_{n}(1) = \pi^{n/2}/\Gamma(n/2 + 1)$, giving
\begin{itemize}[leftmargin=*]
  \item $n = 2$: $\pi/4 \approx 0.7854$. Expected count $\approx
    7854$.
  \item $n = 5$: $V_{5}(1) = 8\pi^{2}/15 \approx 5.264$; ratio
    $5.264/32 \approx 0.1645$. Expected count $\approx 1645$.
  \item $n = 10$: $V_{10}(1) = \pi^{5}/120 \approx 2.550$; ratio
    $2.550/1024 \approx 0.00249$. Expected count $\approx 25$.
  \item $n = 20$: ratio $\approx 2.5 \times 10^{-8}$. Expected
    count $\approx 0$ with $N = 10^{4}$; the rare event is below
    Monte-Carlo resolution.
  \item $n = 50, 100$: ratios $\approx 10^{-28}, 10^{-70}$. Zero
    hits expected; the inscribed ball is effectively empty under
    uniform cube sampling.
\end{itemize}
The student's plot should show empirical fractions that track the
analytical ratio for $n \leq 10$, then drop to zero for $n \geq 20$
because the rare event is below resolution. A standard-error band
$\pm \sqrt{p(1-p)/N}$ widens near zero, and the empirical zero is
consistent with the analytical $10^{-8}$. The qualitative content is
that the inscribed ball is a vanishing fraction of the cube even
at modest dimension, and uniform cube sampling does not put samples
in the ball. To estimate $V_{n}(1)/2^{n}$ at $n = 20$ or beyond, the
student must use importance sampling biased toward the origin, not
uniform cube sampling. The engineering content of this exercise is
the empirical confirmation that high-dimensional volume sits in the
corners of the cube and the analytical formula's prediction
matches.
\end{exercise}

\begin{exercise}
Implement gradient descent on the function $f(x, y) = (x - 1)^{2} +
10 (y - x^{2})^{2}$ (the Rosenbrock function with $a = 1, b = 10$),
starting from $(x_{0}, y_{0}) = (-1, 1)$, with a fixed step size
$\eta = 0.001$. Run for $10000$ iterations, plot the trajectory
on a contour plot of $f$, and report the final position. Identify
qualitatively why the trajectory is slow.

\paragraph{Solution sketch.} The gradient is
\[
\grad f = \begin{pmatrix} 2(x - 1) - 40 x (y - x^{2}) \\
20 (y - x^{2}) \end{pmatrix}.
\]
The minimum sits at $(x^{*}, y^{*}) = (1, 1)$ where both components
vanish. The Hessian at the minimum is
\[
H(1, 1) = \begin{pmatrix} 2 + 80 & -40 \\ -40 & 20 \end{pmatrix} =
\begin{pmatrix} 82 & -40 \\ -40 & 20 \end{pmatrix},
\]
with eigenvalues approximately $101.4$ and $0.59$, giving a
condition number $\kappa = \lambda_{\max}/\lambda_{\min} \approx
172$. Steepest-descent on a quadratic with condition number $\kappa$
converges with error reduction factor $((\kappa - 1)/(\kappa + 1))^
{2} \approx 0.977$ per iteration in the worst direction; the
number of iterations to reduce the error by a factor of ten is
about $\ln(10)/(1 - 0.977) \approx 100$. The reader's $10{,}000$
iterations should reach the curved valley quickly (within
$\sim 50$ iterations) and then crawl down the valley floor toward
$(1, 1)$. Empirically with $\eta = 0.001$ and start $(-1, 1)$,
the iterate lands near $(0.7, 0.5)$ after $10{,}000$ steps, well
short of the minimum but with $\norm{\grad f} \approx 0.5$, small
in absolute terms but stuck in the badly-conditioned valley.

The qualitative reason for slow convergence is the
\engterm{condition-number pathology}. The function has a long
curved valley along the parabola $y = x^{2}$, with steep walls
perpendicular to the valley axis and gentle slope along it. Each
gradient-descent step is dominated by the steep cross-valley
component, taking the iterate down to the valley floor and then
oscillating across the floor at every step instead of moving along
the floor toward the minimum. The step size $\eta = 0.001$ is
tuned to the steep cross-valley curvature ($\lambda \approx 100$);
the along-valley curvature ($\lambda \approx 0.6$) is two-and-a-
half orders of magnitude smaller, so the effective per-step
progress along the valley is two-and-a-half orders smaller than the
cross-valley component. A larger $\eta$ would diverge cross-valley;
a smaller would crawl. The fix is either to scale the step by the
Hessian (Newton's method, $\Delta\vect{x} = -H^{-1}\grad f$, which
gives one-step convergence on a pure quadratic), to add momentum
(Polyak heavy-ball; gradient descent with momentum effectively
reduces $\kappa$ to $\sqrt{\kappa}$), or to use an adaptive method
(Adam, RMSprop) that rescales each coordinate by its running
gradient magnitude. Volume II Chapter 15 develops the algorithmic
remedies. The Rosenbrock function is the canonical test problem
because its ill-conditioned valley is exactly the failure mode that
every general-purpose optimiser must defeat to be taken seriously
\cite{paper:v2c11-rosenbrock-1960}.
\end{exercise}

\begin{exercise}
Numerically locate and classify all critical points of $f(x, y) =
\sin(x) \cos(y)$ on $[0, 2\pi]^{2}$ by solving $\grad f = \vect{0}$
and computing the Hessian at each candidate. Report the eigenvalues
at each critical point.

\paragraph{Solution sketch.} The gradient is $\grad f = (\cos x
\cos y, -\sin x \sin y)$. Setting $\cos x \cos y = 0$ and $\sin x
\sin y = 0$ gives two product factorisations. Either $\cos x = 0$
and $\sin x = 0$ (impossible), or $\cos y = 0$ and $\sin y = 0$
(impossible), or $\cos x = 0$ and $\sin y = 0$, or $\cos y = 0$ and
$\sin x = 0$. The two surviving cases enumerate:
\begin{itemize}[leftmargin=*]
  \item $\cos x = 0$ and $\sin y = 0$: $x \in \set{\pi/2, 3\pi/2}$,
    $y \in \set{0, \pi, 2\pi}$. Six candidates (counting boundary).
  \item $\cos y = 0$ and $\sin x = 0$: $x \in \set{0, \pi, 2\pi}$,
    $y \in \set{\pi/2, 3\pi/2}$. Six candidates.
\end{itemize}
The Hessian is
\[
H = \begin{pmatrix} -\sin x \cos y & -\cos x \sin y \\
-\cos x \sin y & -\sin x \cos y \end{pmatrix}.
\]
At an interior critical point of the first kind ($\cos x = 0$,
$\sin y = 0$), the off-diagonal vanishes ($\cos x = 0$), and the
diagonal is $-\sin x \cos y$. At $(x, y) = (\pi/2, \pi)$, $\sin x =
1$, $\cos y = -1$, so $H = \diag(1, 1)$, eigenvalues $1, 1$, local
minimum. At $(\pi/2, 0)$: $H = \diag(-1, -1)$, local maximum.
At $(3\pi/2, 0)$: $\sin x = -1$, $\cos y = 1$, $H = \diag(1, 1)$,
local minimum. At $(3\pi/2, \pi)$: $H = \diag(-1, -1)$, local
maximum. At a critical point of the second kind ($\sin x = 0$,
$\cos y = 0$), the diagonal vanishes and the off-diagonal is
$-\cos x \sin y$. At $(0, \pi/2)$: $\cos x = 1$, $\sin y = 1$,
$H = \begin{pmatrix} 0 & -1 \\ -1 & 0 \end{pmatrix}$, eigenvalues
$\pm 1$, saddle. The same off-diagonal structure makes every
critical point of the second kind a saddle, with eigenvalues $\pm
1$. The two kinds alternate on a $\pi$-by-$\pi$ grid: minima and
maxima at $(\pi/2, 0), (\pi/2, \pi), (3\pi/2, 0), (3\pi/2, \pi)$,
saddles at $(0, \pi/2), (\pi, \pi/2), (0, 3\pi/2), (\pi, 3\pi/2)$.
The level-set picture is the standard "egg-carton" surface that is
the canonical visualisation of a doubly-periodic critical-point
pattern \cite{text:v2ch11-marsden-tromba-2011}.
\end{exercise}

\begin{exercise}
Numerically integrate $\iiint_{V} (x^{2} + y^{2} + z^{2})\,\dd V$
over the unit ball $V = \set{\norm{\vect{x}} \leq 1}$ by Monte
Carlo, using $N = 10^{4}, 10^{5}, 10^{6}$ samples. Report the
estimate and a $95\%$ confidence interval at each sample size, and
compare to the analytical answer $4\pi/5$ obtained by spherical
coordinates.

\paragraph{Solution sketch.} The analytical answer in spherical
coordinates: $\iiint_{V} r^{2} \cdot r^{2}\sin\theta\,\dd r\,\dd
\theta\,\dd\phi = 2\pi \cdot 2 \cdot R^{5}/5 = 4\pi/5 \approx 2.513$
for $R = 1$. The Monte Carlo estimator uses two pieces. First,
generate uniform samples in the cube $[-1, 1]^{3}$ and reject those
outside the ball; the acceptance rate is $V_{3}(1)/8 = (4\pi/3)/8
= \pi/6 \approx 0.524$. Second, for accepted samples, average
$(x^{2} + y^{2} + z^{2})$ and multiply by the ball volume:
\[
\hat{I} = V_{3}(1) \cdot \frac{1}{N_{\text{accepted}}}
\sum_{\text{accepted}} (x_{i}^{2} + y_{i}^{2} + z_{i}^{2}).
\]
The sample mean has standard error $\sigma/\sqrt{N_{\text{accepted}}
}$ where $\sigma^{2}$ is the variance of $r^{2}$ over the uniform
ball. Numerically, $\Var[r^{2}] = \E[r^{4}] - (\E[r^{2}])^{2} =
3/7 - (3/5)^{2} = 3/7 - 9/25 = 12/175 \approx 0.0686$ (taking
$R = 1$). For accepted-sample counts $N_{\text{accepted}} \approx
0.524 N$:
\begin{itemize}[leftmargin=*]
  \item $N = 10^{4}$: $N_{\text{accepted}} \approx 5240$, standard
    error of $\hat{I}/V_{3}(1) \approx \sqrt{0.0686/5240} \approx
    0.00362$, so $95\%$ CI half-width $\approx V_{3}(1) \cdot
    1.96 \cdot 0.00362 \approx 0.030$. Expect $\hat{I} = 2.513
    \pm 0.030$.
  \item $N = 10^{5}$: CI half-width $\approx 0.0095$. Expect
    $\hat{I} = 2.513 \pm 0.010$.
  \item $N = 10^{6}$: CI half-width $\approx 0.0030$. Expect
    $\hat{I} = 2.513 \pm 0.003$.
\end{itemize}
The Monte Carlo standard error falls as $1/\sqrt{N}$, the canonical
convergence rate independent of dimension; this is the saving
grace of Monte Carlo in moderate-dimensional integration. The
$1/\sqrt{N}$ scaling means a four-orders-of-magnitude reduction in
error costs eight orders of magnitude in samples; for high-
precision integration of a fixed problem, deterministic methods
(adaptive cubature, sparse grids) are usually faster. Volume II
Chapter 16 develops the deterministic alternatives.
\end{exercise}

\subsection*{Design}\pathtag{standard}

\begin{exercise}
Design the optimal proportions of an open-topped rectangular box
of fixed volume $V_{0}$, minimising the total material area
(bottom plus four sides). Set up as a multivariable Lagrange
problem, solve, and report the ratios of length, width, and height.

\paragraph{Discussion.} Let $\ell, w, h$ denote length, width, and
height. The material area is $A = \ell w + 2 \ell h + 2 w h$ (one
bottom, two front-back sides, two left-right sides). The volume
constraint is $V = \ell w h = V_{0}$. The Lagrangian:
\[
\mathcal{L} = \ell w + 2\ell h + 2 w h - \lambda (\ell w h - V_{0}).
\]
The first-order conditions are
\begin{align*}
\pdv{\mathcal{L}}{\ell} &= w + 2 h - \lambda w h = 0, \\
\pdv{\mathcal{L}}{w} &= \ell + 2 h - \lambda \ell h = 0, \\
\pdv{\mathcal{L}}{h} &= 2\ell + 2 w - \lambda \ell w = 0, \\
\pdv{\mathcal{L}}{\lambda} &= -(\ell w h - V_{0}) = 0.
\end{align*}
Subtract the first two: $w - \ell + 2 h - 2 h - \lambda h(w - \ell)
= (w - \ell)(1 - \lambda h) = 0$. Either $w = \ell$, or $\lambda h
= 1$. Suppose $w = \ell$. Then the first equation becomes $\ell + 2
h = \lambda \ell h$, and the third becomes $4\ell = \lambda \ell^{2}$,
giving $\lambda = 4/\ell$. Substituting: $\ell + 2 h = (4/\ell)
\cdot \ell h = 4 h$, so $\ell = 2 h$, that is, $h = \ell/2$. The
constraint $\ell^{2} h = V_{0}$ gives $\ell^{3}/2 = V_{0}$, $\ell =
(2 V_{0})^{1/3}$, $w = \ell$, $h = \ell/2$. The optimal proportions
are $\ell : w : h = 2 : 2 : 1$, a square-base box twice as long as
it is tall. The bordered-Hessian (or direct second-derivative)
check confirms this is a constrained minimum, not a saddle.

The alternative branch $\lambda h = 1$ together with the first
equation $w + 2h = \lambda w h = w$ forces $h = 0$, infeasible
under $V = V_{0} > 0$. So the constrained minimum is unique up to
permutation: $\ell = w = (2 V_{0})^{1/3}$, $h = (V_{0}/4)^{1/3}$.

A reader who solves this problem unconstrained by substitution
($h = V_{0}/(\ell w)$, area $A(\ell, w) = \ell w + 2 V_{0}/w +
2 V_{0}/\ell$) and sets $\grad A = 0$ recovers the same answer.
The Lagrange method generalises better when the constraint is not
solvable for one variable, but the problems are equivalent here.
The optimal aspect ratio $\ell : w : h = 2 : 2 : 1$ is the
prescription that the perimeter material (sides) and floor material
(bottom) trade against each other at equal marginal cost.
\end{exercise}

\begin{exercise}
Design a thermal-insulation layer thickness profile for a cylindrical
pipe carrying fluid at $80\,\si{\degreeCelsius}$ in an environment
at $10\,\si{\degreeCelsius}$. The pipe length is fixed; the
insulation cost per unit volume is $c_{1}$; the heat-loss penalty
per watt-year is $c_{2}$. Set up the optimisation, identify the
constraint, and solve for the optimal thickness as a function of
the cost ratio $c_{2}/c_{1}$.

\paragraph{Discussion.} Let $r_{i}$ be the pipe outer radius (fixed),
$r_{o}$ the outer radius of the insulation, and $t = r_{o} - r_{i}$
the insulation thickness (the design variable). Per unit pipe
length, the insulation cross-section volume is $\pi(r_{o}^{2} -
r_{i}^{2})$, and the steady-state radial heat-loss rate through
cylindrical conduction (Fourier's law in cylindrical geometry) is
\[
\dot{Q} = \frac{2\pi k\,\Delta T}{\ln(r_{o}/r_{i})},
\]
where $k$ is the insulation conductivity and $\Delta T = T_{\text{
fluid}} - T_{\text{env}} = 70\,\si{\kelvin}$. The total annualised
cost per unit pipe length is
\[
C(r_{o}) = c_{1}\,\pi(r_{o}^{2} - r_{i}^{2}) + c_{2}\,
\frac{2\pi k\,\Delta T}{\ln(r_{o}/r_{i})}.
\]
There is no equality constraint; the problem is an unconstrained
scalar optimisation in $r_{o}$. Set $\dv{C}{r_{o}} = 0$:
\[
2 \pi c_{1} r_{o} - \frac{2\pi c_{2} k\,\Delta T}{r_{o}\,\ln^{2}
(r_{o}/r_{i})} = 0,
\]
which rearranges to the implicit equation
\[
r_{o}^{2}\,\ln^{2}(r_{o}/r_{i}) = \frac{c_{2} k\,\Delta T}{c_{1}}.
\]
The optimal outer radius scales as the square root of the cost
ratio $c_{2}/c_{1}$, modulated by a logarithmic factor. The
qualitative reading: when heat is expensive ($c_{2}/c_{1}$ large)
the optimal insulation is thick; when insulation is expensive
relative to heat ($c_{2}/c_{1}$ small), thin insulation suffices.
A worked numerical example: pipe radius $r_{i} =
2.5\,\si{\centi\meter}$, $k = 0.04\,\si{\watt\per\meter\per\kelvin}$
(glass wool), $c_{2}/c_{1} = 100\,\si{\watt\squared\meter\squared
\per\watt\cubed} \cdot 70$ (a unit-consistent value for a $20$-year
amortisation at $\$0.10/\si{\kilo\watt\hour}$), giving $r_{o}^{2}
\ln^{2}(r_{o}/0.025) \approx 0.28$. By iteration, $r_{o} \approx
0.18\,\si{\meter}$, $t \approx 15\,\si{\centi\meter}$ of glass
wool. The reader sanity-checks against typical industrial insulation
thicknesses of $5$ to $20\,\si{\centi\meter}$.

A second-order check confirms $\ddot{C}(r_{o}^{*}) > 0$: the cost
function is convex on the feasible domain $r_{o} > r_{i}$, so the
critical point is the global minimum. One engineering warning: in
geometries where the insulation also reduces the convective heat-
transfer area, there is a \engterm{critical radius of insulation}
below which adding insulation increases heat loss. For a pipe with
external convection coefficient $h_{\text{conv}}$, the critical
radius is $r_{\text{crit}} = k/h_{\text{conv}}$. For glass wool
($k = 0.04$) in still air ($h_{\text{conv}} \approx 10$), $r_{\text{
crit}} \approx 4\,\si{\milli\meter}$, well below the pipe's own
radius, so the warning is moot here. For a small wire ($r_{i} <
r_{\text{crit}}$) the warning bites.
\end{exercise}

\begin{exercise}
Design a least-squares fit of a quadratic surface $z = a x^{2} +
b x y + c y^{2} + d x + e y + f$ to a dataset of $N$ points
$(x_{i}, y_{i}, z_{i})$. Set up the optimisation, derive the normal
equations from the gradient condition, and identify the rank
condition on the design matrix that guarantees a unique solution.

\paragraph{Discussion.} Collect the unknowns into $\boldsymbol{
\beta} = (a, b, c, d, e, f)^{\mathsf{T}} \in \R^{6}$ and the data
into the $N \times 6$ design matrix
\[
X = \begin{pmatrix}
x_{1}^{2} & x_{1} y_{1} & y_{1}^{2} & x_{1} & y_{1} & 1 \\
x_{2}^{2} & x_{2} y_{2} & y_{2}^{2} & x_{2} & y_{2} & 1 \\
\vdots & \vdots & \vdots & \vdots & \vdots & \vdots \\
x_{N}^{2} & x_{N} y_{N} & y_{N}^{2} & x_{N} & y_{N} & 1
\end{pmatrix},
\qquad
\vect{z} = (z_{1}, \ldots, z_{N})^{\mathsf{T}}.
\]
The fit minimises the squared residual
$L(\boldsymbol{\beta}) = \norm{X \boldsymbol{\beta} - \vect{z}}^{2}$.
The gradient with respect to $\boldsymbol{\beta}$ is
\[
\grad L = 2 X^{\mathsf{T}}(X \boldsymbol{\beta} - \vect{z}),
\]
and setting $\grad L = \vect{0}$ gives the \engterm{normal
equations}
\[
X^{\mathsf{T}} X \boldsymbol{\beta} = X^{\mathsf{T}} \vect{z}.
\]
The Hessian is $H = 2 X^{\mathsf{T}} X$, which is positive
semi-definite; positive-definite if and only if $X$ has full column
rank, that is, $\rank X = 6$. The rank condition requires $N \geq
6$ (necessary, not sufficient) and a sufficiently spread-out data
set so that no six of the columns of $X$ are linearly dependent.
The pathological cases:
\begin{itemize}[leftmargin=*]
  \item All $N$ points lie on a line $y = m x + b$. Then $y_{i} =
    m x_{i} + b$ and the columns of $X$ acquire a linear relation,
    making $\rank X < 6$. Geometrically: a quadratic surface
    restricted to a line is a quadratic in one variable, with three
    parameters, not six, so the fit is under-determined.
  \item All $N$ points lie on a quadratic curve. Similar collapse.
  \item Fewer than six distinct $(x_{i}, y_{i})$ pairs. Trivially
    under-determined.
\end{itemize}
The rank condition is the same as the condition for the
least-squares minimiser to be unique. When $X$ is rank-deficient,
the normal equations have infinitely many solutions; the standard
remedy is the pseudo-inverse $\boldsymbol{\beta}^{*} =
X^{+}\vect{z}$, which selects the minimum-norm solution among the
infinitely many. Volume II Chapter 10 develops the SVD pseudo-
inverse machinery. The numerical-conditioning warning of section
11.3 applies: a nearly rank-deficient $X$ (one whose smallest
singular value is small relative to its largest) produces a fit
whose coefficients are sensitive to noise in $\vect{z}$, and the
practical remedy is regularisation (ridge regression, Volume II
Chapter 14), not just solving the normal equations with floating-
point hope.
\end{exercise}

\subsection*{Diagnosis and reverse engineering}\pathtag{standard}

\begin{exercise}
A colleague reports an iterative optimiser that has converged on a
critical point with $\grad f \approx \vect{0}$ but for which the
final loss is much higher than expected. Propose three diagnostic
tests, in order of cost, to determine whether the iterate is at a
local minimum, a saddle, or a local maximum.

\paragraph{Discussion.} Three tests in order of cost.

\textit{Test 1: random perturbation probe} (cheap, $O(n)$ per
sample). Perturb the iterate $\vect{x}^{*}$ by a small random
$\vect{\delta}$ of norm $\varepsilon$ and evaluate $\Delta f =
f(\vect{x}^{*} + \vect{\delta}) - f(\vect{x}^{*})$. Repeat for $k$
random directions (say, $k = 20$ to $50$). At a local minimum,
every $\Delta f > 0$ to second order in $\varepsilon$; at a local
maximum, every $\Delta f < 0$; at a saddle, $\Delta f$ takes both
signs, with the fraction of negative directions reflecting the
number of negative Hessian eigenvalues. A cheap and pragmatic first
look that requires only function evaluations, no derivatives. The
diagnostic confounds: small $\varepsilon$ is noisy at floating
precision; large $\varepsilon$ probes the global shape of $f$, not
the local critical-point structure.

\textit{Test 2: Hessian-vector probe} (moderate, $O(n)$ per probe).
Use a forward-mode automatic-differentiation or finite-difference
implementation of the Hessian-vector product $H \vect{v}$ for a
random unit vector $\vect{v}$. Compute the Rayleigh quotient
$\vect{v}^{\mathsf{T}} H \vect{v}$; its sign reports the sign of
the projection of $H$ onto $\vect{v}$. Sample $\vect{v}$ from many
directions to estimate the eigenvalue spectrum. The Lanczos
iteration on $H$ recovers the few extreme eigenvalues at cost
$O(n \cdot k)$ for $k$ iterations, without ever forming $H$
explicitly. The diagnostic resolves the test-1 confounds: it
reports the exact second-derivative structure.

\textit{Test 3: full Hessian eigendecomposition} (expensive, $O(n^{3})$).
Compute $H$ explicitly and find all eigenvalues by standard linear
algebra. Required when the spectrum's full structure matters (a
saddle with one negative direction is different operationally from
a saddle with half-negative directions) or when the iterate is
suspected to sit in a near-degenerate basin where many eigenvalues
are near zero. For $n > 10^{3}$ to $10^{4}$ this is impractical;
for smaller problems it is definitive.

The ordering is by cost and by diagnostic resolution. In modern
machine-learning practice, test 1 is the default cheap check; test
2 (especially via implicit Hessian-vector products) is the working
standard for non-trivial models \cite{paper:v2c11-yao-hessian-2020};
test 3 is reserved for problems where the full spectrum carries the
diagnosis. A saddle diagnosed at test 2 admits the remedy of adding
a negative-curvature step (Newton step along the most-negative
eigenvector) to escape, which gradient descent cannot do.
\end{exercise}

\begin{exercise}
A double integral $\iint_{D} f \,\dd A$ is computed by an iterated
integral with the orders swapped, and the answer differs from a
direct Cartesian computation by a sign. Identify the most likely
cause of the discrepancy and state the rule the engineer applies
to recover the correct sign.

\paragraph{Discussion.} The Fubini statement of section 11.6.1
swaps the order of iterated integration on a rectangular domain
without sign change. A sign discrepancy almost always traces to one
of two errors. First, mis-stated bounds: when swapping the order on
a non-rectangular region, the inner-integral bounds (which depended
on the outer variable) must be re-derived for the new outer
variable. A region described by $0 \leq x \leq 1$, $x \leq y \leq
1$ (above the line $y = x$ in the unit square) re-describes as
$0 \leq y \leq 1$, $0 \leq x \leq y$. Reversing the bounds without
the re-description produces a smaller region (or a different
region), giving a different and often wrong answer. The sign flip
is the symptom; the cause is a region described from the wrong
direction.

Second, the bounds may be written in the wrong order: $\int_{a}^{b}
f\,\dd x$ versus $\int_{b}^{a} f\,\dd x = -\int_{a}^{b} f\,\dd x$.
A symbolic-computation system that does not enforce $a < b$ in its
bound convention can silently produce the negated answer. The
engineering rule is: when swapping the order of iterated
integration on a non-rectangular region, re-draw the region of
integration, re-determine which variable bounds which, and write
the new outer integral first with constant bounds and the new inner
integral with bounds that depend on the new outer variable. If both
orderings produce the same answer, the swap is correct. If they
differ in sign, the region was re-described incorrectly. The
worked example of section 11.6.1 (the triangle $0 \leq y \leq x$
in the unit square) shows the procedure in operation. A purely
algebraic check, useful as a sanity-test on automated systems, is
to evaluate the integral on a small rectangular subregion where
Fubini's theorem applies trivially; if the orientations agree there,
the systematic sign error sits in the re-description of the
non-rectangular region.
\end{exercise}

\begin{exercise}
A change-of-variables routine in a numerical-integration package
returns the wrong answer when the user supplies a transformation
whose Jacobian determinant changes sign across the integration
region. Identify the failure mode and propose the correct treatment.

\paragraph{Discussion.} The change-of-variables theorem of section
11.7.1 carries an absolute value: $\dd V_{\vect{x}} = \abs{\det
J(\vect{u})}\,\dd V_{\vect{u}}$. A library implementation that
omits the absolute value (or computes a signed determinant and
forgets to take its modulus) returns negative-area contributions in
regions where $\det J < 0$, partially cancelling the positive-area
contributions where $\det J > 0$. The reported integral is a signed
sum rather than an unsigned area integral; for an integrand
$f \equiv 1$, the routine returns the algebraic area, not the
geometric area.

A second confound: a one-to-one hypothesis is required for the
change-of-variables theorem to apply. A transformation whose
Jacobian determinant changes sign across the region typically fails
the one-to-one hypothesis: the map folds the region $D'$ over
itself, so a point in $D$ is the image of two or more points in
$D'$. The folded region's contribution to the integral is
double-counted by a naive routine that does not respect the folding.

The correct treatment is twofold. First, partition the $\vect{u}$
domain into subregions on which $\det J$ has constant sign. On each
subregion, the map is locally one-to-one and the change-of-
variables theorem applies with $\abs{\det J}$. Sum the subregion
integrals to recover the total. Second, when the map is not
one-to-one across the partition, account for the multiplicity:
$\int_{D} f\,\dd V = \int_{D'} f(\vect{T}(\vect{u}))\,\abs{\det J}\,
\dd V_{\vect{u}}$ counts each point of $D$ once per preimage. If
two subregions of $D'$ map to the same subregion of $D$, the
integral is over-counted by the multiplicity. The correct working
rule is to use $D'$ subregions that are individually one-to-one
images of disjoint subregions of $D$, then sum.

In practice the safest test is to compute the same integral by two
different change-of-variable routes and require agreement. If they
disagree, the offending routine has lost the absolute value, the
one-to-one hypothesis, or both. A second sanity-check: evaluate the
integral with $f \equiv 1$; the result must be the geometric area
of $D$, computable by an independent measure.
\end{exercise}

\subsection*{Failure analysis}\pathtag{core}

\begin{exercise}
A team trains a high-dimensional model with $p = 10^{6}$ parameters
and reports a "local minimum" diagnosed by the small magnitude of
the loss gradient. Argue, in 200 to 300 words, why this diagnosis
is not by itself enough to conclude that the iterate is at a local
minimum, and propose an additional check using the Hessian or its
estimates.

\paragraph{Discussion.} A small gradient identifies a critical
point, not a local minimum. The critical-point condition $\grad L =
\vect{0}$ admits local minima, local maxima, and saddles; in high
dimensions, the relative frequency of these three populations is
heavily skewed toward saddles. The result of
\textcite{paper:v2ch11-dauphin-saddle-2014} establishes that for
random Gaussian loss surfaces in $n$ dimensions, the expected
fraction of critical points at error level $\varepsilon$ above the
global minimum that are local minima decays exponentially in $n$;
the rest are saddles, with the index (number of negative Hessian
eigenvalues) growing with the error. For $n = 10^{6}$, this skew
is overwhelming: a critical point at non-trivial loss is almost
certainly a saddle, not a local minimum.

The additional check uses the Hessian. The exact full Hessian is
intractable at $n = 10^{6}$ (the matrix has $10^{12}$ entries), but
two practical estimates are available. First, the
Lanczos-iteration estimate of the few extreme eigenvalues, using
only Hessian-vector products implementable via reverse-mode
autodiff at twice the cost of one gradient. A negative
eigenvalue at any direction proves the iterate is a saddle. Second,
the stochastic Hutchinson estimator of the Hessian trace and
spectrum-density, using random probe vectors. A trace much smaller
than $n \cdot \norm{H}_{\max}$ signals near-flat directions; a
spectrum with non-trivial negative mass signals a saddle. A team
that reports a "local minimum" from gradient norm alone has
produced a critical-point diagnosis, not a minimum diagnosis. The
Hessian probe is the working check for high-dimensional minima, and
in practice the answer is usually "saddle, not minimum"
\cite{paper:v2c11-yao-hessian-2020}.
\end{exercise}

\begin{exercise}
A nearest-neighbour classifier achieves $99\%$ accuracy on a
$5$-feature problem and $52\%$ accuracy on a $500$-feature
problem with the same number of training points, despite the
$500$-feature problem nominally containing more information.
Argue, in 200 to 300 words, why the curse of dimensionality
predicts this outcome and propose a working remedy.

\paragraph{Discussion.} The nearest-neighbour classifier relies on
the assumption that nearby points in feature space carry the same
class label. Section 11.8.4 collapsed this assumption in high
dimensions: the ratio of the farthest-pair distance to the closest-
pair distance among $N$ uniform random points in $[0, 1]^{n}$ tends
to $1$ as $n$ grows. In $5$ dimensions with a few hundred training
points, "nearby" is still a meaningful notion and a query point
typically sits closer to a same-class training point than to any
other; the classifier achieves near-perfect accuracy. In $500$
dimensions, every training point is approximately equidistant from
every other, the query point sits approximately equidistant from
all of them, and the "nearest" neighbour is barely closer than the
average and carries no special predictive power. The accuracy
collapses to chance ($\approx 50\%$ for a balanced binary problem),
which is consistent with the reported $52\%$. The $500$-feature
problem nominally contains more information; in practice the
extra features dilute the signal across an exponentially-growing
volume in which the classifier cannot tell signal from noise. This
is the working content of the curse of dimensionality
\cite{paper:v2ch11-blum-curse-1997}.

The remedy uses the principle of section 11.8.5: engineering data
sits on a much smaller intrinsic manifold than the ambient feature
space. The classifier should operate on the intrinsic-dimensional
representation, not the ambient one. Three working tools: principal
component analysis (Volume II Chapter 10) projects onto the top-$d$
SVD components, recovering an intrinsic-dimension $d \ll 500$
representation; a learned embedding (a neural network bottleneck or
an autoencoder, Volume VII) does the same with non-linear
projections; feature selection by lasso (Volume II Chapter 14)
selects the few features that carry signal and discards the rest.
On the intrinsic representation, nearest-neighbour recovers its
meaning and its accuracy.
\end{exercise}

\begin{exercise}
A Monte Carlo simulation of a six-dimensional integral over the
unit cube uses uniform sampling of $N = 10^{4}$ points and reports
an answer that systematically underestimates the true value
(known from an analytical reference). Argue, in 200 to 300 words,
why uniform sampling in six dimensions is likely to undersample the
boundary region and propose a sampling correction.

\paragraph{Discussion.} Uniform sampling in a six-dimensional unit
cube places equal expected density everywhere; on the question of
"where the volume sits," it gives the right answer (most volume is
in the corners). The systematic underestimation arises when the
integrand is sharply peaked or rapidly varying somewhere other than
the deep interior, and the small sample size cannot resolve the
local peak. With $N = 10^{4}$ in $n = 6$, the average point
spacing in any coordinate direction is $N^{-1/n} = 10^{-4/6}
\approx 0.046$, that is, about $20$ points per unit edge in any
single direction. Resolving a boundary layer of width $0.05$ around
each face of the cube needs samples near that face; uniform
sampling places samples there in proportion to the boundary-layer
volume, but the contribution to the integral from that layer can
be disproportionately large (the integrand is peaked there). The
sample variance is therefore dominated by the rare-event
contribution from samples that land in the boundary layer, and
the empirical mean systematically undershoots when those rare
events are missed.

The sampling correction is \engterm{importance sampling}. Draw
samples from a proposal distribution $q(\vect{x})$ that places
higher density where $f$ is large, and reweight: $\hat{I} =
(1/N)\sum f(\vect{x}_{i})/q(\vect{x}_{i})$. For a boundary-peaked
integrand, choose $q$ to concentrate samples near the boundary
(for instance, a Beta-distribution-weighted product measure that
favours $x_{i}$ near $0$ or $1$). A second corrective: stratified
sampling partitions the cube into smaller cells and draws a fixed
number of samples per cell, guaranteeing coverage of the boundary
regions. Volume II Chapter 16 develops quasi-Monte Carlo and
adaptive sampling as the systematic remedies. The Sobol sequence
of low-discrepancy points is the working default for moderate-
dimensional integration of smooth integrands.
\end{exercise}

\subsection*{Judgment}\pathtag{standard}

\begin{exercise}
A regulatory agency requires that an engineered system optimise an
objective subject to a single safety constraint. The contractor
proposes converting the constraint into a penalty term added to
the objective, then optimising unconstrained. Argue, in 200 to 300
words, when this conversion is defensible (the penalty method) and
when it is not (when the constraint must be exactly satisfied or
when the penalty is too small to bind).

\paragraph{Discussion.} The penalty method replaces $\text{
minimise } f \text{ subject to } g(\vect{x}) \leq 0$ with
$\text{minimise } f(\vect{x}) + \rho\,\max(0, g(\vect{x}))^{2}$ for
a penalty weight $\rho > 0$. The method is defensible in three
settings. First, when the constraint represents a soft preference
that the engineer can trade off against the objective; the penalty
weight then encodes the trade-off rate, and a finite $\rho$ is the
honest representation of the engineering judgement. Second, when
the constraint is asymptotic: as $\rho \to \infty$, the penalty
solution converges to the constrained solution, so the penalty
method is a computationally convenient way to approach the
constrained answer when the algorithm cannot handle constraints
directly. Third, when the constraint is statistical (a probability
of failure $\leq \varepsilon$ rather than a hard event), the
penalty method naturally accommodates the soft nature of the
constraint.

The conversion is not defensible when the constraint is hard. A
nuclear reactor must not exceed its safety temperature; a structural
member must not exceed its allowable stress; a medical device must
not exceed its maximum delivered dose. A penalty method with finite
$\rho$ produces a solution that may violate the constraint by an
amount that the penalty alone permits, which is unsafe by
construction. The Lagrange/KKT framework of section 11.5 enforces
the constraint exactly at the optimum (subject to numerical
tolerance), which is the honest formulation. A second failure mode:
even when the constraint is soft, a $\rho$ too small for the
problem's scale produces a solution at which the constraint is
massively violated; a $\rho$ too large makes the problem ill-
conditioned. The penalty method is honest only when the engineer
has reported and justified the choice of $\rho$, including a
sensitivity check on the answer's dependence on it.
\end{exercise}

\begin{exercise}
A textbook claims that "gradient descent always converges to a
local minimum." Identify the assumption under which the claim is
true (it is false in general). Cite a counter-example from the
chapter and propose a corrected statement.

\paragraph{Discussion.} The claim is false in general for at least
three reasons that the chapter has named.

First, gradient descent with a fixed positive step size on a saddle
point is stationary in expectation; an iterate that lands at or
near a saddle is not pushed away by the gradient (the gradient is
zero or near-zero). Section 11.3.5 made the point: in high
dimensions, saddles outnumber minima exponentially, and a
gradient-descent run that does not handle saddles gracefully will
spend most of its time near them
\cite{paper:v2ch11-dauphin-saddle-2014}. The counter-example: the
saddle of section 11.3.3, $f(x, y) = x^{2} - y^{2}$. Gradient
descent from the $y$-axis is exactly stationary at the origin and
never moves.

Second, gradient descent with a step size too large diverges; with
a step size too small, it converges arbitrarily slowly. The
Rosenbrock counter-example of Exercise 11.18 above shows the slow
convergence; an even more egregious failure happens when the step
size exceeds $2/\lambda_{\max}(H)$, where the iterate oscillates
and diverges.

Third, gradient descent on a non-smooth or unbounded function may
fail to converge at all; it may diverge to $-\infty$ on an
unbounded function (e.g.\ $f(x) = -x^{2}$, which is "minimised" at
$-\infty$).

A corrected statement: "Gradient descent with sufficiently small
step size on a smooth, bounded-below function with non-degenerate
critical points and a starting point not at a saddle's stable
manifold converges to a critical point. Almost-everywhere with
respect to random initialisation, the limit is a local minimum
rather than a saddle." The almost-everywhere clause is a theorem of
\textcite{paper:v2c11-lee-saddle-2016}: with random initialisation,
the stable manifold of any saddle has measure zero. Even so,
gradient descent's effective speed near saddles is so slow that
practical optimisers add explicit saddle-escape mechanisms (noise,
negative-curvature steps, momentum).
\end{exercise}

\begin{exercise}
\textcite{text:strang2016} treats the Lagrangian as a "saddle-point
problem" in $(\vect{x}, \boldsymbol{\lambda})$. Argue, in one
paragraph, why the saddle-point interpretation is geometrically
correct and why a naive simultaneous-minimisation of $\mathcal{L}$
in both $\vect{x}$ and $\boldsymbol{\lambda}$ does not give the
constrained optimum.

\paragraph{Discussion.} The Lagrangian $\mathcal{L}(\vect{x},
\lambda) = f(\vect{x}) - \lambda\,g(\vect{x})$ has, at the
constrained optimum $(\vect{x}^{*}, \lambda^{*})$, a vanishing
gradient in both $\vect{x}$ and $\lambda$, so $(\vect{x}^{*},
\lambda^{*})$ is a critical point of $\mathcal{L}$ on the joint
space. Geometrically, $\mathcal{L}$ is minimised in $\vect{x}$
(for fixed $\lambda^{*}$, the inner problem is the unconstrained
minimisation of $f - \lambda^{*} g$) and maximised in $\lambda$
(for fixed $\vect{x}^{*}$, the Lagrangian is linear in $\lambda$
with coefficient $-g(\vect{x}^{*})$, and the constraint
$g(\vect{x}^{*}) = 0$ makes the linear coefficient vanish so that
the supremum over $\lambda$ is finite). The constrained optimum is
therefore a saddle of $\mathcal{L}$ in the joint variable: minimum
in $\vect{x}$, maximum in $\lambda$. The reader of this exercise
will recognise that simultaneously minimising in both variables
sends $\lambda \to \infty$ (or $-\infty$) whenever $g(\vect{x})
\neq 0$, because the linear term $-\lambda g$ has no lower bound;
the iterate runs away rather than approaching the constrained
optimum. The correct algorithm is the saddle-point algorithm
(primal-dual ascent-descent, or the method of alternating
projections onto the constraint surface), which minimises in
$\vect{x}$ and maximises in $\lambda$ simultaneously, recovering
the constrained optimum as the joint critical point. Volume II
Chapter 15 develops the algorithmic version (the augmented
Lagrangian method and the primal-dual interior-point method) that
respect the saddle-point structure.
\end{exercise}

\begin{exercise}
A practitioner argues that high-dimensional engineering problems
should be re-cast in terms of the "intrinsic dimension" of the
data manifold rather than the ambient parameter dimension. Argue,
in 200 to 300 words, when this re-casting is justified and what
diagnostic the engineer should apply to estimate the intrinsic
dimension. (Hint: the singular-value decomposition of Volume II
Chapter 10 is the working tool.)

\paragraph{Discussion.} The re-casting is justified whenever the
data sit on (or near) a low-dimensional manifold inside a
high-dimensional ambient space. Section 11.8.5 named the case: the
geometry of the parameter space scales with $n$ in ways that defeat
low-dimensional intuition, but engineering data is almost never
isotropic in the ambient space. A typical example: a $1000$-pixel
greyscale image of a human face appears in $\R^{1000}$, but the
manifold of all human faces is a much smaller subspace, perhaps
$10$ to $50$-dimensional, parameterised by pose, lighting,
expression, and identity. Working on the ambient $\R^{1000}$ is
expensive and loses statistical power to the curse of
dimensionality; working on the intrinsic manifold is cheap and
recovers the geometric structure.

The diagnostic is the singular-value spectrum of the data matrix.
Take $N$ data points, organise them into an $N \times n$ matrix
$X$, centre it, and compute the SVD $X = U \Sigma V^{\mathsf{T}}$
(Volume II Chapter 10). The singular values $\sigma_{1} \geq
\sigma_{2} \geq \cdots \geq \sigma_{\min(N, n)}$ report the
variance captured by each principal direction. The intrinsic
dimension $d$ is estimated as the smallest $k$ such that
$\sum_{i = 1}^{k} \sigma_{i}^{2} / \sum_{i} \sigma_{i}^{2} \geq
1 - \varepsilon$ for a tolerance $\varepsilon$ (commonly $0.01$ or
$0.05$). A spectrum with a sharp elbow at index $k$ signals
intrinsic dimension $k$; a flat spectrum signals genuinely
high-dimensional data with no manifold structure to exploit. The
re-casting is justified when the elbow is sharp; the re-casting is
not justified (or is at best a heuristic) when the spectrum decays
gradually, because the choice of $d$ is then arbitrary and the
"intrinsic manifold" is an artefact of the truncation rather than
a property of the data. Two engineering guards: report the
truncation error, and validate the re-casting on a held-out test
set rather than on the data used to fit the SVD.
\end{exercise}

\begin{exercise}
A student writes the KKT conditions for a problem with one equality
and one inequality, but treats the inequality multiplier $\mu$ as
unconstrained in sign. Argue, in one paragraph, where the analysis
will go wrong and which of the four KKT conditions the student has
violated.

\paragraph{Discussion.} The student has violated dual feasibility
($\mu \geq 0$), the third KKT condition of section 11.5.2. The
geometric content of $\mu \geq 0$ is that an inequality $h \leq 0$
is helpful (slack helps) and an active inequality contributes a
gradient in the direction of the inactive half-space, so the
optimality direction sign of $\mu$ is determined. A student who
ignores the sign restriction will, in some cases, "find" a
candidate $\vect{x}^{*}$ with $\mu^{*} < 0$ at the active set.
Three things go wrong. First, the candidate is not actually a
minimum: the negative multiplier reports that the objective can be
further reduced by moving into the strictly feasible interior
(away from the active constraint), which means the assumed active
set is wrong and the candidate is not optimal. Second, the
shadow-price interpretation breaks: a negative $\mu$ would say
that relaxing the constraint hurts the objective, which contradicts
the meaning of an inequality. Third, the convexity argument that
makes KKT sufficient (for convex programs) fails without the sign
restriction, so a candidate found in this way carries no
sufficiency guarantee. The correct procedure when $\mu^{*} < 0$
shows up is to declare the assumed active set wrong, drop that
inequality from the active set (set $\mu = 0$ by complementary
slackness), and re-solve the equality-only Lagrange system on the
new active set. The vertex-enumeration logic of the worked LP
example in section 11.4's estimation block illustrates the
procedure in operation. The student's error is the most common
one in first encounters with KKT, and the diagnostic is always
to inspect the multiplier signs.
\end{exercise}

\begin{exercise}
A vendor of an optimisation library reports that "the package solves
arbitrary constrained-optimisation problems." Argue, in 200 to
300 words, what classes of problem the vendor's claim is honest
for, what classes it is misleading for, and what the engineer
should ask before selecting the package. \textcite{text:boyd-vandenberghe2004}
gives the working taxonomy of convex programs against which the
claim should be tested.

\paragraph{Discussion.} The claim is honest for problems in the
\engterm{convex} class: convex objective $f$, affine equality
constraints, and convex inequality constraints. The KKT conditions
are necessary and sufficient for convex programs, and a wide range
of polynomial-time algorithms (interior-point methods,
operator-splitting methods, projected-gradient methods) solve them
to global optimality. The taxonomy \cite{text:boyd-vandenberghe2004}
identifies linear programs, quadratic programs, second-order cone
programs, semi-definite programs, and geometric programs as
convex-program subclasses for which there are mature, reliable
solvers. A vendor whose package implements interior-point methods
for these classes can honestly claim "solves the problem."

The claim is misleading for general non-convex programs. Three
sub-classes name the trouble. Mixed-integer programs (some variables
discrete) are NP-hard in general; a vendor cannot guarantee a
solution in any reasonable time. Non-convex continuous programs
(e.g., neural-network training) have many local minima and saddles;
a solver returns a critical point that may be neither global nor
even a local minimum. Combinatorial programs (e.g., the travelling
salesman) admit only heuristic or branch-and-bound approaches whose
runtime is exponential in the worst case.

The engineer should ask: (i) is the problem convex, and what
solver does the package use? (ii) what guarantees does the package
give on convergence, and to what (global, local, critical point)?
(iii) what is the package's behaviour on rank-deficient, ill-
conditioned, or degenerate problems? (iv) what is the package's
record on the engineer's specific class of problem (e.g., LP with
$10^{6}$ variables, semi-definite programs with sparse structure)?
A vendor unable to answer (i)-(iv) is selling a black box, not a
tool. The discipline is to match the package to the convex-program
class and to inspect the convergence diagnostics on every problem.
\end{exercise}
